\documentclass[11pt]{article}

% ---------- Packages ----------
\usepackage[margin=1in]{geometry}
\usepackage{setspace}
\usepackage{amsmath, amssymb, amsthm, mathtools}
\usepackage{graphicx}
\usepackage{booktabs}
\usepackage{float}
\usepackage{hyperref}
\usepackage{xcolor}
\usepackage{enumitem}
\usepackage{fancyhdr}
\usepackage{svg}

\setstretch{1.15}

% ---------- Hyperref Setup ----------
\hypersetup{
    colorlinks=true,
    linkcolor=blue,
    citecolor=blue,
    urlcolor=blue
}

% ---------- Header ----------
\pagestyle{fancy}
\fancyhf{}
\lhead{Quantitative Data Analysis in Finance}
\rhead{\thepage}
\lfoot{Liam O'Shaughnessy}

% ---------- Theorem Environments ----------
\theoremstyle{definition}
\newtheorem{definition}{Definition}[section]

\theoremstyle{plain}
\newtheorem{theorem}{Theorem}[section]
\newtheorem{proposition}{Proposition}[section]
\newtheorem{corollary}{Corollary}[section]

\theoremstyle{remark}
\newtheorem{remark}{Remark}[section]
\newtheorem{example}{Example}[section]

% ---------- Custom Commands ----------
\newcommand{\E}{\mathbb{E}}
\newcommand{\Var}{\mathrm{Var}}
\newcommand{\Cov}{\mathrm{Cov}}
\newcommand{\R}{\mathbb{R}}
\newcommand{\bs}{\boldsymbol}

% ---------- Title ----------
\title{\textbf{Quantitative Data Analysis in Finance}}
\author{Liam O'Shaughnessy}
\date{\today}

\begin{document}

\maketitle
\tableofcontents
\newpage

\section*{Preface}

These are my notes from when I took ECO 480/FIN 580 in Spring 2026 at Princeton University. I have elaborated considerably on the main content with supplementary material and proofs.

\section{Introduction}

Can machines learn finance? Artificial intelligence (AI) has exploded in success, with areas such as ML and NLP. Successes with DeepBlue, Watson, real-time translation, AlphaGo, AlphaFold, and the ensemble of LLMs have shown great promise. ``Statistical AI'' has relied on giving computers subproblems and rules for solving specific problems, whereas ``Neural AI'' has relied on complex, often unexplainable architectures equipped with big data to solve problems. The main drivers of these have been hardware (CPUs, GPUs, TPUs), software (architectures), and big data (now collecting incredible volumes daily). Three subtypes of machine learning are supervised (labeled data, such as regression and classification), unsupervised (unlabeled data, such as PCA and anomaly detection), and reinforcement (trial and error to maximize rewards).

\begin{definition}[Supervised Learning Setup]
    Let $\{(\vec{x}_i, y_i)\}_{i=1}^n$ be a dataset of $n$ observations, where $\vec{x}_i \in \mathbb{R}^p$ is a vector of $p$ features (covariates) and $y_i \in \mathbb{R}$ is a scalar response. Collect observations into a data matrix $X \in \mathbb{R}^{n \times p}$ and response vector $\vec{y} \in \mathbb{R}^n$. The goal of \textit{supervised learning} is to find a function $\hat{f}: \mathbb{R}^p \to \mathbb{R}$ minimizing the empirical risk $\frac{1}{n}\sum_i L(y_i, f(\vec{x}_i))$ over some function class, such that the \emph{population} risk $\mathbb{E}[L(y,f(\vec{x}))]$ is small on new, unseen inputs. The gap between the two is the subject of the bias--variance decomposition below.
\end{definition}

% [CORRECTED] The original defined only the simple (net) return but then used
% X_t = R_t P_{t-1} in the Euler equation, which requires the GROSS return.
% Both are now defined and the gross convention is fixed for this section.
\begin{definition}[Payoff, Simple Return, Gross Return]
    An asset held from $t-1$ to $t$ delivers the \textit{payoff} $X_t = P_t + D_t$, where $P_t$ is the ex-dividend price and $D_t$ the dividend paid at $t$. Its \textit{simple (net) return} is
    \begin{align}
        R^{\mathrm{net}}_t = \frac{P_t + D_t - P_{t-1}}{P_{t-1}} = \frac{X_t - P_{t-1}}{P_{t-1}},
    \end{align}
    and its \textit{gross return} is
    \begin{align}
        R_t = \frac{X_t}{P_{t-1}} = 1 + R^{\mathrm{net}}_t.
    \end{align}
    Unless stated otherwise, $R$ denotes the \emph{gross} return throughout this section; the Euler equation below is stated in gross returns. The \textit{log return} is $r_t = \log R_t$, which is additive across periods ($r_{t\to t+k} = \sum_{j=1}^k r_{t+j}$) and is the natural object for time-aggregation and for Gaussian modelling.
\end{definition}

\begin{definition}[Euler Equation]
    The canonical asset pricing equation in discrete time is
    \begin{align}
        P_{i,t} = \mathbb{E}\!\left[M_{t+1} X_{i,t+1} \mid \mathcal{F}_t\right],
        \qquad\text{equivalently}\qquad
        1 = \mathbb{E}\!\left[M_{t+1} R_{i,t+1} \mid \mathcal{F}_t\right],
    \end{align}
    the second form following by dividing through by the $\mathcal{F}_t$-measurable price $P_{i,t}$ and using $R_{i,t+1} = X_{i,t+1}/P_{i,t}$.
\end{definition}

% [CORRECTED] An SDF must price ALL assets simultaneously (otherwise the
% definition is vacuous), must be strictly positive and F_{t+1}-measurable,
% and the CCAPM discount factor needs beta^s, not beta, for an s-period SDF.
\begin{definition}[Stochastic Discount Factor]
    A \textit{stochastic discount factor} (SDF), or pricing kernel, is a strictly positive, $\mathcal{F}_{t+1}$-measurable random variable $M_{t+1}$ such that the Euler equation holds \emph{simultaneously for every traded asset} $i$. The requirement that a single $M$ price all assets at once is what gives the definition content: for one asset in isolation, many such random variables exist.

    By the fundamental theorem of asset pricing, a strictly positive SDF exists if and only if the market admits no arbitrage, and it is unique if and only if the market is complete. Under incompleteness the set of admissible SDFs is convex, and the choice of $M$ within it is a modelling decision.

    The CCAPM proposes a \emph{consumption-based} SDF: for a representative agent with period utility $u$ over consumption and subjective discount factor $\beta \in (0,1)$, the intertemporal first-order condition yields
    \begin{align}
        M_{t+s} = \beta^s \frac{u'(c_{t+s})}{u'(c_t)}.
    \end{align}
    The SDF is high in states where marginal utility is high, i.e.\ where consumption is scarce. Under CRRA utility $u'(c) = c^{-\gamma}$, this becomes $M_{t+1} = \beta (c_{t+1}/c_t)^{-\gamma}$.
\end{definition}

\begin{remark}[Relation to Risk-Neutral Pricing]
    The SDF is the discrete-time counterpart of the change of measure used in martingale pricing. Writing $\mathbb{Q}$ for the risk-neutral measure,
    \begin{align}
        M_{t+1} = \frac{1}{R_{f,t}} \left.\frac{d\mathbb{Q}}{d\mathbb{P}}\right|_{\mathcal{F}_{t+1}},
    \end{align}
    so that $P_{i,t} = \mathbb{E}^{\mathbb{P}}_t[M_{t+1}X_{i,t+1}] = R_{f,t}^{-1}\,\mathbb{E}^{\mathbb{Q}}_t[X_{i,t+1}]$: discounting with the SDF under the physical measure and discounting at the risk-free rate under $\mathbb{Q}$ are the same operation. Strict positivity of $M$ is exactly what makes $\mathbb{Q}$ a probability measure equivalent to $\mathbb{P}$.
\end{remark}

\begin{theorem}[Beta Representation of Expected Returns]\label{thm:beta-rep}
    Let $R_{f,t}$ be the gross return on a risk-free asset (known at $t$, hence $\mathcal{F}_t$-measurable). Then for any asset $i$,
    \begin{align}
        \mathbb{E}[R_{i,t+1}\mid\mathcal{F}_t] - R_{f,t}
        &= -R_{f,t}\,\mathrm{Cov}_t(M_{t+1}, R_{i,t+1}) \\
        &= \underbrace{\frac{\mathrm{Cov}_t(M_{t+1},R_{i,t+1})}{\mathrm{Var}_t(M_{t+1})}}_{\textstyle \beta_{i,t}}
           \cdot
           \underbrace{\left(-R_{f,t}\,\mathrm{Var}_t(M_{t+1})\right)}_{\textstyle \lambda_t}
        = \beta_{i,t}\lambda_t.
    \end{align}
    Here $\beta_{i,t}$ is the coefficient from a conditional regression of the asset's return on the SDF (a \emph{quantity} of risk, asset-specific), and $\lambda_t$ is the \emph{price} of risk, common to all assets.
\end{theorem}
\begin{proof}
    Since $P_{i,t}$ is $\mathcal{F}_t$-measurable and $X_{i,t+1} = R_{i,t+1}P_{i,t}$, the Euler equation gives
    \begin{align}
        P_{i,t} = \mathbb{E}[M_{t+1}R_{i,t+1}P_{i,t}\mid\mathcal{F}_t] = P_{i,t}\,\mathbb{E}[M_{t+1}R_{i,t+1}\mid\mathcal{F}_t],
    \end{align}
    so $1 = \mathbb{E}[M_{t+1}R_{i,t+1}\mid\mathcal{F}_t]$. Applying this to a risk-free asset, whose gross return $R_{f,t}$ is $\mathcal{F}_t$-measurable and can be pulled out of the conditional expectation,
    \begin{align}
        1 = R_{f,t}\,\mathbb{E}[M_{t+1}\mid\mathcal{F}_t]
        \quad\Longrightarrow\quad
        R_{f,t} = \frac{1}{\mathbb{E}[M_{t+1}\mid\mathcal{F}_t]}.
    \end{align}
    Decomposing the expectation of the product,
    \begin{align}
        1 = \mathbb{E}[M_{t+1}\mid\mathcal{F}_t]\,\mathbb{E}[R_{i,t+1}\mid\mathcal{F}_t] + \mathrm{Cov}_t(M_{t+1},R_{i,t+1}),
    \end{align}
    and multiplying through by $R_{f,t} = 1/\mathbb{E}_t[M_{t+1}]$,
    \begin{align}
        R_{f,t} = \mathbb{E}[R_{i,t+1}\mid\mathcal{F}_t] + R_{f,t}\,\mathrm{Cov}_t(M_{t+1},R_{i,t+1}).
    \end{align}
    Rearranging gives the first line; multiplying and dividing by $\mathrm{Var}_t(M_{t+1})$ gives the second.
\end{proof}

\begin{remark}[Sign of the Price of Risk]
    Note that $\lambda_t = -R_{f,t}\mathrm{Var}_t(M_{t+1}) < 0$. This is not an error but the entire economic content of the theorem: an asset whose payoff covaries \emph{positively} with the SDF pays off precisely in states where marginal utility is high (bad states). It is insurance, investors bid its price up, and it earns a return \emph{below} the risk-free rate. Assets that covary negatively with the SDF --- procyclical assets, which fail you in a recession --- must compensate with a positive risk premium. Conventional factor models (CAPM, Fama--French) restate this with a factor $f$ that is a decreasing linear function of $M$, which flips the sign and yields the familiar $\lambda > 0$.
\end{remark}

\begin{corollary}[Hansen--Jagannathan Bound]
    For any asset $i$, the conditional Sharpe ratio is bounded by the coefficient of variation of the SDF:
    \begin{align}
        \left|\frac{\mathbb{E}_t[R_{i,t+1}] - R_{f,t}}{\sigma_t(R_{i,t+1})}\right|
        \;\leq\; \frac{\sigma_t(M_{t+1})}{\mathbb{E}_t[M_{t+1}]}
        \;=\; R_{f,t}\,\sigma_t(M_{t+1}).
    \end{align}
\end{corollary}
\begin{proof}
    From Theorem~\ref{thm:beta-rep}, $\mathbb{E}_t[R_{i,t+1}] - R_{f,t} = -R_{f,t}\,\mathrm{Cov}_t(M_{t+1},R_{i,t+1}) = -R_{f,t}\,\rho_t\,\sigma_t(M_{t+1})\sigma_t(R_{i,t+1})$, where $\rho_t$ is the conditional correlation. Since $|\rho_t| \leq 1$, dividing by $\sigma_t(R_{i,t+1})$ and using $R_{f,t} = 1/\mathbb{E}_t[M_{t+1}]$ gives the bound.
\end{proof}

\begin{remark}[Why the HJ Bound Matters in Practice]
    The bound converts any candidate SDF into a testable restriction: a model whose SDF has volatility $\sigma(M)$ cannot rationalize a market Sharpe ratio above $R_f\,\sigma(M)$. This is the sharp form of the equity premium puzzle --- a CRRA consumption SDF calibrated to observed consumption volatility produces a $\sigma(M)$ roughly an order of magnitude too small to support the historical equity Sharpe ratio of $\approx 0.4$, unless $\gamma$ is implausibly large. Read in the other direction, the bound says the maximum attainable Sharpe ratio in a market is a property of the pricing kernel, and a strategy claiming to exceed it is either exploiting a genuinely richer information set or is mismeasured.
\end{remark}

\begin{remark}[CAPM as a Single-Factor Model]
    The Capital Asset Pricing Model (CAPM) is the simplest instance of Theorem~\ref{thm:beta-rep}. Let the SDF be linear in the market return $R_{m,t+1}$ (equivalently, take $M_{t+1} = a - bR_{m,t+1}$ with $b>0$, which flips the sign of $\lambda$ as described above). Then the expected excess return on any asset $i$ is
    \begin{align}
        \mathbb{E}[R_{i,t+1}] - R_{f} = \beta_i (\mathbb{E}[R_{m,t+1}] - R_f), \quad \beta_i = \frac{\text{Cov}(R_i, R_m)}{\text{Var}(R_m)}.
    \end{align}
    Empirically, ``high-beta'' stocks have historically underperformed their CAPM-predicted returns (the ``betting-against-beta'' anomaly).
\end{remark}

\begin{figure}[h]
    \centering
    \includegraphics[width=0.5\linewidth]{capm.png}
    \caption{CAPM}
    \label{fig:placeholder}
\end{figure}

\begin{remark}[Factor Models in Practice]
    The Fama--French three-factor model extends the CAPM by adding size (SMB: Small Minus Big) and value (HML: High Minus Low book-to-market) factors. A time-series regression of excess returns on these factors,
    \begin{align}
        R_{i,t} - R_{f,t} = \alpha_i + \beta_{i,1}(R_{m,t}-R_{f,t}) + \beta_{i,2}\mathrm{SMB}_t + \beta_{i,3}\mathrm{HML}_t + \varepsilon_{i,t},
    \end{align}
    tests whether $\alpha_i = 0$ (the asset is correctly priced by the model). A statistically significant $\alpha_i$ indicates abnormal return unexplained by the factors.
\end{remark}

\begin{theorem}[Bias--Variance]
    For $\hat{f}(x)$ a function predicting $\{y\}$ observed data, generated by an underlying process, the mean squared error decomposes as
    \begin{align}
        \text{MSE} = \text{Bias}[\hat{f}(x)]^2 + \text{Var}[\hat{f}(x)] + \sigma^2,
    \end{align}
    establishing a fundamental trade-off: reducing bias (by fitting a more complex model) tends to increase variance, and vice versa. Regularization methods navigate this trade-off by deliberately introducing bias to achieve a larger reduction in variance.
\end{theorem}
\begin{proof}
    Consider the observed data as generated by some underlying process with Gaussian errors $f(x)+\varepsilon$, then the error is
    \begin{align}
        \text{MSE} &= \mathbb{E}[(y-\hat{f}(x))^2] = \mathbb{E}[(\varepsilon +f(x)-\hat{f}(x))^2] \\
        &= \mathbb{E}[(f(x)-\hat{f}(x))^2] + 2\mathbb{E}[(f(x)-\hat{f}(x))\varepsilon] + \mathbb{E}[\varepsilon^2]\\
        &=\mathbb{E}[(f(x) - \mathbb{E}[\hat{f}(x)]+\mathbb{E}[\hat{f}(x)]-\hat{f}(x))^2] + 0 + \sigma^2.
    \end{align}
    Expanding this and simplifying yields
    \begin{align}
        \text{MSE} &= (f(x)-\mathbb{E}[\hat{f}(x)])^2 + \mathbb{E}[(\mathbb{E}[\hat{f}(x)]-\hat{f}(x))^2]+\sigma^2\\
        &= \text{Bias}[\hat{f}(x)]^2 + \text{Var}[\hat{f}(x)] + \sigma^2.
    \end{align}
\end{proof}

\begin{definition}[Random Vectors]

Let $\vec{X} \in \mathbb{R}^n$ be a \textit{random vector} with components $\vec{X} = (X_1,\dots,X_n)^T$, each random variables. The expectation is defined componentwise:
\begin{align}
\mathbb{E}[\vec{X}]
&=
\begin{pmatrix}
\mathbb{E}[X_1] \\
\vdots \\
\mathbb{E}[X_n]
\end{pmatrix}.
\end{align}

The variance of a random vector is its covariance matrix:
\begin{align}
\mathrm{Var}(\vec{X})
&= \mathbb{E}\!\left[(\vec{X}-\mathbb{E}[\vec{X}])(\vec{X}-\mathbb{E}[\vec{X}])^T\right].
\end{align}

Equivalently,
\begin{align}
\mathrm{Var}(\vec{X})
&= \mathbb{E}[\vec{X}\vec{X}^T] - \mathbb{E}[\vec{X}]\mathbb{E}[\vec{X}]^T.
\end{align}

Considering $\mathrm{Var}(\vec{X})_{ij} = \mathrm{Cov}(X_i,X_j)$, the diagonal entries are variances and the off-diagonal entries are covariances.
\end{definition}

\begin{proposition}[Linear Transformations]
If $\vec{Y} = B\vec{X} + \vec{c}$ for matrix $B$ and vector $\vec{c}$, then
\begin{align}
\mathbb{E}[\vec{Y}] &= B\,\mathbb{E}[\vec{X}] + \vec{c}, \\
\mathrm{Var}(\vec{Y}) &= B\,\mathrm{Var}(\vec{X})\,B^T.
\end{align}
\end{proposition}
\begin{proof}
    Linearity of expectation gives $\mathbb{E}[\vec{Y}] = B\,\mathbb{E}[\vec{X}] + \vec{c}$ directly. For the variance,
    \begin{align}
        \mathrm{Var}(\vec{Y}) &= \mathbb{E}[(\vec{Y} - \mathbb{E}[\vec{Y}])(\vec{Y} - \mathbb{E}[\vec{Y}])^T] \\
        &= \mathbb{E}[B(\vec{X} - \mathbb{E}[\vec{X}])(B(\vec{X} - \mathbb{E}[\vec{X}]))^T] \\
        &= B\,\mathbb{E}[(\vec{X} - \mathbb{E}[\vec{X}])(\vec{X} - \mathbb{E}[\vec{X}])^T]\,B^T = B\,\mathrm{Var}(\vec{X})\,B^T,
    \end{align}
    since the constant $\vec{c}$ cancels in the centering and $B$ can be pulled out of the expectation.
\end{proof}


\section{Estimation}

Before studying any particular estimator, it is worth being precise about what an estimator \emph{is}, and about the two limit theorems that govern essentially all of the asymptotic theory in these notes.

\begin{definition}[Statistic and Estimator]
    Let $\mathcal{D}_n = \{y_1,\dots,y_n\}$ be a sample. A \emph{statistic} is any measurable function $T(\mathcal{D}_n)$ of the sample that does not depend on unknown parameters. An \emph{estimator} of a parameter $\theta$ is a statistic $\hat{\theta}_n = \hat{\theta}(\mathcal{D}_n)$ intended to approximate it.

    An estimator is a \emph{random variable}: it is a rule, a function of a sample that has not yet been drawn, and it has a distribution. An estimate is the single realized number obtained by applying that rule to the data actually observed.
\end{definition}

\begin{definition}[Sampling Distribution and Standard Error]
    The \emph{sampling distribution} of $\hat{\theta}_n$ is its distribution induced by the randomness of the sample. Its standard deviation is the \emph{standard error},
    \begin{align}
        \mathrm{se}(\hat{\theta}_n) = \sqrt{\mathrm{Var}(\hat{\theta}_n)}.
    \end{align}
    A standard error is describing how much the estimate would move across hypothetical resamples. This is why standard errors can be estimated by resampling (the bootstrap) or read off an asymptotic formula.
\end{definition}

\begin{definition}[Bias and Mean Squared Error]
    For a scalar parameter $\theta$ and estimator $\hat{\theta}_n$,
    \begin{align}
        \mathrm{Bias}(\hat{\theta}_n) = \mathbb{E}[\hat{\theta}_n] - \theta, \qquad
        \mathrm{MSE}(\hat{\theta}_n) = \mathbb{E}\!\left[(\hat{\theta}_n - \theta)^2\right].
    \end{align}
    The estimator is \emph{unbiased} if $\mathrm{Bias}(\hat{\theta}_n) = 0$ for every $\theta \in \Theta$. In the vector case, $\mathrm{MSE}$ is replaced by the matrix $\mathbb{E}[(\hat{\theta}_n-\theta)(\hat{\theta}_n-\theta)^T]$ or by its trace.
\end{definition}

\begin{proposition}[Bias--Variance Decomposition for an Estimator]\label{prop:mse-est}
    \begin{align}
        \mathrm{MSE}(\hat{\theta}_n) = \mathrm{Bias}(\hat{\theta}_n)^2 + \mathrm{Var}(\hat{\theta}_n).
    \end{align}
\end{proposition}
\begin{proof}
    Add and subtract $\mathbb{E}[\hat{\theta}_n]$ inside the square:
    \begin{align}
        \mathbb{E}\!\left[(\hat{\theta}_n - \theta)^2\right]
        &= \mathbb{E}\!\left[\left((\hat{\theta}_n - \mathbb{E}[\hat{\theta}_n]) + (\mathbb{E}[\hat{\theta}_n] - \theta)\right)^2\right] \\
        &= \mathbb{E}\!\left[(\hat{\theta}_n - \mathbb{E}[\hat{\theta}_n])^2\right] + 2\underbrace{\mathbb{E}\!\left[\hat{\theta}_n - \mathbb{E}[\hat{\theta}_n]\right]}_{=\,0}\left(\mathbb{E}[\hat{\theta}_n]-\theta\right) + \left(\mathbb{E}[\hat{\theta}_n]-\theta\right)^2,
    \end{align}
    the cross term vanishing because $\mathbb{E}[\hat{\theta}_n] - \theta$ is a constant.
\end{proof}

\begin{remark}[Two Different Bias--Variance Decompositions]
    Proposition~\ref{prop:mse-est} looks like the Bias--Variance theorem of Section~1 but decomposes a different quantity, and the two should not be conflated.
    \begin{itemize}[noitemsep]
        \item \textbf{Estimation.} $\mathbb{E}[(\hat{\theta}-\theta)^2] = \mathrm{Bias}^2 + \mathrm{Var}$. The target $\theta$ is a fixed unknown constant. There is \emph{no} irreducible term: with enough data the error can in principle be driven to zero.
        \item \textbf{Prediction.} $\mathbb{E}[(y - \hat{f}(x))^2] = \mathrm{Bias}^2 + \mathrm{Var} + \sigma^2$. The target $y$ is itself random. The extra $\sigma^2$ is irreducible noise and no estimator, however good, removes it.
    \end{itemize}
    The distinction matters when interpreting an $R^2$: the prediction floor is set by $\sigma^2$, which in return prediction is enormous relative to the signal.
\end{remark}

\begin{definition}[Dominance and Admissibility]
    An estimator $\hat{\theta}$ \emph{dominates} $\tilde{\theta}$ if $\mathrm{MSE}(\hat{\theta}) \leq \mathrm{MSE}(\tilde{\theta})$ for every $\theta \in \Theta$, with strict inequality for at least one $\theta$. An estimator is \emph{admissible} if no estimator dominates it, and \emph{inadmissible} otherwise.
\end{definition}

\begin{example}[Sample Mean and Sample Variance]
    Let $y_1,\dots,y_n$ be i.i.d.\ with mean $\mu$ and variance $\sigma^2$. The sample mean $\bar{y} = \frac{1}{n}\sum_i y_i$ satisfies
    \begin{align}
        \mathbb{E}[\bar{y}] = \mu, \qquad \mathrm{Var}(\bar{y}) = \frac{\sigma^2}{n}, \qquad \mathrm{se}(\bar{y}) = \frac{\sigma}{\sqrt{n}},
    \end{align}
    so it is unbiased, with error shrinking at the rate $n^{-1/2}$. This rate is the fundamental one in statistics: to halve a standard error one must \emph{quadruple} the sample.

    The variance estimator $\frac{1}{n}\sum_i (y_i - \bar{y})^2$ is biased downward, with expectation $\frac{n-1}{n}\sigma^2$, because the deviations are taken about the \emph{estimated} mean rather than the true one. Dividing by $n-1$ corrects this; the lost degree of freedom is spent estimating $\mu$.
\end{example}

\begin{definition}[Modes of Convergence]
    Let $\{X_n\}$ be a sequence of random variables and $X$ a random variable.
    \begin{itemize}[noitemsep]
        \item \emph{Almost surely}, $X_n \xrightarrow{a.s.} X$, if $\mathbb{P}\!\left(\lim_{n\to\infty} X_n = X\right) = 1$. The realized sequence/path settles down.
        \item \emph{In probability}, $X_n \xrightarrow{p} X$, if $\mathbb{P}(|X_n - X| > \varepsilon) \to 0$ for every $\varepsilon > 0$. For large $n$, the estimator is unlikely to be far from the target. 
        \item \emph{In $L^q$}, $X_n \xrightarrow{L^q} X$, if $\mathbb{E}[|X_n - X|^q] \to 0$. The case $q=2$ is convergence in mean square.
        \item \emph{In distribution}, $X_n \xrightarrow{d} X$, if $F_n(x) \to F(x)$ at every continuity point $x$ of $F$. The shape of the sampling distribution stabilizes.
    \end{itemize}
\end{definition}

\begin{proposition}[Hierarchy]
    \begin{align}
        X_n \xrightarrow{a.s.} X \;\Longrightarrow\; X_n \xrightarrow{p} X \;\Longrightarrow\; X_n \xrightarrow{d} X,
        \qquad
        X_n \xrightarrow{L^q} X \;\Longrightarrow\; X_n \xrightarrow{p} X .
    \end{align}
\end{proposition}

\begin{theorem}[Continuous Mapping]
    If $X_n \xrightarrow{d} X$ and $g$ is continuous on a set of $\mathbb{P}$-measure one under the law of $X$, then $g(X_n) \xrightarrow{d} g(X)$. The same holds with $\xrightarrow{d}$ replaced by $\xrightarrow{p}$ or $\xrightarrow{a.s.}$ throughout.
\end{theorem}

\begin{theorem}[Slutsky]
    If $X_n \xrightarrow{d} X$ and $Y_n \xrightarrow{p} c$ for a constant $c$, then
    \begin{align}
        X_n + Y_n \xrightarrow{d} X + c, \qquad Y_n X_n \xrightarrow{d} cX, \qquad X_n / Y_n \xrightarrow{d} X/c \;\;(c \neq 0).
    \end{align}
\end{theorem}

\begin{theorem}[Weak Law of Large Numbers]
    Let $X_1, X_2, \dots$ be i.i.d.\ with $\mathbb{E}[|X_1|] < \infty$ and $\mathbb{E}[X_1] = \mu$. Then
    \begin{align}
        \bar{X}_n = \frac{1}{n}\sum_{i=1}^n X_i \;\xrightarrow{p}\; \mu .
    \end{align}
\end{theorem}
\begin{proof}[Proof under the additional assumption $\mathrm{Var}(X_1) = \sigma^2 < \infty$]
    $\mathbb{E}[\bar{X}_n] = \mu$ and, by independence, $\mathrm{Var}(\bar{X}_n) = \sigma^2/n$. Chebyshev's inequality gives, for any $\varepsilon > 0$,
    \begin{align}
        \mathbb{P}\!\left(|\bar{X}_n - \mu| > \varepsilon\right) \leq \frac{\mathrm{Var}(\bar{X}_n)}{\varepsilon^2} = \frac{\sigma^2}{n\varepsilon^2} \;\longrightarrow\; 0.
    \end{align}
    The result under the weaker assumption $\mathbb{E}|X_1| < \infty$ requires a truncation argument.
\end{proof}

\begin{theorem}[Strong Law of Large Numbers (Kolmogorov)]
    Under the same hypotheses, $\bar{X}_n \xrightarrow{a.s.} \mu$. Conversely, if $\mathbb{E}[|X_1|] = \infty$ then $\limsup_n |\bar{X}_n| = \infty$ almost surely: the sample mean does not converge to anything.
\end{theorem}

\begin{figure}[h]
    \centering
    \includesvg[width=0.5\linewidth]{Lawoflargenumbers}
    \caption{Law of Large Numbers}
    \label{fig:placeholder}
\end{figure}

\begin{remark}[What the LLN Actually Requires]
    The only requirement is a \emph{finite first moment}. Finite variance is convenient (it gives the one-line Chebyshev proof and a rate) but is not necessary. What \emph{is} necessary is $\mathbb{E}|X| < \infty$.

    The canonical failure is the Cauchy distribution, whose density $\frac{1}{\pi(1+x^2)}$ has tails so heavy that $\mathbb{E}|X| = \infty$. For some assets and many intraday or high-frequency series, the variance is infinite and the LLN, while still valid, converges arbitrarily slowly. When a sample average refuses to stabilize as the window grows, the correct diagnosis is often not ``more data needed'' but ``this moment may not exist.''
\end{remark}

\begin{theorem}[Uniform Law of Large Numbers]
    Let $\Theta$ be compact, let $q(y;\theta)$ be continuous in $\theta$ for almost every $y$ and measurable in $y$ for each $\theta$, and suppose the dominance condition $\mathbb{E}[\sup_{\theta\in\Theta}|q(y;\theta)|] < \infty$ holds. Then, with $Q(\theta) := \mathbb{E}[q(y;\theta)]$,
    \begin{align}
        \sup_{\theta \in \Theta}\left|\frac{1}{n}\sum_{i=1}^n q(y_i;\theta) - Q(\theta)\right| \;\xrightarrow{p}\; 0,
    \end{align}
    and $Q$ is continuous on $\Theta$.
\end{theorem}

\begin{remark}[Why Uniformity Is Needed]
    The pointwise LLN gives convergence at each \emph{fixed} $\theta$. That is not enough to conclude anything about the \emph{maximizer} of the sample objective, because the maximizer is itself data-dependent and moves with $n$: pointwise convergence permits the sample objective to develop a spurious spike at some $\theta$ that drifts as the sample grows, so that $\hat{\theta}$ never settles. Uniform convergence rules this out by controlling the entire surface at once.
\end{remark}

The LLN says the sample mean converges. The CLT describes the distribution of the residual error.

\begin{theorem}[Lindeberg--L\'evy Central Limit Theorem]
    Let $X_1,X_2,\dots$ be i.i.d.\ with mean $\mu$ and \emph{finite} variance $\sigma^2 > 0$. Then
    \begin{align}
        \sqrt{n}\left(\bar{X}_n - \mu\right) \;\xrightarrow{d}\; \mathcal{N}(0,\sigma^2).
    \end{align}
\end{theorem}

\begin{theorem}[Multivariate CLT]
    Let $\vec{X}_1,\vec{X}_2,\dots$ be i.i.d.\ random vectors in $\mathbb{R}^k$ with mean $\vec{\mu}$ and finite covariance matrix $\Sigma$. Then
    \begin{align}
        \sqrt{n}\left(\bar{\vec{X}}_n - \vec{\mu}\right) \;\xrightarrow{d}\; \mathcal{N}_k(\vec{0}, \Sigma).
    \end{align}
\end{theorem}
\begin{proof}[Proof sketch]
    By the Cram\'er--Wold device, $\vec{Z}_n \xrightarrow{d} \vec{Z}$ in $\mathbb{R}^k$ if and only if $\vec{a}^T\vec{Z}_n \xrightarrow{d} \vec{a}^T\vec{Z}$ for every fixed $\vec{a} \in \mathbb{R}^k$. Each projection $\vec{a}^T\vec{X}_i$ is a scalar i.i.d.\ sequence with variance $\vec{a}^T\Sigma\vec{a} < \infty$, so the univariate CLT applies and gives $\mathcal{N}(0, \vec{a}^T\Sigma\vec{a})$, which is the law of $\vec{a}^T\vec{Z}$ for $\vec{Z} \sim \mathcal{N}_k(\vec{0},\Sigma)$.
\end{proof}

\begin{figure}[h]
    \centering
    \includegraphics[width=0.5\linewidth]{sample-mean-distribution.png}
    \caption{Sampling distribution of sample means converging to Gaussian}
    \label{fig:placeholder}
\end{figure}

\begin{theorem}[Delta Method]
    Suppose $\sqrt{n}(\hat{\theta}_n - \theta) \xrightarrow{d} \mathcal{N}_k(\vec{0}, V)$ and $g : \mathbb{R}^k \to \mathbb{R}^m$ is continuously differentiable at $\theta$ with Jacobian $G = \nabla_\theta g(\theta)$. Then
    \begin{align}
        \sqrt{n}\left(g(\hat{\theta}_n) - g(\theta)\right) \;\xrightarrow{d}\; \mathcal{N}_m\!\left(\vec{0},\; G V G^T\right).
    \end{align}
\end{theorem}
\begin{proof}[Proof sketch]
    A first-order Taylor expansion gives $g(\hat{\theta}_n) - g(\theta) = G(\hat{\theta}_n - \theta) + o_p(\|\hat{\theta}_n-\theta\|)$. Multiplying by $\sqrt{n}$, the remainder is $o_p(1)$ since $\sqrt{n}(\hat{\theta}_n-\theta) = O_p(1)$, and the result follows from the continuous mapping theorem applied to the linear map $G$, together with Slutsky.
\end{proof}

\begin{remark}[Using the Delta Method]
    The delta method is how standard errors are propagated through nonlinear transformations. First, it is a \emph{local linearization}: it is unreliable when $g$ is strongly curved over the region where $\hat{\theta}$ has appreciable mass, which in practice means when the standard error is large relative to the curvature scale. Second, it degenerates when $G = 0$ --- the first-order term vanishes and the limit is no longer Gaussian but a quadratic form in normals, i.e.\ a $\chi^2$-type limit at rate $n$ rather than $\sqrt{n}$.
\end{remark}

\begin{theorem}[Berry--Esseen]
    Let $X_i$ be i.i.d.\ with mean $\mu$, variance $\sigma^2$, and finite third absolute central moment $\rho = \mathbb{E}|X_1-\mu|^3$. Then
    \begin{align}
        \sup_{x \in \mathbb{R}} \left| \mathbb{P}\!\left(\frac{\sqrt{n}(\bar{X}_n - \mu)}{\sigma} \leq x\right) - \Phi(x) \right| \;\leq\; \frac{C\rho}{\sigma^3\sqrt{n}},
    \end{align}
    where $C$ is an absolute constant (the best known bound is below $0.48$).
\end{theorem}

\begin{remark}[The CLT Converges Slowly, and Not Where You Need It]
    Berry--Esseen quantifies the CLT's accuracy, and delivers two pieces of bad news for risk work.

    First, the error decays only at rate $n^{-1/2}$, with a constant proportional to $\rho/\sigma^3$ --- essentially the standardized third moment. For strongly skewed data, the effective sample size required for a good normal approximation can be an order of magnitude larger than intuition suggests. 

    Second, and more seriously, Berry--Esseen bounds the \emph{uniform absolute} error. That is the wrong metric for tail risk. In the tails, both the true probability and the Gaussian approximation are small, so a small absolute error can be a catastrophic \emph{relative} error. The precise statement (Cram\'er's moderate-deviation theorem) is that
    \begin{align}
        \frac{\mathbb{P}\left(\sqrt{n}(\bar{X}_n-\mu)/\sigma > x\right)}{1 - \Phi(x)} \to 1
        \qquad \text{only for } x = o(n^{1/6}).
    \end{align}
    Beyond that range, the sum's tail is governed by the tail of the summands themselves, not by the Gaussian limit. The CLT is thus a theorem about the centre of the distribution, and licenses nothing about the tails.
\end{remark}

\begin{remark}[Generalized CLT in Extreme Value Theory]
    When the variance is infinite, the Gaussian limit fails, but a limit theorem survives. If the $X_i$ lie in the domain of attraction of an $\alpha$-stable law with tail index $\alpha \in (0,2)$ --- i.e.\ $\mathbb{P}(|X| > x) \sim Lx^{-\alpha}$ --- then for suitable centering $a_n$,
    \begin{align}
        n^{-1/\alpha}\left(\sum_{i=1}^n X_i - a_n\right) \;\xrightarrow{d}\; S_\alpha,
    \end{align}
    where $S_\alpha$ is $\alpha$-stable. The convergence rate is $n^{-1/\alpha}$, which is \emph{slower} than $n^{-1/2}$ and slows further as tails fatten; and the limit retains the heavy tails of the summands rather than washing them out. Aggregation then reproduces the stable law, not normality.
\end{remark}

\begin{remark}[Dependence]
    Both limit theorems above assume independence, which financial time series violate. The results extend, with conditions:
    \begin{itemize}[noitemsep]
        \item \textbf{Martingale difference CLT.} If $\{X_t,\mathcal{F}_t\}$ is a martingale difference sequence ($\mathbb{E}[X_t|\mathcal{F}_{t-1}]=0$) with suitable moment and stability conditions, the CLT holds with the same $\sqrt{n}$ rate. This is the relevant version for asset pricing, since unpredictable excess returns are martingale differences by construction.
        \item \textbf{Mixing processes.} For stationary sequences with sufficiently fast decay of dependence, the CLT holds but the asymptotic variance becomes the \emph{long-run variance}
        \begin{align}
            \Omega = \sum_{j=-\infty}^{\infty} \mathrm{Cov}(X_t, X_{t-j}) \;\neq\; \mathrm{Var}(X_t),
        \end{align}
        which must be estimated by a HAC (Newey--West) estimator.
        \item \textbf{Long memory.} If autocovariances decay so slowly that $\sum_j |\mathrm{Cov}(X_t,X_{t-j})| = \infty$, the $\sqrt{n}$ rate itself fails and is replaced by $n^{H}$ for a Hurst exponent $H > 1/2$.
    \end{itemize}
\end{remark}


\begin{definition}[$t$-Statistic]
    Let $\hat{\theta}$ be an estimator of $\theta$ with estimated standard error $\widehat{\mathrm{se}}(\hat{\theta})$. The \emph{$t$-statistic} for the null hypothesis $H_0 : \theta = \theta_0$ is
    \begin{align}
        t = \frac{\hat{\theta} - \theta_0}{\widehat{\mathrm{se}}(\hat{\theta})}.
    \end{align}
    The near-universal convention $|t| > 2$ is the approximate $97.5$th percentile of the standard normal, corresponding to a two-sided test at the $5\%$ level.
\end{definition}

Substituting $\hat{\sigma}$ for the unknown $\sigma$ ought, on its face, to be a problem: the statistic now contains two sources of randomness rather than one, and its distribution ought to depend on the very quantity $\sigma$ that was eliminated. Remarkably, it does not, as will be shown.

\begin{definition}[Chi-Square Distribution]
    If $W_1,\dots,W_r$ are i.i.d.\ $\mathcal{N}(0,1)$, then $V = \sum_{i=1}^r W_i^2$ has the \emph{chi-square distribution with $r$ degrees of freedom}, written $V \sim \chi^2_r$, with density
    \begin{align}
        f_V(v) = \frac{1}{2^{r/2}\Gamma(r/2)}\,v^{r/2 - 1}e^{-v/2}, \qquad v > 0.
    \end{align}
\end{definition}

\begin{proposition}[Gaussian Projection]\label{prop:proj}
    Let $\vec{\varepsilon} \sim \mathcal{N}_n(\vec{0}, \sigma^2 I_n)$ and let $P$ be an orthogonal projection matrix (idempotent and self-adjoint) of rank $r$. Then
    \begin{enumerate}[noitemsep]
        \item $\displaystyle \frac{\|P\vec{\varepsilon}\|^2}{\sigma^2} \sim \chi^2_r$;
        \item for any fixed vector $\vec{c}$ with $P\vec{c} = \vec{0}$, the scalar $\vec{c}^T\vec{\varepsilon}$ and the vector $P\vec{\varepsilon}$ are independent.
    \end{enumerate}
\end{proposition}
\begin{proof}
    For the first claim, by the spectral theorem, write $P = U\Lambda U^T$ with $U$ orthogonal and $\Lambda = \mathrm{diag}(1,\dots,1,0,\dots,0)$ carrying exactly $r$ ones, since the eigenvalues of an idempotent symmetric matrix are $0$ and $1$ with multiplicity $r = \mathrm{rank}(P)$. The standard normal law is rotationally invariant: $\vec{W} = U^T\vec{\varepsilon}/\sigma \sim \mathcal{N}_n(\vec{0}, I_n)$, since $\mathrm{Var}(U^T\vec{\varepsilon}) = U^T(\sigma^2 I)U = \sigma^2 I$. Then
    \begin{align}
        \frac{\|P\vec{\varepsilon}\|^2}{\sigma^2} = \frac{\vec{\varepsilon}^TP\vec{\varepsilon}}{\sigma^2} = \vec{W}^T\Lambda\vec{W} = \sum_{i=1}^r W_i^2 \sim \chi^2_r.
    \end{align}
    For the second claim, the pair $(\vec{c}^T\vec{\varepsilon}, P\vec{\varepsilon})$ is jointly Gaussian, being a linear image of $\vec{\varepsilon}$. Their covariance is
    \begin{align}
        \mathrm{Cov}\!\left(\vec{c}^T\vec{\varepsilon},\, P\vec{\varepsilon}\right) = \mathbb{E}\!\left[P\vec{\varepsilon}\,\vec{\varepsilon}^T\vec{c}\right] = P(\sigma^2 I)\vec{c} = \sigma^2 P\vec{c} = \vec{0}.
    \end{align}
    For jointly Gaussian variables, zero covariance implies independence.
\end{proof}

\begin{theorem}[Student's $t$ Distribution]\label{thm:t-density}
    Let $Z \sim \mathcal{N}(0,1)$ and $V \sim \chi^2_\nu$ be independent. Then
    \begin{align}
        T = \frac{Z}{\sqrt{V/\nu}}
    \end{align}
    has density
    \begin{align}
        f_\nu(t) = \frac{\Gamma\!\left(\frac{\nu+1}{2}\right)}{\sqrt{\nu\pi}\,\Gamma\!\left(\frac{\nu}{2}\right)}\left(1 + \frac{t^2}{\nu}\right)^{-\frac{\nu+1}{2}}, \qquad t \in \mathbb{R},
    \end{align}
    called \emph{Student's $t$ distribution with $\nu$ degrees of freedom}, written $T \sim t_\nu$.
\end{theorem}

\begin{figure}[h]
    \centering
    \includegraphics[width=0.5\linewidth]{Student_t_pdf.svg.png}
    \caption{t-distribution with different degrees of freedom}
    \label{fig:placeholder}
\end{figure}

\begin{theorem}[Distribution of a $t$-Statistic]\label{thm:general-t}
    Let $\vec{Y} \sim \mathcal{N}_n(\vec{\mu}, \sigma^2 I_n)$ with $\sigma^2 > 0$ unknown. Suppose there is a linear estimator $\hat{\theta} = \vec{c}^T\vec{Y}$ for fixed $\vec{c} \neq \vec{0}$, estimating $\theta = \vec{c}^T\vec{\mu}$; a variance estimator $\hat{\sigma}^2 = \|P\vec{Y}\|^2/r$, where $P$ is an orthogonal projection of rank $r \geq 1$ satisfying $P\vec{\mu} = \vec{0}$; and an orthogonality condition $P\vec{c} = \vec{0}$.

    Then $\hat{\sigma}^2$ is unbiased for $\sigma^2$, and
    \begin{align}
        T \;=\; \frac{\hat{\theta} - \theta}{\hat{\sigma}\,\|\vec{c}\|} \;\sim\; t_{\,r}.
    \end{align}
    In particular, the distribution of $T$ does not depend on $\sigma$, on $\vec{\mu}$, or on $n$.
\end{theorem}
\begin{proof}
    Write $\vec{\varepsilon} = \vec{Y} - \vec{\mu} \sim \mathcal{N}_n(\vec{0},\sigma^2 I_n)$. By linearity, $\hat{\theta} - \theta = \vec{c}^T\vec{Y} - \vec{c}^T\vec{\mu} = \vec{c}^T\vec{\varepsilon}$, which is Gaussian with mean zero and variance $\sigma^2\|\vec{c}\|^2$. Hence, call
    \begin{align}
        Z = \frac{\vec{c}^T\vec{\varepsilon}}{\sigma\|\vec{c}\|} \sim \mathcal{N}(0,1).
    \end{align}
    The variance estimator assumption gives $P\vec{Y} = P\vec{\mu} + P\vec{\varepsilon} = P\vec{\varepsilon}$, so the variance estimator sees only noise. By Proposition~\ref{prop:proj}, call
    \begin{align}
        V = \frac{\|P\vec{\varepsilon}\|^2}{\sigma^2} = \frac{r\hat{\sigma}^2}{\sigma^2} \sim \chi^2_r,
    \end{align}
    and since $\mathbb{E}[\chi^2_r] = r$, we get $\mathbb{E}[\hat{\sigma}^2] = \sigma^2$, the claimed unbiasedness.

    The orthogonality assumption and Proposition~\ref{prop:proj} give $\vec{c}^T\vec{\varepsilon} \perp P\vec{\varepsilon}$, hence $Z \perp V$. Now assemble T, multiplying numerator and denominator by $1/\sigma$:
    \begin{align}
        T = \frac{\hat{\theta} - \theta}{\hat{\sigma}\,\|\vec{c}\|} =\frac{\vec{c}^T\vec{\varepsilon}}{\|\vec{c}\|\sqrt{\|P\vec{\varepsilon}\|^2/r}}
          = \frac{\vec{c}^T\vec{\varepsilon}\,/\,(\sigma\|\vec{c}\|)}{\sqrt{\|P\vec{\varepsilon}\|^2/(\sigma^2 r)}}
          = \frac{Z}{\sqrt{V/r}}.
    \end{align}
    The unknown $\sigma$ has cancelled identically. By Theorem~\ref{thm:t-density}, $T \sim t_r$.
\end{proof}

\begin{corollary}[One-Sample $t$-Test]
    Let $Y_1,\dots,Y_n$ be i.i.d.\ $\mathcal{N}(\mu,\sigma^2)$, and let $S^2 = \frac{1}{n-1}\sum_i (Y_i - \bar{Y})^2$. Then
    \begin{align}
        \frac{\sqrt{n}\left(\bar{Y}-\mu\right)}{S} \sim t_{\,n-1}.
    \end{align}
\end{corollary}
\begin{proof}
    Apply Theorem~\ref{thm:general-t} with $\vec{c} = \vec{1}/n$, so $\hat{\theta} = \bar{Y}$, $\theta = \mu$ and $\|\vec{c}\| = 1/\sqrt{n}$; and with $P = I_n - \frac{1}{n}\vec{1}\vec{1}^T$, the centering projection, of rank $n-1$. Then $\vec{\mu} = \mu\vec{1}$ gives $P\vec{\mu} = \vec{0}$, and $P\vec{c} = \frac{1}{n}P\vec{1} = \vec{0}$, so both hypotheses hold. Finally $\|P\vec{Y}\|^2 = \sum_i (Y_i-\bar{Y})^2 = (n-1)S^2$, so $\hat{\sigma}^2 = S^2$ and $T = (\bar{Y}-\mu)/(S/\sqrt{n})$.
\end{proof}

\begin{remark}[Necessity of Normality]
    It is tempting to read the Gaussian assumption as a technical simplification that a more careful argument could remove. The independence of $\bar{Y}$ and $S^2$, which is what makes the ratio a $t$ rather than something intractable, \emph{characterizes} the normal distribution: for an i.i.d.\ sample, the sample mean and sample variance are independent if and only if the population is normal (Geary, 1936; Lukacs, 1942). The rotational-invariance step in Proposition~\ref{prop:proj} is where this enters: the isotropic Gaussian is the unique distribution with independent components that is invariant under all rotations, so it is the unique law for which projections onto orthogonal subspaces are independent.
\end{remark}

\begin{remark}[Recovering the Normal]
    The excess tail weight relative to the normal has a single source: the denominator is estimated. Occasionally $\hat{\sigma}$ comes out too small, inflating $|T|$, and that additional randomness fattens the tails. Formally, since $(1+t^2/\nu)^{-(\nu+1)/2} \to e^{-t^2/2}$ as $\nu\to\infty$ and the normalizing constant converges to $1/\sqrt{2\pi}$ by Stirling,
    \begin{align}
        f_\nu(t) \;\longrightarrow\; \frac{1}{\sqrt{2\pi}}e^{-t^2/2},
    \end{align}
    so $t_\nu \Rightarrow \mathcal{N}(0,1)$. This is the exact-distribution counterpart of the Slutsky argument: as $r\to\infty$, $\hat{\sigma}\xrightarrow{p}\sigma$ and the estimated denominator behaves like a constant.
\end{remark}

\begin{remark}[What Survives Without Normality]
    Drop normality and the exact result collapses entirely: $\hat{\theta}$ is no longer Gaussian, $\hat{\sigma}^2$ is no longer chi-square, and the two are no longer independent. What survives is the asymptotic statement --- by the CLT the numerator is asymptotically normal, by the LLN the denominator is consistent, and Slutsky combines them to give $T \xrightarrow{d}\mathcal{N}(0,1)$. \textbf{The finite-sample justification is lost; the large-sample one is not.}

    How large is ``large'' depends on the shape of the underlying distribution, and by Berry--Esseen the governing quantity is the standardized third moment. For the skewed distributions endemic to finance, the normal approximation to a $t$-statistic can be poor at sample sizes where one would expect it to be excellent, and it errs in a specific direction: with left-skewed returns the $t$-statistic is right-skewed, so \emph{positive} findings are over-represented at any fixed threshold.

    Because $T$ is a \emph{self-normalized} sum --- the same data appear in numerator and denominator --- it is far better behaved than the raw sum. Gin\'e, G\"otze, and Mason (1997) proved that the self-normalized sum is asymptotically standard normal \emph{if and only if} the summands lie in the domain of attraction of the normal law, and self-normalized sums admit Gaussian-type tail bounds under weaker moment conditions than raw sums do. Studentizing therefore buys robustness as well as pivotality. It does not, however, rescue the case of genuinely infinite variance, where the limit is not normal at all.
\end{remark}

\begin{proposition}[Duality of Tests and Confidence Intervals]
    The set of null values not rejected by a two-sided $t$-test at level $\alpha$ is exactly the $(1-\alpha)$ confidence interval:
    \begin{align}
        \left\{\theta_0 : \left|\frac{\hat{\theta}-\theta_0}{\widehat{\mathrm{se}}}\right| \leq c_{1-\alpha/2}\right\}
        = \left[\hat{\theta} - c_{1-\alpha/2}\,\widehat{\mathrm{se}},\;\; \hat{\theta} + c_{1-\alpha/2}\,\widehat{\mathrm{se}}\right].
    \end{align}
\end{proposition}

\begin{proposition}[The $t$-Statistic of a Sharpe Ratio]\label{prop:sharpe-t}
    Let a strategy have i.i.d.\ excess returns with mean $\mu$ and volatility $\sigma$, observed over $T$ years, and let $S = \mu/\sigma$ be its annualized Sharpe ratio. The $t$-statistic for $H_0 : \mu = 0$ is
    \begin{align}
        t = \frac{\hat{\mu}}{\hat{\sigma}/\sqrt{T}} = \hat{S}\,\sqrt{T}.
    \end{align}
\end{proposition}
\begin{proof}
    Immediate: the standard error of the mean is $\hat{\sigma}/\sqrt{T}$, so $t = \hat{\mu}\sqrt{T}/\hat{\sigma} = \hat{S}\sqrt{T}$. Note the result is invariant to observation frequency --- sampling daily rather than monthly changes both $\hat\mu$ and $\hat\sigma$ by the same $\sqrt{\text{scale}}$ factor --- so only the calendar span $T$ matters.
\end{proof}

\begin{remark}[The Backtest Length Constraint]
    Proposition~\ref{prop:sharpe-t} is the most useful single identity in empirical finance, and deserves to be memorized. It says the $t$-statistic of a track record is the Sharpe ratio times the square root of its length in years. Rearranged, the years required to establish a given Sharpe at a given hurdle are
    \begin{align}
        T = \left(\frac{t^{*}}{S}\right)^{2}.
    \end{align}
    A genuine Sharpe-$0.5$ strategy cannot be statistically distinguished from noise at the conventional threshold with less than sixteen years of live data. Most backtests are far shorter than the horizon needed to validate the effects they claim to find. Consequently, when a five-year backtest \emph{does} produce $t = 3$, the most probable explanations are, in order: the sample was selected after the fact, the strategy was chosen from among many, or the returns are not i.i.d.\ --- not that a Sharpe-$1.3$ strategy has been discovered.
\end{remark}

\begin{remark}[Multiple Testing: Why the Hurdle Is Not 2]
    The $|t| > 2$ convention controls the false positive rate for \emph{one} pre-specified hypothesis. It provides no protection whatever when a hypothesis is selected from many, which is the actual practice in factor research and strategy development.

    If $N$ independent strategies with \emph{true} Sharpe of zero are tested, their $t$-statistics are approximately i.i.d.\ standard normal, and the expected maximum satisfies
    \begin{align}
        \mathbb{E}\left[\max_{1\leq i\leq N} t_i\right] \;\approx\; \sqrt{2\ln N}.
    \end{align}
    For $N = 100$ this is $3.03$; for $N = 1000$ it is $3.72$. In other words, testing a hundred worthless signals is expected to produce one with $t \approx 3$.

    Several formal corrections exist:
    \begin{itemize}[noitemsep]
        \item \textbf{Bonferroni}: test each at level $\alpha/N$. Controls the family-wise error rate; very conservative when tests are correlated, as factor tests invariably are.
        \item \textbf{Benjamini--Hochberg}: controls the \emph{false discovery rate}, the expected proportion of rejections that are false. Usually the more sensible target in research settings, where some false positives are tolerable but their fraction must be bounded.
        \item \textbf{Deflated Sharpe ratio} (Bailey and L\'opez de Prado): adjusts an observed Sharpe for the number of configurations tried, the length of the sample, and the skew and kurtosis of returns, returning the probability the true Sharpe exceeds zero.
    \end{itemize}
\end{remark}

\begin{remark}[A $t$-Statistic Is Not an Effect Size]
    Since $t = \hat{S}\sqrt{T}$ and more generally $t \propto \sqrt{n}$, significance can always be manufactured with enough data. In panel return prediction, where $n$ runs to millions of stock-months, effects far too small to survive transaction costs will register enormous $t$-statistics. Statistical significance answers ``is this distinguishable from zero,'' which is rarely the question of interest; the question of interest is ``is this large enough to trade after costs,'' and no $t$-statistic addresses it.
\end{remark}

\begin{remark}[The $t$-Statistic Inherits Assumption in Its Denominator]
    A $t$-statistic is only as trustworthy as the standard error beneath it, and inflated $t$-statistics almost always originate in an understated denominator rather than an overstated numerator.
    \begin{itemize}[noitemsep]
        \item \textbf{Serial correlation.} Using i.i.d.\ standard errors on persistent series understates the long-run variance and inflates $t$. HAC (Newey--West) standard errors are required, as discussed above. This is the single most common defect in published backtests of slow-moving signals.
        \item \textbf{Overlapping observations.} Regressing $k$-period returns sampled every period induces an MA($k-1$) structure by construction. The effective sample size is roughly $T/k$, so naive standard errors are too small by a factor of about $\sqrt{k}$ --- which is how long-horizon predictability regressions produce spectacular $t$-statistics from very little independent information.
        \item \textbf{Cross-sectional correlation.} Pooled regressions on stock-level panels treat correlated residuals as independent, overstating the effective number of observations by orders of magnitude. Fama--MacBeth (estimating cross-sectional regressions period by period, then computing the $t$-statistic from the \emph{time series} of coefficients) or clustered standard errors correct for this.
        \item \textbf{Misspecification.} As shown in the quasi-MLE discussion below, the information-matrix-based standard error is wrong whenever the model is misspecified; the sandwich form is required.
    \end{itemize}
\end{remark}

\section{Maximum Likelihood}

\begin{definition}[Likelihood and Score]
    Let $\{y_i\}_{i=1}^n$ be an i.i.d.\ sample from a parametric family $\{p(y;\theta) : \theta \in \Theta \subseteq \mathbb{R}^k\}$, and let $\theta_0$ denote the \emph{true} parameter, so that the data are actually drawn from $p(\cdot\,;\theta_0)$. The \emph{log-likelihood} is
    \begin{align}
        \ell(\theta) = \sum_{i=1}^n \log p(y_i;\theta).
    \end{align}
    The \emph{score function} is the gradient of the log-likelihood,
    \begin{align}
        s(\theta) = \nabla_\theta \ell(\theta) = \sum_{i=1}^n \nabla_\theta \log p(y_i;\theta).
    \end{align}
    The \emph{maximum likelihood estimator} (MLE) is $\hat{\theta}_{MLE} = \arg\max_{\theta \in \Theta} \ell(\theta)$, which satisfies the first-order condition $s(\hat{\theta}) = 0$ at an interior solution.
\end{definition}

\begin{definition}[Identification]
    The parameter is \emph{identified} if distinct parameters induce distinct distributions:
    \begin{align}
        \theta \neq \theta_0 \quad\Longrightarrow\quad p(\cdot\,;\theta) \neq p(\cdot\,;\theta_0) \text{ on a set of positive measure.}
    \end{align}
\end{definition}

\begin{definition}[Large-Sample Properties of an Estimator]
    Let $\hat{\theta}_n$ be an estimator of $\theta_0$ computed from a sample of size $n$.
    \begin{itemize}[noitemsep]
        \item $\hat{\theta}_n$ is \emph{consistent} if $\hat{\theta}_n \xrightarrow{p} \theta_0$: for every $\varepsilon > 0$, $\mathbb{P}(\|\hat{\theta}_n - \theta_0\| > \varepsilon) \to 0$. This is a statement about \emph{eventual correctness}, and is the minimum one should demand of an estimator. Note that consistency is neither implied by nor implies unbiasedness: an estimator can be unbiased at every $n$ yet never concentrate (its variance need not shrink), and it can be biased at every $n$ yet be consistent (if the bias vanishes asymptotically).
        \item $\hat{\theta}_n$ is \emph{asymptotically normal} at rate $\sqrt{n}$ if $\sqrt{n}(\hat{\theta}_n - \theta_0) \xrightarrow{d} \mathcal{N}(0, V)$ for some positive-definite $V$, the \emph{asymptotic variance}. This is what licenses confidence intervals and $t$-statistics; the $\sqrt{n}$ rate says the sampling error shrinks like $n^{-1/2}$.
        \item $\hat{\theta}_n$ is \emph{asymptotically efficient} if it is asymptotically normal with the smallest attainable $V$ among a suitable class of competing estimators. The content of the theory below is that this floor is $V = \mathcal{I}(\theta_0)^{-1}$, and that the MLE attains it.
    \end{itemize}
\end{definition}

The MLE maximizes $\frac{1}{n}\ell(\theta)$, a sample average. By the law of large numbers this converges to a population objective, and the question of consistency reduces to asking whether that population objective is maximized at the truth. It is, and the reason is information-theoretic.

\begin{proposition}[The Population Log-Likelihood Is Maximized at the Truth]\label{prop:kl-max}
    For every $\theta \in \Theta$,
    \begin{align}
        \mathbb{E}_{\theta_0}\!\left[\log p(y;\theta)\right] \;\leq\; \mathbb{E}_{\theta_0}\!\left[\log p(y;\theta_0)\right],
    \end{align}
    with equality if and only if $p(\cdot\,;\theta) = p(\cdot\,;\theta_0)$ almost everywhere. Equivalently, the gap between the two sides is exactly the Kullback--Leibler divergence,
    \begin{align}
        \mathbb{E}_{\theta_0}\!\left[\log p(y;\theta_0)\right] - \mathbb{E}_{\theta_0}\!\left[\log p(y;\theta)\right] = D_{\mathrm{KL}}\!\left(p_{\theta_0} \,\|\, p_\theta\right) \geq 0.
    \end{align}
\end{proposition}
\begin{proof}
    Consider the expected log-likelihood ratio and apply Jensen's inequality, using the concavity of $\log$:
    \begin{align}
        \mathbb{E}_{\theta_0}\!\left[\log \frac{p(y;\theta)}{p(y;\theta_0)}\right]
        \;\leq\; \log \mathbb{E}_{\theta_0}\!\left[\frac{p(y;\theta)}{p(y;\theta_0)}\right]
        = \log \int \frac{p(y;\theta)}{p(y;\theta_0)}\,p(y;\theta_0)\,dy
        = \log \int p(y;\theta)\,dy = \log 1 = 0.
    \end{align}
    Rearranging gives the inequality. Since $\log$ is \emph{strictly} concave, equality in Jensen holds if and only if the ratio $p(y;\theta)/p(y;\theta_0)$ is almost surely constant; as it integrates to one against $p_{\theta_0}$, that constant must be $1$.
\end{proof}

\begin{theorem}[Consistency of the MLE]
    Suppose (i) $\theta_0$ is identified; (ii) $\Theta$ is compact; (iii) $\log p(y;\theta)$ is continuous in $\theta$ for almost every $y$; and (iv) a dominance condition holds, $\mathbb{E}_{\theta_0}[\sup_{\theta\in\Theta}|\log p(y;\theta)|] < \infty$, which together with (ii) and (iii) delivers the uniform law of large numbers
    \begin{align}
        \sup_{\theta \in \Theta}\left| \tfrac{1}{n}\ell(\theta) - Q(\theta) \right| \xrightarrow{p} 0, \qquad Q(\theta) := \mathbb{E}_{\theta_0}[\log p(y;\theta)].
    \end{align}
    Then $\hat{\theta}_{MLE} \xrightarrow{p} \theta_0$.
\end{theorem}
\begin{proof}
    By Proposition~\ref{prop:kl-max} together with identification, $Q$ is \emph{uniquely} maximized at $\theta_0$. Fix $\varepsilon > 0$ and let $\mathcal{N}_\varepsilon = \{\theta \in \Theta : \|\theta - \theta_0\| \geq \varepsilon\}$, which is compact. Since $Q$ is continuous with a unique maximum at $\theta_0$,
    \begin{align}
        \delta = Q(\theta_0) - \sup_{\theta \in \mathcal{N}_\varepsilon} Q(\theta) > 0.
    \end{align}
    Write $Q_n(\theta) = \frac{1}{n}\ell(\theta)$. On the event $\{\sup_\theta |Q_n(\theta) - Q(\theta)| < \delta/2\}$, which has probability tending to one by the uniform LLN, we have for any $\theta \in \mathcal{N}_\varepsilon$:
    \begin{align}
        Q_n(\theta) < Q(\theta) + \tfrac{\delta}{2} \leq Q(\theta_0) - \delta + \tfrac{\delta}{2} = Q(\theta_0) - \tfrac{\delta}{2} < Q_n(\theta_0).
    \end{align}
    So every point at distance $\geq \varepsilon$ from $\theta_0$ is strictly beaten by $\theta_0$ itself, and the maximizer $\hat{\theta}$ cannot lie in $\mathcal{N}_\varepsilon$. Hence $\mathbb{P}(\|\hat{\theta}-\theta_0\| \geq \varepsilon) \to 0$.
\end{proof}

\begin{remark}[MLE as Minimum-Divergence Estimation]
    Proposition~\ref{prop:kl-max} says something stronger than ``the MLE works.'' Maximizing the likelihood is \emph{identically} minimizing an estimate of $D_{\mathrm{KL}}(p_{\theta_0}\|p_\theta)$: the MLE is the KL projection of the truth onto the model family. This single observation unifies several results elsewhere in these notes:
    \begin{itemize}[noitemsep]
        \item Under Gaussian noise, minimizing KL is minimizing squared error, so OLS is the MLE (Section~2).
        \item Under a Bernoulli model, minimizing KL is minimizing cross-entropy, so logistic regression is the MLE (Section~2).
        \item Training a language model by minimizing cross-entropy is minimizing the KL divergence from the true conditional distribution over tokens, i.e.\ it is maximum likelihood (Section~11).
        \item When the model is \emph{misspecified} --- the truth is not in the family at all --- the KL projection still exists and is still what the estimator finds. This is the pseudo-true parameter discussed below.
    \end{itemize}
\end{remark}

Consistency says the estimator eventually gets there. Efficiency asks how fast, and whether anything could be faster. The answer is governed by the curvature of the log-likelihood.

\begin{definition}[Fisher Information]
    The \emph{Fisher information matrix} at $\theta$ is
    \begin{align}
        \mathcal{I}(\theta) = \mathbb{E}\!\left[\left(\nabla_\theta \log p(y;\theta)\right)\left(\nabla_\theta \log p(y;\theta)\right)^T\right] = -\mathbb{E}\!\left[\nabla^2_\theta \log p(y;\theta)\right],
    \end{align}
    where the equality between the two expressions --- the \emph{information matrix equality} --- holds under regularity conditions permitting differentiation under the integral sign, and requires that the model be correctly specified. The Fisher information measures the curvature of the expected log-likelihood: high curvature means the data are highly informative about $\theta$, so the likelihood surface has a sharp peak and $\hat{\theta}$ is precisely estimated. A flat likelihood, by contrast, means many parameter values explain the data nearly equally well, and $\hat{\theta}$ is correspondingly imprecise.
\end{definition}

\begin{proof}[Proof of the information matrix equality]
    Since $\int p(y;\theta)\,dy = 1$ for every $\theta$, differentiating with respect to $\theta$ and exchanging differentiation and integration gives
    \begin{align}
        0 = \int \nabla_\theta p(y;\theta)\,dy = \int p(y;\theta)\,\nabla_\theta\log p(y;\theta)\,dy,
    \end{align}
    using $\nabla_\theta p = p\,\nabla_\theta \log p$. Hence $\mathbb{E}[\nabla_\theta\log p(y;\theta)] = 0$: the score has mean zero at the true $\theta$, so $\mathcal{I}(\theta)$ is also the \emph{variance} of the score. Differentiating this identity once more,
    \begin{align}
        0 = \int \nabla_\theta\!\left[p(y;\theta)\nabla_\theta\log p(y;\theta)\right]dy = \mathbb{E}\!\left[\nabla^2_\theta\log p(y;\theta)\right] + \mathbb{E}\!\left[\nabla_\theta\log p(y;\theta)\,\nabla_\theta\log p(y;\theta)^T\right],
    \end{align}
    which rearranges to the stated equality. Both steps require that the support of $p(\cdot\,;\theta)$ not depend on $\theta$, so that the limits of integration are fixed.
\end{proof}

% [CORRECTED] Statement now matches the proof (the 1/n is required, since
% I(theta) is defined PER OBSERVATION above). The Cauchy-Schwarz step also
% needed fixing: one must project the score onto a second vector b and then
% CHOOSE b = I^{-1} a. That choice is the entire trick and was missing.
\begin{theorem}[Cramér--Rao Lower Bound]
    Let $\tilde{\theta}$ be any unbiased estimator of $\theta$ based on an i.i.d.\ sample of size $n$, and assume the regularity conditions above (in particular, that the support of $p(\cdot\,;\theta)$ does not depend on $\theta$). Then
    \begin{align}
        \mathrm{Var}(\tilde{\theta}) \;\geq\; \frac{1}{n}\,\mathcal{I}(\theta)^{-1},
    \end{align}
    in the sense that $\mathrm{Var}(\tilde{\theta}) - \frac{1}{n}\mathcal{I}(\theta)^{-1}$ is positive semi-definite, where $\mathcal{I}(\theta)$ is the \emph{per-observation} Fisher information defined above. No unbiased estimator can have variance smaller than this bound.
\end{theorem}
\begin{proof}
    Write $s(\theta) = \sum_{i=1}^n \nabla_\theta \log p(y_i;\theta)$ for the total score, which has mean zero and variance $\mathrm{Var}(s(\theta)) = n\mathcal{I}(\theta)$. Unbiasedness means $\mathbb{E}[\tilde{\theta}] = \theta$ for all $\theta$; differentiating under the integral sign gives the key covariance identity
    \begin{align}
        \mathrm{Cov}(\tilde{\theta}, s(\theta)) = I_k.
    \end{align}
    Fix arbitrary $\vec{a}, \vec{b} \in \mathbb{R}^k$ and consider the two scalars $\vec{a}\cdot\tilde{\theta}$ and $\vec{b}\cdot s(\theta)$. Their covariance is $\vec{a}^TI_k\vec{b} = \vec{a}^T\vec{b}$, so Cauchy--Schwarz gives
    \begin{align}
        (\vec{a}^T\vec{b})^2 \;\leq\; \mathrm{Var}(\vec{a}^T\tilde{\theta})\cdot \mathrm{Var}(\vec{b}^Ts(\theta)) \;=\; \left(\vec{a}^T\mathrm{Var}(\tilde{\theta})\vec{a}\right)\cdot n\,\vec{b}^T\mathcal{I}(\theta)\vec{b}.
    \end{align}
    This holds for \emph{every} $\vec{b}$, so we are free to choose the one that makes the bound tightest. Taking $\vec{b} = \mathcal{I}(\theta)^{-1}\vec{a}$ yields $\vec{a}^T\vec{b} = \vec{a}^T\mathcal{I}^{-1}\vec{a}$ and $\vec{b}^T\mathcal{I}\vec{b} = \vec{a}^T\mathcal{I}^{-1}\vec{a}$, hence
    \begin{align}
        \left(\vec{a}^T\mathcal{I}^{-1}\vec{a}\right)^2 \leq \left(\vec{a}^T\mathrm{Var}(\tilde{\theta})\vec{a}\right)\cdot n\,\vec{a}^T\mathcal{I}^{-1}\vec{a}.
    \end{align}
    Dividing by $n\,\vec{a}^T\mathcal{I}^{-1}\vec{a} > 0$ gives $\vec{a}^T\mathrm{Var}(\tilde{\theta})\vec{a} \geq \frac{1}{n}\vec{a}^T\mathcal{I}(\theta)^{-1}\vec{a}$. Since $\vec{a}$ was arbitrary, the matrix difference is positive semi-definite.
\end{proof}

\begin{remark}[The Support Condition]
    The requirement that the support of $p(\cdot\,;\theta)$ not depend on $\theta$ is important. Take $y_i \sim \mathrm{Unif}[0,\theta]$. The estimator $\tilde{\theta} = \frac{n+1}{n}\max_i y_i$ is unbiased with variance of order $\theta^2/n^2$, which converges at rate $n^{-1}$ rather than $n^{-1/2}$ and violates any naive application of the bound. Nothing is wrong: the support $[0,\theta]$ moves with $\theta$, so differentiation under the integral sign is invalid and the theorem simply does not apply. Models with parameter-dependent support fall outside the standard theory and need separate treatment.
\end{remark}

\begin{remark}[Efficiency Is Not the Same as Low Error]
    The Cramér--Rao bound constrains \emph{unbiased} estimators only, and that restriction is severe. Since
    \begin{align}
        \mathrm{MSE}(\hat{\theta}) = \mathrm{Bias}(\hat{\theta})^2 + \mathrm{Var}(\hat{\theta}),
    \end{align}
    an estimator willing to accept bias can have strictly lower mean squared error than \emph{any} unbiased estimator, the Cramér--Rao bound notwithstanding.

    The sharpest statement of this is the James--Stein phenomenon. For estimating the mean $\vec{\mu}$ of a $k$-dimensional Gaussian with $k \geq 3$, the MLE $\hat{\vec{\mu}} = \vec{y}$ is unbiased and attains the Cramér--Rao bound --- and is nonetheless \emph{inadmissible}: the shrinkage estimator
    \begin{align}
        \hat{\vec{\mu}}^{JS} = \left(1 - \frac{(k-2)\sigma^2}{\|\vec{y}\|^2}\right)\vec{y}
    \end{align}
    has strictly lower MSE for every $\vec{\mu}$. Efficiency in the Cramér--Rao sense is therefore a statement about a restricted competition, not about being the best available estimator. When the objective is prediction rather than inference on a parameter, unbiasedness is a constraint one should be willing to abandon.
\end{remark}

\begin{proposition}[Asymptotic Normality of the MLE]
    Under regularity conditions --- consistency, $\theta_0$ in the interior of $\Theta$, twice continuous differentiability of the log-likelihood, and a nonsingular $\mathcal{I}(\theta_0)$ --- as $n \to \infty$,
    \begin{align}
        \sqrt{n}(\hat{\theta}_{MLE} - \theta_0) \xrightarrow{d} \mathcal{N}\!\left(0,\, \mathcal{I}(\theta_0)^{-1}\right).
    \end{align}
    The MLE is therefore consistent and asymptotically efficient: its asymptotic variance attains the Cramér--Rao floor.
\end{proposition}
\begin{proof}[Proof sketch]
    Taylor-expand the score around the true parameter $\theta_0$:
    \begin{align}
        0 = s(\hat{\theta}) = s(\theta_0) + \nabla_\theta s(\theta_0)(\hat{\theta}-\theta_0) + o_p(\|\hat{\theta}-\theta_0\|),
    \end{align}
    where $s(\hat{\theta}) = 0$ by the first-order condition (which requires $\hat{\theta}$ to be interior, hence the need for consistency first). Rearranging,
    \begin{align}
        \sqrt{n}(\hat{\theta}-\theta_0) = \left[-\frac{1}{n}\nabla_\theta s(\theta_0)\right]^{-1} \frac{s(\theta_0)}{\sqrt{n}} + o_p(1).
    \end{align}
    By the law of large numbers, $-\frac{1}{n}\nabla_\theta s(\theta_0) = \frac{1}{n}\sum_i[-\nabla^2_\theta\log p(y_i;\theta_0)] \xrightarrow{p} \mathcal{I}(\theta_0)$, using the information matrix equality. By the central limit theorem, since the individual scores $\nabla_\theta\log p(y_i;\theta_0)$ are i.i.d.\ with mean zero and variance $\mathcal{I}(\theta_0)$,
    \begin{align}
        \frac{s(\theta_0)}{\sqrt{n}} = \frac{1}{\sqrt{n}}\sum_{i=1}^n \nabla_\theta\log p(y_i;\theta_0) \xrightarrow{d} \mathcal{N}(0,\mathcal{I}(\theta_0)).
    \end{align}
    Combining via Slutsky's theorem gives $\sqrt{n}(\hat{\theta}-\theta_0) \xrightarrow{d} \mathcal{N}(0,\mathcal{I}(\theta_0)^{-1})$, since $\mathcal{I}^{-1}\mathcal{I}\,\mathcal{I}^{-1} = \mathcal{I}^{-1}$.
\end{proof}

\begin{remark}[Asymptotic Efficiency Does Not Mean Finite-Sample Efficiency]
    The MLE attains the Cramér--Rao bound \emph{in the limit}. In finite samples, it is typically biased, and there is no contradiction with the bound, which applies only to unbiased estimators. The canonical example is the Gaussian variance: the MLE is
    \begin{align}
        \hat{\sigma}^2_{MLE} = \frac{1}{n}\sum_{i=1}^n (y_i - \bar{y})^2, \qquad \mathbb{E}[\hat{\sigma}^2_{MLE}] = \frac{n-1}{n}\sigma^2 \neq \sigma^2,
    \end{align}
    which is why the unbiased sample variance divides by $n-1$. The same phenomenon appears in linear regression, where the MLE of the error variance is $\mathrm{RSS}/n$ rather than the unbiased $\mathrm{RSS}/(n-k)$. The bias is $O(1/n)$ and vanishes asymptotically, but with $k$ comparable to $n$ --- the regime of most cross-sectional return prediction --- it is not negligible.
\end{remark}

\begin{proposition}[Invariance]
    If $\hat{\theta}$ is the MLE of $\theta$ and $g$ is any function, then $g(\hat{\theta})$ is the MLE of $g(\theta)$. Moreover, if $g$ is continuously differentiable with Jacobian $G = \nabla_\theta g(\theta_0)$, the delta method gives
    \begin{align}
        \sqrt{n}\left(g(\hat{\theta}) - g(\theta_0)\right) \xrightarrow{d} \mathcal{N}\!\left(0,\; G\,\mathcal{I}(\theta_0)^{-1}G^T\right).
    \end{align}
\end{proposition}

\begin{remark}
    Invariance is a genuine convenience and is not shared by all estimators (it fails for unbiased estimation: if $\hat{\theta}$ is unbiased for $\theta$, $g(\hat{\theta})$ is generally biased for $g(\theta)$ by Jensen). It means one may parametrize a model in whatever way makes optimization best-conditioned --- estimating $\log\sigma$ rather than $\sigma$ to enforce positivity, or a Cholesky factor rather than a covariance matrix to enforce positive-definiteness --- and simply transform back at the end.
\end{remark}

Everything above assumes the truth lies in the model family, yet this is usually not the case.

\begin{theorem}[Quasi-MLE and the Sandwich Covariance]
    Suppose the data are generated by a density $g$ that need \emph{not} belong to the family $\{p(\cdot\,;\theta)\}$. Then, under regularity conditions, the MLE remains consistent --- but for the \emph{pseudo-true} parameter
    \begin{align}
        \theta^* = \arg\min_{\theta\in\Theta} D_{\mathrm{KL}}\!\left(g \,\|\, p_\theta\right) = \arg\max_{\theta\in\Theta}\mathbb{E}_g[\log p(y;\theta)],
    \end{align}
    the member of the family closest to the truth in KL divergence. Its asymptotic distribution is
    \begin{align}
        \sqrt{n}(\hat{\theta} - \theta^*) \xrightarrow{d} \mathcal{N}\!\left(0,\; H^{-1} S H^{-1}\right),
    \end{align}
    where
    \begin{align}
        H = -\mathbb{E}_g\!\left[\nabla^2_\theta \log p(y;\theta^*)\right], \qquad
        S = \mathbb{E}_g\!\left[\nabla_\theta \log p(y;\theta^*)\,\nabla_\theta \log p(y;\theta^*)^T\right].
    \end{align}
\end{theorem}

\begin{remark}[Why the Sandwich Matters]
    Under correct specification the information matrix equality gives $H = S = \mathcal{I}(\theta_0)$, and the sandwich collapses to $\mathcal{I}^{-1}S\,\mathcal{I}^{-1} = \mathcal{I}(\theta_0)^{-1}$: the familiar result. Under misspecification the equality fails, $H \neq S$, and the sandwich does not simplify. Reporting $\mathcal{I}^{-1}/n$ standard errors when the model is wrong therefore produces standard errors that are simply the wrong size --- usually too small, hence $t$-statistics that are too large and spuriously significant.

    This is the theoretical basis for \emph{robust} (Huber--White, ``sandwich'') standard errors, and given that no financial model is correctly specified, the sandwich should be regarded as the default rather than an occasional robustness check.
    \begin{itemize}[noitemsep]
        \item Comparing the Hessian estimate $\hat{H}$ against the outer-product-of-gradients estimate $\hat{S}$ is itself a specification test --- this is White's information matrix test. A large discrepancy is evidence the model is wrong.
        \item With dependent data the argument extends, but $S$ must be estimated with a HAC (Newey--West) estimator that accounts for serial correlation in the scores; the i.i.d.\ outer-product form understates it, often badly, for persistent series.
        \item Consistency for $\theta^*$ is often good enough. Estimating a Gaussian model on fat-tailed returns still recovers the correct mean and variance, because the KL projection of any distribution onto the Gaussian family matches its first two moments. What one loses is any claim about tail probabilities, and the validity of the na\"ive standard errors.
    \end{itemize}
\end{remark}

\begin{remark}[Computing the MLE]
    The generic approach is Newton--Raphson,
    \begin{align}
        \theta^{(m+1)} = \theta^{(m)} - \left[\nabla^2_\theta \ell(\theta^{(m)})\right]^{-1} s(\theta^{(m)}),
    \end{align}
    which converges quadratically near the optimum but requires the Hessian and can misbehave far from it. Two standard substitutes for the Hessian exploit the information matrix equality: \emph{Fisher scoring} replaces $-\nabla^2\ell$ by its expectation $n\mathcal{I}(\theta)$, which is guaranteed positive-definite and so always yields an ascent direction; and \emph{BHHH} replaces it by the outer product of gradients $\sum_i \nabla\log p_i \nabla\log p_i^T$, which requires only first derivatives.
\end{remark}


\begin{remark}[The Classical Trinity]
    The asymptotic normality of the MLE generates a framework for hypothesis testing. For a null hypothesis $H_0 : R\theta_0 = r$ (where $R$ is a $q \times k$ restriction matrix), three asymptotically equivalent tests with the same $\chi^2_q$ null distribution are available:

    \begin{itemize}[noitemsep]
        \item \textbf{Wald test:} $(R\hat{\theta} - r)^T[R\,\widehat{\mathcal{I}}^{-1}R^T/n]^{-1}(R\hat{\theta}-r) \xrightarrow{d} \chi^2_q$. Requires only the unrestricted MLE $\hat{\theta}$; this is the test implicitly used when reporting $t$-statistics on regression coefficients (the scalar case $q=1$ gives $t = \hat{\theta}_j/\widehat{\mathrm{se}}(\hat{\theta}_j) \xrightarrow{d} \mathcal{N}(0,1)$).
        \item \textbf{Likelihood ratio test:} $2[\ell(\hat{\theta}) - \ell(\tilde{\theta})] \xrightarrow{d} \chi^2_q$, where $\tilde{\theta}$ is the restricted MLE. Requires fitting both models; invariant to reparametrization.
        \item \textbf{Score (Lagrange multiplier) test:} $s(\tilde{\theta})^T\widehat{\mathcal{I}}^{-1}s(\tilde{\theta})/n \xrightarrow{d} \chi^2_q$. Requires only the restricted MLE $\tilde{\theta}$; useful when the unrestricted model is expensive to fit.
    \end{itemize}

    Wald asks how far $\hat{\theta}$ is from the restricted set, LR asks how much log-likelihood is given up by imposing the restriction, and Score asks how steep the likelihood still is at the restricted estimate. All three are equivalent to first order, which is why they share a limiting distribution.
\end{remark}

\begin{remark}[The Three Tests Do Not Agree in Finite Samples]
    Asymptotic equivalence is a statement about the limit, and the finite-sample differences are systematic rather than random. In the Gaussian linear regression model with \emph{known} $\sigma^2$, the three statistics coincide exactly. With $\sigma^2$ \emph{unknown} --- the realistic case --- they do not. Writing $\tilde{\sigma}^2$ and $\hat{\sigma}^2$ for the restricted and unrestricted variance estimates and $x = \tilde{\sigma}^2/\hat{\sigma}^2 \geq 1$,
    \begin{align}
        W = n(x-1), \qquad \mathrm{LR} = n\log x, \qquad \mathrm{LM} = n(1 - x^{-1}),
    \end{align}
    and since $x - 1 \geq \log x \geq 1 - x^{-1}$ for all $x \geq 1$, these obey the \emph{Berndt--Savin inequality}
    \begin{align}
        W \;\geq\; \mathrm{LR} \;\geq\; \mathrm{LM}.
    \end{align}
\end{remark}

\begin{remark}[The Wald Test Is Not Invariant to Reparametrization]
    A further and more serious defect of the Wald test: its value depends on how the restriction is algebraically written. Testing $H_0 : \theta = 1$ and testing the equivalent $H_0 : \log\theta = 0$, or $H_0 : \theta^2 = 1$, produce \emph{different} Wald statistics from the same data and can produce different conclusions. The likelihood ratio test is invariant to any smooth reparametrization, since it compares maximized likelihood values, which are coordinate-free.
\end{remark}

\section{Regression}

\begin{theorem}[Singular Value Decomposition]
    For a complex $m \times n$ matrix $A$, there exists a decomposition
    \begin{align}
        A = U\Sigma V^\dagger ,
    \end{align}
    where $U$ is an $m \times m$ unitary matrix, $\Sigma$ is a rectangular diagonal $m \times n$ matrix with non-negative real entries $\sigma_1 \geq \sigma_2 \geq \cdots \geq 0$, and $V$ is an $n \times n$ unitary matrix.
\end{theorem}
\begin{proof}
    Consider the $n \times n$ matrix $A^\dagger A$. It is Hermitian and positive semi-definite ($\vec{x}^\dagger A^\dagger A\vec{x} = \|A\vec{x}\|^2 \geq 0$), so by the spectral theorem it is unitarily diagonalizable with non-negative real eigenvalues $\{\lambda_i\}$. Therefore $A^\dagger A = V\Lambda V^\dagger$. Define the singular values $\sigma_i = \sqrt{\lambda_i}$, which are real and non-negative, so that $A^\dagger A\vec{v}_i = \sigma_i^2\vec{v}_i$ for orthonormal eigenvectors $\vec{v}_i$, ordered so that $\sigma_1 \geq \sigma_2 \geq \cdots$.

    For each $i$ with $\sigma_i > 0$, set $\vec{u}_i = \frac{1}{\sigma_i}A\vec{v}_i$. These are orthonormal, since
    \begin{align}
        \vec{u}_i^\dagger \vec{u}_j = \frac{1}{\sigma_i\sigma_j}\vec{v}_i^\dagger A^\dagger A \vec{v}_j = \frac{\sigma_j^2}{\sigma_i\sigma_j}\vec{v}_i^\dagger\vec{v}_j = \delta_{ij}.
    \end{align}
    The eigenvectors with $\sigma_i = 0$ span $\ker(A)$ and contribute nothing; the corresponding $\vec{u}_i$ are obtained by extending $\{\vec{u}_i : \sigma_i > 0\}$ to an orthonormal basis of $\mathbb{C}^m$ via Gram--Schmidt. Hence $A\vec{v}_i = \sigma_i\vec{u}_i$ for all $i$, i.e.\ $AV = U\Sigma$, giving $A = U\Sigma V^\dagger$.
\end{proof}

\begin{figure}[h]
    \centering
    \includegraphics[width=0.4\linewidth]{svd.png}
    \caption{SVD mapping}
    \label{fig:placeholder}
\end{figure}

\begin{remark}
    If $X$ is a matrix of data, the singular value decomposition can be written $X = \sum_{i=1}^{\min(m,n)}\sigma_i\vec{u}_i\vec{v}_i^T$. A rank $k$ approximation (or ``reduced SVD'') of the data can thus be written
    \begin{align}
        X_k = \sum_{i=1}^{k}\sigma_i\vec{u}_i\vec{v}_i^T.
    \end{align}
\end{remark}

\begin{theorem}[Eckart-Young] The singular value decomposition is the best rank-$k$ approximation of a matrix $X$ --- that is, for the Frobenius norm,
\begin{align}
    X_k &= \arg \min_{\text{rank}(Y)\leq k}||X-Y||_F,\\
    ||X-X_k||^2_F &= \sum_{i>k}\sigma^2_i.
\end{align}
    
\end{theorem}
\begin{proof}
    Write the full SVD as $X = U\Sigma V^T = \sum_{i=1}^r \sigma_i \vec{u}_i\vec{v}_i^T$, where $r = \mathrm{rank}(X)$. For any matrix $Y$ with $\mathrm{rank}(Y) \leq k$, the null space of $Y$ has dimension at least $n - k$. Consider the $(k+1)$-dimensional subspace $W = \mathrm{span}(\vec{v}_1,\dots,\vec{v}_{k+1})$. By a dimension argument, $W \cap \ker(Y) \neq \{0\}$, so there exists a unit vector $\vec{w} \in W$ with $Y\vec{w} = 0$. Writing $\vec{w} = \sum_{i=1}^{k+1} c_i \vec{v}_i$ with $\sum c_i^2 = 1$,
    \begin{align}
        \|X - Y\|_F^2 &\geq \|(X-Y)\vec{w}\|_2^2 = \|X\vec{w}\|_2^2 = \left\|\sum_{i=1}^{k+1} c_i \sigma_i \vec{u}_i\right\|_2^2 = \sum_{i=1}^{k+1} c_i^2 \sigma_i^2 \geq \sum_{i=1}^{k+1} c_i^2 \sigma_{k+1}^2 = \sigma_{k+1}^2.
    \end{align}
    More carefully, iterating this argument for each discarded component gives $\|X - Y\|_F^2 \geq \sum_{i>k}\sigma_i^2$. Meanwhile, the truncated SVD $X_k$ achieves equality:
    \begin{align}
        \|X - X_k\|_F^2 = \left\|\sum_{i=k+1}^r \sigma_i \vec{u}_i \vec{v}_i^T\right\|_F^2 = \sum_{i=k+1}^r \sigma_i^2,
    \end{align}
    using the fact that the Frobenius norm of a sum of rank-1 matrices with orthonormal columns and rows equals the sum of the squares of the coefficients.
\end{proof}

\begin{proposition}[Maximum Variance]
    The singular values are ordered with descending variance. If sample covariance of centered data is defined as
    \begin{align}
        \xi = \frac 1n X^TX,
    \end{align}
    then $\max_{||\vec{v}||=1} \text{Var}(X\vec{v}) = \max_{||\vec{v}||=1} \vec{v}^T \xi \vec{v}$ has solution vector equal to the eigenvector with the largest eigenvalue, and has a solution equal to that eigenvalue. Hence, the first SVD vector points in the direction of greatest data variance.
\end{proposition}
\begin{proof}
    The problem $\max_{\|\vec{v}\|=1} \vec{v}^T \xi \vec{v}$ is solved via the method of Lagrange multipliers. Define $\mathcal{L} = \vec{v}^T\xi\vec{v} - \mu(\vec{v}^T\vec{v} - 1)$. Setting $\nabla_{\vec{v}}\mathcal{L} = 0$ gives $2\xi\vec{v} = 2\mu\vec{v}$, i.e., $\vec{v}$ must be an eigenvector of $\xi$ with eigenvalue $\mu$. Substituting back, $\vec{v}^T\xi\vec{v} = \mu\vec{v}^T\vec{v} = \mu$, so the maximum is the largest eigenvalue $\lambda_1$ and the maximizer is the corresponding eigenvector $\vec{v}_1$. For subsequent components, imposing the additional constraint $\vec{v}^T\vec{v}_j = 0$ for $j < i$ at the $i$th step (which adds Lagrange terms that vanish by orthogonality of eigenvectors) gives $\vec{v}_i$ as the $i$th eigenvector with value $\lambda_i$.
\end{proof}

\begin{theorem}[Best Estimator]
    Let $\{(\vec{x}_i,y_i)\}$ be observed data. Then the conditional expectation is the best predictor of $y$ under squared error loss --- that is,
    \begin{align}
    \mathbb{E}[y|x] = \arg \min_{f \in \mathbb{F}}\mathbb{E}[(y-f(x))^2],
    \end{align}
    for the class of measurable functions.
\end{theorem}
\begin{proof}
    Via the tower property,
    \begin{align}
        \mathbb{E}[(y-f(x))^2] = \mathbb{E}[\mathbb{E}[(y-f(x))^2|x]].
    \end{align}
    Massaging the inside yields
    \begin{align}
        (y-f)^2 &= (y -\mathbb{E}[y|x] + \mathbb{E}[y|x] -f)^2\\
        & = (y -\mathbb{E}[y|x])^2 +2(y -\mathbb{E}[y|x])(\mathbb{E}[y|x] -f) + (\mathbb{E}[y|x] -f)^2.
    \end{align}
    Taking the conditional expectation gives
    \begin{align}
        \mathbb{E}[(y-f(x))^2|x]= \mathbb{E}[(y -\mathbb{E}[y|x])^2|x] + (\mathbb{E}[y|x] -f)^2,
    \end{align}
    which is minimized if $f = \mathbb{E}[y|x]$. Therefore, this choice pointwise minimizes the unconditional case, and is the unconditional minimum.
\end{proof}

\begin{figure}[h]
    \centering
    \includegraphics[width=0.6\linewidth]{linreg.jpeg}
    \caption{Gaussian linear regression}
    \label{fig:placeholder}
\end{figure}


\begin{definition}[Gaussian Linear Regression]
Call a \textit{linear regression} a model of the form
\begin{align}
    y = \vec{x}\cdot\vec{\beta} + \varepsilon.
\end{align}
The model is said to be a \textit{Gaussian linear regression} if
\begin{align}
    y | \vec{x} \sim \mathcal{N}(\vec{x}\cdot \vec{\beta}, \sigma^2)
\end{align}
for homoskedastic error.
\end{definition}

\begin{proposition}[Maximum Entropy of the Gaussian]
    Among all continuous distributions on $\mathbb{R}$ with fixed mean $\mu$ and variance $\sigma^2$, the Gaussian $\mathcal{N}(\mu,\sigma^2)$ uniquely maximizes the differential entropy
    \begin{align}
        h(p) = -\int_{-\infty}^\infty p(x)\log p(x)\,dx.
    \end{align}
\end{proposition}
\begin{proof}
    Let $p$ be any density with mean $\mu$ and variance $\sigma^2$, and let $\phi = \mathcal{N}(\mu,\sigma^2)$ be the Gaussian with the same moments. The non-negativity of KL divergence gives
    \begin{align}
        0 \leq D_{\mathrm{KL}}(p\,\|\,\phi) = \int p(x)\log\frac{p(x)}{\phi(x)}\,dx = -h(p) - \int p(x)\log\phi(x)\,dx.
    \end{align}
    Since $\log\phi(x) = -\frac{1}{2}\log(2\pi\sigma^2) - \frac{(x-\mu)^2}{2\sigma^2}$, and both $p$ and $\phi$ share the same mean and variance,
    \begin{align}
        \int p(x)\log\phi(x)\,dx = -\tfrac{1}{2}\log(2\pi\sigma^2) - \tfrac{1}{2\sigma^2}\int p(x)(x-\mu)^2\,dx = -\tfrac{1}{2}\log(2\pi\sigma^2) - \tfrac{1}{2} = \int \phi(x)\log\phi(x)\,dx,
    \end{align}
    which is $-h(\phi)$. Substituting back into the KL inequality gives $h(p) \leq h(\phi)$, with equality if and only if $p = \phi$ almost everywhere.
\end{proof}

\begin{remark}[Choice of Noise Model]
    The maximum entropy result gives the Gaussian noise assumption its privileged status. Suppose the true data-generating process is $y = \vec{x}\cdot\vec{\beta} + \varepsilon$, and all that we know about $\varepsilon$ is its mean (zero) and variance ($\sigma^2$). The maximum entropy principle instructs us to model $\varepsilon$ as Gaussian: it is the distribution that encodes exactly those two constraints and nothing else, injecting no additional structure beyond what the data support. Any other choice of noise model implicitly asserts additional knowledge about the error distribution's shape.

    Under Gaussian noise, the log-likelihood of the sample is
    \begin{align}
        \ell(\vec{\beta}) = -\frac{n}{2}\log(2\pi\sigma^2) - \frac{1}{2\sigma^2}\sum_{i=1}^n(y_i - \vec{x}_i\cdot\vec{\beta})^2,
    \end{align}
    so maximizing the likelihood is exactly equivalent to minimizing the sum of squared residuals. The Gaussian assumption is therefore the theoretical justification for least squares as the default loss: ordinary least squares is MLE under the noise model that makes the fewest assumptions given fixed variance.

    Different assumptions about the error distribution yield different loss functions. The following table summarizes the correspondence, which unifies the regression methods covered throughout these notes:

    \begin{table}[H]
    \centering
    \renewcommand{\arraystretch}{1.3}
    \begin{tabular}{@{}llll@{}}
    \toprule
    \textbf{Noise model} & \textbf{Density} $p(\varepsilon)$ & \textbf{MLE loss} & \textbf{Method} \\
    \midrule
    Gaussian & $\propto \exp\!\left(-\frac{\varepsilon^2}{2\sigma^2}\right)$ & $\sum(y_i - \hat{y}_i)^2$ & OLS, Ridge, LASSO \\
    Laplace & $\propto \exp\!\left(-\frac{|\varepsilon|}{\sigma}\right)$ & $\sum|y_i - \hat{y}_i|$ & Least absolute deviations \\
    Bernoulli & $p^y(1-p)^{1-y}$, $p = \sigma(\vec{x}\cdot\vec{\beta})$ & Cross-entropy & Logistic regression \\
    Poisson & $\frac{\lambda^y e^{-\lambda}}{y!}$, $\lambda = e^{\vec{x}\cdot\vec{\beta}}$ & Poisson deviance & Log-linear GLM \\
    Asymmetric Laplace & $\propto \exp\!\left(-\frac{\varepsilon(\tau - \mathbf{1}_{\varepsilon<0})}{\sigma}\right)$ & Pinball loss & Quantile regression \\
    \bottomrule
    \end{tabular}
    \caption{Noise model, implied loss function, and regression method.}
    \end{table}
\end{remark}

\begin{proposition}[Best Linear Predictor]
    The best linear predictor of $y$ given $\vec{x}$ under squared error loss, i.e.\ the minimizer of $\mathbb{E}[(y - \vec{x}\cdot\vec{\beta})^2]$ over $\vec{\beta}$, is the \emph{linear projection} coefficient
    \begin{align}
        \vec{\beta}^* = \mathbb{E}[\vec{x}\vec{x}^T]^{-1}\,\mathbb{E}[\vec{x}y],
    \end{align}
    whose sample analogue is the ordinary least-squares (OLS) estimator
    \begin{align}
        \hat{\beta}= (X^TX)^{-1}X^T\vec{y}.
    \end{align}
    This requires only that $\mathbb{E}[\vec{x}\vec{x}^T]$ be invertible; \emph{no exogeneity assumption is needed}.
\end{proposition}
\begin{proof}
    Setting the gradient of the population objective to zero,
    \begin{align}
        \nabla_{\vec{\beta}} \mathbb{E}[(y-\vec{x}\cdot \vec{\beta})^2] &= -2\,\mathbb{E}[\vec{x}(y-\vec{x}\cdot \vec{\beta})] = 0 \\
        \Longrightarrow \quad \mathbb{E}[\vec{x}\vec{x}^T]\vec{\beta} &= \mathbb{E}[\vec{x}y],
    \end{align}
    which inverts to the stated $\vec{\beta}^*$. The Hessian is $2\,\mathbb{E}[\vec{x}\vec{x}^T] \succ 0$, so this is the global minimum. Defining the projection error $\varepsilon = y - \vec{x}\cdot\vec{\beta}^*$, the first-order condition reads $\mathbb{E}[\vec{x}\varepsilon] = 0$: the projection error is orthogonal to the regressors \emph{by construction}. Replacing population moments by sample moments gives $\hat{\beta} = \mathbb{E}_n[\vec{x}_i\vec{x}_i^T]^{-1}\mathbb{E}_n[\vec{x}_i y_i] = (X^TX)^{-1}X^T\vec{y}$.
\end{proof}

\begin{remark}[Geometric Interpretation]

To elaborate further on the use of the word ``projection'', the problem $\hat{\beta} = \arg \min_\beta ||\vec{y} - X\vec{\beta}||^2$ is finding the closest vector to $\vec{y}$ that can be constructed out of the columns of $X$. Thus, we are finding a vector in the column space of $X$. For $\hat{y}=X\hat{\beta} \in C(X)$ to be closest to $\vec{y}$, the residual $\vec{y} - \hat{y}$ must be orthogonal to every vector in $C(X)$ (otherwise, moving in an allowed direction could decrease the distance). 

Therefore, the columns of $X$, being a basis for $C(X)$, must each be orthogonal to the residual, so, if $X^TX$ is full-rank,
\begin{align}
    X^T(\vec{y}-X\hat{\beta})=\vec{0} \quad \Rightarrow \quad\hat{\beta} = (X^TX)^{-1}X^T\vec{y}.
\end{align}
Hence, 
\begin{align}
    \hat{y} = X(X^TX)^{-1}X^T\vec{y} = P\vec{y},
\end{align}
where $P$ is an orthogonal projection matrix (idempotent and self-adjoint) onto the column space of $X$.
\end{remark}

\begin{figure}[h]
    \centering
    \includegraphics[width=0.5\linewidth]{fV5Zd.png}
    \caption{Regression as projection}
    \label{fig:placeholder}
\end{figure}

\begin{remark}[Projection versus Causation]
    The orthogonality $\mathbb{E}[\vec{x}\varepsilon]=0$ is a definitional property of the projection and carries no economic content: it holds even if $\vec{x}$ is correlated with omitted variables, measurement error, or simultaneously determined with $y$. What \emph{is} an assumption is the stronger \emph{exogeneity} (or mean-independence) condition
    \begin{align}
        \mathbb{E}[\varepsilon \mid \vec{x}] = 0,
    \end{align}
    which upgrades the linear projection to the conditional expectation $\mathbb{E}[y\mid\vec{x}] = \vec{x}\cdot\vec{\beta}$ and licenses a structural or causal reading of $\vec{\beta}$. This is the distinction between ``my regression predicts $y$'' and ``my regression identifies the effect of $x$ on $y$.'' In return prediction, the former is usually all that is claimed, and all that is needed.
\end{remark}

\begin{theorem}[Gauss--Markov]
Consider the linear regression model
\begin{align}
\vec{y} = X\vec{\beta} + \vec{\varepsilon},
\end{align}
where $\vec{y} \in \mathbb{R}^n$, $X \in \mathbb{R}^{n\times k}$ has full column rank,
and the error vector satisfies
\begin{align}
E[\vec{\varepsilon}]=0, \qquad \operatorname{Var}(\vec{\varepsilon})=\sigma^2 I_n.
\end{align}
Then the OLS estimator
\begin{align}
\hat{\beta}=(X^T X)^{-1}X^T \vec{y}
\end{align}
is the best linear unbiased estimator (BLUE) of $\vec{\beta}$ --- that is, no other linear unbiased estimator has lower variance.
\end{theorem}

\begin{proof}
Consider another estimator $\Tilde{\beta} = C\vec{y} = ((X^T X)^{-1}X^T+D)\vec{y}$, for some matrix $D$. Then
\begin{align}
    \mathbb{E}[\Tilde{\beta}] &= \mathbb{E}[((X^T X)^{-1}X^T+D)(X\vec{\beta} + \vec{\varepsilon})]\\
    &=(X^T X)^{-1}X^TX\vec{\beta} +DX\vec{\beta} \\
    & = (I_k + DX)\vec{\beta},
\end{align}
so $DX\vec{\beta}$ must be 0 to not bias the estimator (note that $D=0$ gives the usual OLS estimator, which is unbiased). Then,
\begin{align}
    \text{Var}(\Tilde{\beta}) &= \text{Var}(C\vec{y})\\
    & = \mathbb{E}[(C(\vec{y}-\mathbb{E}[\vec{y}]))(C(\vec{y}-\mathbb{E}[\vec{y}]))^T]\\
    &=C\text{Var}(\vec{y})C^T\\
    &= \sigma^2 CC^T\\
    &= \sigma^2 ((X^T X)^{-1}X^T+D)(X(X^T X)^{-1}+D^T)\\
    &= \sigma^2(X^TX)^{-1} + \sigma^2 DD^T\\
    & = \text{Var}(\hat{\beta}) + \sigma^2DD^T.
\end{align}
Since $DD^T$ is positive semi-definite, the OLS variance is exceeded by a positive semi-definite matrix.
\end{proof}

\begin{remark}[MLE, OLS, and Gauss--Markov]
    Gauss--Markov establishes that OLS is the Best Linear Unbiased Estimator (BLUE): no other \emph{linear} unbiased estimator has lower variance. The Cramér--Rao bound is a strictly stronger statement: under Gaussian noise, OLS is MLE, and MLE achieves the lowest variance among \emph{all} regular unbiased estimators, not just linear ones. Gauss--Markov requires only $\mathbb{E}[\vec{\varepsilon}] = 0$ and $\mathrm{Var}(\vec{\varepsilon}) = \sigma^2 I_n$; the MLE efficiency result additionally requires the Gaussian distributional assumption. The two results are therefore complementary: Gauss--Markov is distribution-free but restricted to linear estimators; MLE efficiency is fully efficient but requires correct specification of the noise distribution.
\end{remark}

\begin{proposition}[Regression $t$-Statistic]
    In the model $\vec{Y} = X\vec{\beta}+\vec{\varepsilon}$ with $\vec{\varepsilon}\sim\mathcal{N}(\vec{0},\sigma^2I_n)$ and $X$ of full column rank $k$, with $\hat{\sigma}^2 = \mathrm{RSS}/(n-k)$,
    \begin{align}
        \frac{\hat{\beta}_j - \beta_j}{\hat{\sigma}\sqrt{\left[(X^TX)^{-1}\right]_{jj}}} \sim t_{\,n-k}.
    \end{align}
\end{proposition}
\begin{proof}
    Take $\vec{c} = X(X^TX)^{-1}\vec{e}_j$, so that $\vec{c}^T\vec{Y} = \vec{e}_j^T(X^TX)^{-1}X^T\vec{Y} = \hat{\beta}_j$ and
    \begin{align}
        \|\vec{c}\|^2 = \vec{e}_j^T(X^TX)^{-1}X^TX(X^TX)^{-1}\vec{e}_j = \left[(X^TX)^{-1}\right]_{jj}.
    \end{align}
    Take $P = I_n - H$ with $H = X(X^TX)^{-1}X^T$ the hat matrix, an orthogonal projection of rank $n-k$. Since $HX = X$ we get $P\vec{\mu} = (I-H)X\vec{\beta} = \vec{0}$ and $P\vec{c} = (I-H)X(X^TX)^{-1}\vec{e}_j = \vec{0}$, verifying both hypotheses. Finally $\|P\vec{Y}\|^2 = \mathrm{RSS}$, so $\hat{\sigma}^2 = \mathrm{RSS}/(n-k)$ and Theorem~\ref{thm:general-t} applies with $r = n-k$.
\end{proof}

\begin{proposition}[GLS]
    For Gaussian linear regression $\vec{y} = X\vec{\beta} + \vec{\varepsilon}$, where instead $\Var(\vec{\varepsilon})=\Sigma \neq \sigma^2 I_n$ (heteroskedastic positive-definite covariance), then the BLUE is
    \begin{align}
        \hat{\beta} = (X^T\Sigma^{-1}X)^{-1}X^T\Sigma^{-1}\vec{y}.
    \end{align}
\end{proposition}
\begin{proof}
    Since $\Sigma$ is symmetric and positive-definite, its inverse exists, and its eigenvalues are all real and positive - this permits a ``matrix square root''. Consider
    \begin{align}
        \Sigma^{-1/2}\vec{y} = \Sigma^{-1/2}X\vec{\beta} + \Sigma^{-1/2}\vec{\varepsilon} \quad \Rightarrow \quad \tilde{y} = \tilde{X}\vec{\beta} + \tilde{\varepsilon}.
    \end{align}
    Then $\Var(\tilde{\varepsilon}) = \Sigma^{-1/2}\Sigma\Sigma^{-1/2}=I_n$, so the OLS solution on the tilde variables is the BLUE by Gauss-Markov. Expanding this gives the stated form.

    Geometrically, this is still equivalent to finding the closest vector in $C(X)$ to $\vec{y}$, but under the norm $||\vec{v}||^2_{\Sigma^{-1}}=\vec{v}^T\Sigma^{-1}\vec{v}$ (this is a valid inner product). The covariance is treating residuals in different directions as differently informative under noise, concretized in a norm. Hence, requiring orthogonality under the $\Sigma^{-1}$ inner product between the columns of $X$ and $\vec{y} - \hat{y}$ gives
    \begin{align}
        X^T\Sigma^{-1}(\vec{y} - \hat{y})=\vec{0},
    \end{align}
    which can be rearranged to give the above result.
    
\end{proof}

\begin{definition}[Multicollinearity]
    For data matrix $X$, with the $j$th column being the $j$th covariate's value across all data samples,
    \begin{align}
        X = \left[\vec{x}_1 \cdots \vec{x}_p\right],
    \end{align}
    then two regressors being perfectly correlated is referred to as \textit{multicollinearity}. In that case, the columns of $X$ are linearly dependent, so $X^TX$ is non-invertible, and $\hat{\beta}$ does not exist. If it is almost singular, the entries in the inverse blow up.
\end{definition}

\begin{definition}[Logistic Regression]
    Linear regression is ill-suited for binary outcomes: the linear predictor $\vec{x}_i \cdot \vec{\beta}$ is unbounded, whereas a probability must lie in $[0,1]$. \textit{Logistic regression} resolves this. Consider $\{(\vec{x}_i,Y_i)\}$, where the output $Y \in \{0,1\}$, for classification, is distributed $Y_i |\vec{x}_i \sim \text{Bernoulli}(p_i)$. Then
    \begin{align}
        \mathbb{P}(Y_i = y | \vec{x}_i) = p_i^y(1-p_i)^{1-y}.
    \end{align}
    Define $p_i$ with a \textit{sigmoid} function:
    \begin{align}
        p_i = \frac{1}{1+e^{-\vec{x}_i \cdot \vec{\beta}}}.
    \end{align}
    The likelihood of observing the data is then:
    \begin{align}
        \mathcal{L}(\vec{\beta}) &= \prod_i p_i^{Y_i}(1-p_i)^{1-Y_i}\\
        l(\vec{\beta})&= \sum_i \left[Y_i \vec{x}_i \cdot \vec{\beta} - \log (1 + e^{\vec{x}_i \cdot \vec{\beta}})\right]\\
        \nabla_{\vec{\beta}}l(\vec{\beta}) &= \sum_i \vec{x}_i\left(Y_i - \frac{1}{1+e^{-\vec{x}_i \cdot \vec{\beta}}}\right) = X^T(\vec{y}-\sigma(X\beta))= 0.
    \end{align}
\end{definition}

\begin{figure}[h]
    \centering
    \includegraphics[width=0.6\linewidth]{logreg.png}
    \caption{Logistic and linear regression}
    \label{fig:placeholder}
\end{figure}

\begin{remark}
    The decision boundary occurs where $\vec{x}_i \cdot \vec{\beta} = 0$, as that is where $p=1/2$ occurs. The maximum likelihood solution has no closed form, since the gradient equation is non-linear, and requires an iterative method such as gradient descent. Gradient descent updates parameters in the direction of the negative gradient of the loss:
    \begin{align}
        \vec{\beta}^{(t+1)} = \vec{\beta}^{(t)} + \eta \nabla_{\vec{\beta}} l(\vec{\beta}^{(t)}),
    \end{align}
    where $\eta > 0$ is the learning rate. Convergence to the global maximum is guaranteed here since the log-likelihood is concave.
\end{remark}

\begin{definition}[Deviance]
    For a model with log-likelihood as a function of parameters, and an observed dataset, define deviance as a goodness-of-fit measure:
    \begin{align}
        \text{Dev}(\vec{\beta}) = -2(l(\vec{\beta})-l(\text{saturated})),
    \end{align}
    where the saturated model sets one parameter per data point. Minimizing deviance is hence equivalent to maximizing likelihood.
\end{definition}

\begin{definition}[Model $R^2$]
    Define the \textit{residual deviance} $D = \text{Dev}(\vec{\beta})$ of a model as the deviance where the prediction comes from $\vec{x} \cdot \vec{\beta}$, and define the \textit{null deviance} $D_0 = \text{Dev}(\hat{y} = \Bar{y})$ as the deviance where the prediction is replaced by the average of the response data. Call the following value $R^2$, a measure of the improvement in explanation of the data over the naive model:
    \begin{align}
        R^2 = 1-\frac{D}{D_0}.
    \end{align}
    For linear regression, this is
    \begin{align}
        R^2 = 1-\frac{\sum_i(y_i-\vec{x}_i \cdot \hat{\beta})^2}{\sum_i(y_i-\Bar{y})^2} = 1-\frac{\text{RSS}}{\text{TSS}}.
    \end{align}
\end{definition}



\begin{remark}[OOS $R^2$ in Finance]
    In stock return prediction, in-sample $R^2$ values of 5--10\% are common, but the relevant metric is out-of-sample (OOS) $R^2$. An OOS $R^2$ of just 0.5\% per month for individual stock returns is economically significant and can generate substantial portfolio profits. This starkly contrasts other domains where $R^2 < 30\%$ would be considered a poor model. The reason is that even tiny predictability, compounded over many assets and time periods with proper portfolio construction, translates into meaningful Sharpe ratios.
\end{remark}

\begin{figure}[h]
    \centering
    \includegraphics[width=0.6\linewidth]{kfold.png}
    \caption{K-fold cross-validation}
    \label{fig:placeholder}
\end{figure}

\begin{definition}[K-Fold Cross-Validation]
Another method of validating a model consists in splitting the dataset into $k$ folds (subsections) and training on $k-1$ folds while testing on the remaining fold by computing MSE. This ``leave one out'' is repeated for each fold, and the average MSE is given as the validation error.
\end{definition}


\section{Model Selection}

Given a set of outcomes and a large amount of data per outcome, how to select variables to explain the outcomes? What is the best subset of covariates that predicts the outcomes? There are $2^\Omega$ possibilities to check, yet this is often too costly to use.

\begin{proposition}
    Adding regressor covariates to a model will increase or keep constant the $R^2$.
\end{proposition}
% [CORRECTED] The original argued "the model will set beta_2 = 0" if the new
% regressor doesn't help. It won't: OLS generically returns a nonzero estimate
% even for pure noise. The correct argument is a nested-minimization one.
\begin{proof}
    Let $\mathcal{B}_1 = \{(\beta_1, 0)\}$ and $\mathcal{B}_2 = \{(\beta_1,\beta_2)\}$ be the parameter spaces of the small and large models, so $\mathcal{B}_1 \subset \mathcal{B}_2$. Adding a regressor does not change $\vec{y}$, hence does not change TSS. The larger model minimizes the same RSS objective over a \emph{superset} of the parameter space, so
    \begin{align}
        \text{RSS}' = \min_{\beta \in \mathcal{B}_2} \|\vec{y} - X\beta\|^2 \;\leq\; \min_{\beta \in \mathcal{B}_1}\|\vec{y}-X\beta\|^2 = \text{RSS},
    \end{align}
    and therefore $R^2{}' \geq R^2$. Note this is not because OLS ``sets $\hat{\beta}_2 = 0$'' when $x_2$ is useless --- it generically does not, even when $x_2$ is pure noise. It is simply that the minimum over a larger set cannot be larger. Equality holds only in the knife-edge case where $x_2$ is orthogonal to the residual from the small model. This motivates further diagnostics.
\end{proof}

\begin{definition}[Screening]
    Call \textit{screening} the method whereby covariates are selected by inspecting their correlation with the output variable, and selecting those with correlation above a certain threshold. Notably, this process does not address multicollinearity of regressors.
\end{definition}

\begin{definition}[Information Criteria]
    Define the \textit{Akaike information criterion (AIC)} as a measure of model fit that penalizes complexity:
    \begin{align}
        \text{AIC} = 2k - 2\ln (\hat{L}),
    \end{align}
    where $k$ is the number of estimated parameters, and $\hat{L}$ is the maximized likelihood value. Likewise, define the \textit{Bayesian information criterion (BIC)} as:
    \begin{align}
        \text{BIC} = k \ln(n)- 2\ln (\hat{L}),
    \end{align}
    where $n$ is the number of data points.
\end{definition}

\begin{remark}
    BIC penalizes model complexity more heavily than AIC for $n \geq 8$ (since $\ln(n) > 2$), and therefore tends to select simpler models. In finance, BIC is often preferred for selecting the number of factors in a factor model because overfitting factors can lead to spurious in-sample performance that vanishes out of sample.
\end{remark}

\begin{definition}[Forward Stepwise]
    Rather than checking all subsets of regressors, begin with the null model, and iteratively check $(\text{model} + \beta_jx_j)$ over remaining covariates for the minimum AIC, adding the best covariate until the next addition does not offer a lower model AIC. This is far faster than evaluating the entire power set, yet the MLE variance is non-negligible, so the ``best model'' might change with data resampling.
\end{definition}

\begin{theorem}[Karush-Kuhn-Tucker Conditions]
    Consider the optimization problem
    \begin{align}
        \min_x f(x), \quad g_i(x) \leq 0, \quad h_j(x) = 0
    \end{align}
    for arbitrary numbers of constraints. Define an objective function to be minimized
    \begin{align}
        \mathcal{L} = f(x) + \sum_i \lambda_i g_i(x) + \sum_j \nu_jh_j(x).
    \end{align}
    Assume the problem is convex ($f$ and $g_i$ convex, $h_j$ affine) and that a constraint qualification holds --- for instance Slater's condition, that there exists a strictly feasible point with $g_i(x) < 0$ for all $i$. Then at the optimal point $x^*$:
    \begin{enumerate}
        \item (stationarity) $\nabla_x \mathcal{L} = 0$, as the objective cannot be locally improved.
        \item (primal feasibility) $g_i(x^*) \leq 0, h_j(x^*) = 0$, as the constraints are obeyed.
        \item (dual feasibility) $\lambda_i \geq  0$, as making $g>0$ must increase the objective.
        \item (complementary slackness) $\lambda_ig_i(x^*) = 0$, as constraint $g_i$ must be active whenever the unconstrained minimizer lies outside the feasible region, as otherwise you could move further downhill while remaining feasible. If the constraint isn't forcing, then $\lambda=0$, and this is unconstrained.
    \end{enumerate}
\end{theorem}
\begin{proof}[Proof sketch]
    Stationarity: at any interior optimum the gradient of the objective must vanish; at a boundary optimum the gradient of $f$ must be a non-negative combination of the constraint gradients (otherwise one could move along the boundary to decrease $f$). Primal feasibility is the requirement that $x^*$ is in the feasible set. Dual feasibility follows because the Lagrange multiplier for an inequality constraint measures the rate of change of the optimal value with respect to relaxing the constraint: if $\lambda_i < 0$, tightening $g_i$ would decrease $f$, contradicting optimality. Complementary slackness: if $g_i(x^*) < 0$ (the constraint is slack), then perturbing $g_i$ locally does not affect the optimum, so $\lambda_i = 0$; if $\lambda_i > 0$, the constraint must be active. Together, $\lambda_i g_i(x^*) = 0$. For a convex problem, KKT conditions are both necessary and sufficient for global optimality.
\end{proof}

\begin{definition}[LASSO]
    \textit{LASSO regression} is a regularization to linear regression that imposes an $\ell^1$ norm penalty on $\vec{\beta}$:
    \begin{align}
        \min_\beta\left(||\vec{y}-X\vec{\beta}||_2^2\right), \quad ||\vec{\beta}||_1 \leq c
    \end{align}
    as a constrained optimization problem, or 
    \begin{align}
        \min_\beta\left(||\vec{y}-X\vec{\beta}||_2^2 + \lambda||\vec{\beta}||_1\right)
    \end{align}
    as an unconstrained problem.
\end{definition}

\begin{remark}
    LASSO is not differentiable at 0, so there is no closed-form solution. Geometrically, this constrains $\vec{\beta}$ to lie within a hyper-diamond, whose corners lie on the coordinate axes. Because the RSS contours are ellipsoidal, the constrained minimum typically occurs at a corner, setting some coefficients exactly to zero. LASSO therefore performs \emph{variable selection}: it produces sparse solutions, setting irrelevant regressors to exactly 0. So-called ``LASSO paths'' are used for variable selection by setting $\lambda$ very large, then gradually decreasing it to watch which regressors ``turn on'' as LASSO deems them important.
\end{remark}

\begin{theorem}[LASSO Soft Thresholding]
    For orthonormal design $X$ ($X^TX = I$), let $\hat{\beta}^{\mathrm{OLS}}_j = \vec{x}_j^T\vec{y}$ be the OLS estimator for the $j$th coefficient. Then the LASSO solution is given componentwise by the \emph{soft thresholding operator}:
    \begin{align}
        \hat{\beta}^{\mathrm{LASSO}}_j = \mathrm{sign}(\hat{\beta}^{\mathrm{OLS}}_j)\max\!\left(|\hat{\beta}^{\mathrm{OLS}}_j| - \frac{\lambda}{2},\; 0\right).
    \end{align}
    Coefficients with $|\hat{\beta}^{\mathrm{OLS}}_j| \leq \lambda/2$ are set exactly to zero; the remaining coefficients are shrunk toward zero by $\lambda/2$.
\end{theorem}
\begin{proof}
    With orthonormal design matrices, the LASSO objective separates into independent scalar problems:
    \begin{align}
        \min_{\beta_j}\left[(\hat{\beta}^{\mathrm{OLS}}_j - \beta_j)^2 + \lambda|\beta_j|\right].
    \end{align}
    For $\beta_j > 0$, the derivative is $-2(\hat{\beta}^{\mathrm{OLS}}_j - \beta_j) + \lambda = 0$, giving $\beta_j = \hat{\beta}^{\mathrm{OLS}}_j - \lambda/2$, which is valid only if $\hat{\beta}^{\mathrm{OLS}}_j > \lambda/2$. Similarly for $\beta_j < 0$. If $|\hat{\beta}^{\mathrm{OLS}}_j| \leq \lambda/2$, the minimum occurs at $\beta_j = 0$ (verified by checking the subdifferential condition $0 \in \partial_{\beta_j}\mathcal{L}$, which requires $|2\hat{\beta}^{\mathrm{OLS}}_j| \leq \lambda$). Combining the cases gives the soft thresholding operator.
\end{proof}

\begin{figure}[h]
    \centering
    \includegraphics[width=0.6\linewidth]{lasso_ridge.png}
    \caption{LASSO and Ridge error surfaces}
    \label{fig:placeholder}
\end{figure}

\begin{definition}[Ridge]
    \textit{Ridge regression} is a regularization to linear regression that imposes an $\ell^2$ norm penalty on $\vec{\beta}$:
    \begin{align}
        \min_\beta\left(||\vec{y}-X\vec{\beta}||_2^2\right), \quad ||\vec{\beta}||_2 \leq c
    \end{align}
    as a constrained optimization problem, or 
    \begin{align}
        \min_\beta\left(||\vec{y}-X\vec{\beta}||_2^2 + \lambda||\vec{\beta}||_2^2\right)
    \end{align}
    as an unconstrained problem.
\end{definition}

\begin{remark}
    Geometrically, this constrains $\vec{\beta}$ to a hyper-sphere. Unlike LASSO, the sphere has no corners, so the constrained minimum generically occurs at a smooth point on the boundary where all coordinates are nonzero. Ridge therefore \emph{shrinks} all coefficients toward zero but never sets any exactly to zero: it does not perform variable selection. The shrinkage is not uniform across directions --- as Theorem~\ref{thm:ridge-svd} below shows, it is strongest along the low-variance principal directions and weakest along the high-variance ones.
\end{remark}

\begin{theorem}[Ridge $\vec{\beta}$]
    The $\vec{\beta}$ for Ridge is
    \begin{align}
        \hat{\beta} = (X^TX + \lambda I)^{-1}X^T \vec{y}.
    \end{align}
    Note that even when $X^TX$ is singular, the $\lambda I$ term allows for matrix inversion, making Ridge a natural remedy for multicollinearity.
\end{theorem}
\begin{proof}
    Take the derivative of the objective function to minimize:
    \begin{align}
        \mathcal{L}(\vec{\beta}) &= (\vec{y}-X\vec{\beta})^T(\vec{y}-X\vec{\beta}) + \lambda \vec{\beta}^T\vec{\beta} \\
        \nabla_{\vec{\beta}}\mathcal{L}(\vec{\beta}) &= -2X^T\vec{y} + 2 X^TX\vec{\beta} + 2 \lambda I \vec{\beta} = 0.
    \end{align}
    Rearranging gives the theorem.
\end{proof}

\begin{theorem}[Ridge Shrinkage in the SVD Basis]\label{thm:ridge-svd}
    Let $X = U\Sigma V^T$ be the SVD of the data matrix, with singular values $\sigma_1 \geq \cdots \geq \sigma_p$. Then the Ridge prediction decomposes in the principal component basis as
    \begin{align}
        X\hat{\beta}^{\mathrm{Ridge}} = \sum_{j=1}^p \frac{\sigma_j^2}{\sigma_j^2 + \lambda}\,\vec{u}_j(\vec{u}_j^T\vec{y}).
    \end{align}
    The shrinkage factor $d_j = \sigma_j^2/(\sigma_j^2 + \lambda)$ satisfies $0 < d_j < 1$. Directions of large variance ($\sigma_j^2 \gg \lambda$) are barely shrunk ($d_j \approx 1$), while directions of small variance ($\sigma_j^2 \ll \lambda$) are heavily shrunk ($d_j \approx 0$). Ridge therefore performs \emph{soft dimensionality reduction}: it nearly eliminates components where the signal-to-noise ratio is low.
\end{theorem}
\begin{proof}
    Using $X = U\Sigma V^T$, we have $X^TX = V\Sigma^2 V^T$ and $X^T\vec{y} = V\Sigma U^T\vec{y}$. Then
    \begin{align}
        \hat{\beta}^{\mathrm{Ridge}} &= (V\Sigma^2 V^T + \lambda I)^{-1}V\Sigma U^T\vec{y} = V(\Sigma^2 + \lambda I)^{-1}\Sigma U^T\vec{y}.
    \end{align}
    Multiplying by $X = U\Sigma V^T$:
    \begin{align}
        X\hat{\beta}^{\mathrm{Ridge}} = U\Sigma (\Sigma^2 + \lambda I)^{-1}\Sigma U^T\vec{y} = \sum_{j=1}^p \frac{\sigma_j^2}{\sigma_j^2 + \lambda}\vec{u}_j(\vec{u}_j^T\vec{y}).
    \end{align}
    Each $d_j = \sigma_j^2/(\sigma_j^2+\lambda) \in (0,1)$ since both $\sigma_j^2$ and $\lambda$ are positive. The monotonicity $d_j > d_{j+1}$ when $\sigma_j > \sigma_{j+1}$ follows directly.
\end{proof}

\begin{corollary}[Effective Degrees of Freedom]
    One metric for effective degrees of freedom in regressions is the trace of the hat/projection matrix. The effective degrees of freedom of Ridge regression are
    \begin{align}
        \mathrm{df}(\lambda) = \mathrm{tr}\left[X(X^TX + \lambda I)^{-1}X^T\right] = \sum_{j=1}^p \frac{\sigma_j^2}{\sigma_j^2 + \lambda}.
    \end{align}
    When $\lambda = 0$, $\mathrm{df} = p$ (OLS). As $\lambda \to \infty$, $\mathrm{df} \to 0$ (null model). This provides a continuous measure of model complexity governed by $\lambda$.
\end{corollary}

\begin{definition}[Elastic Net]
    The \textit{elastic net} combines the $\ell^1$ and $\ell^2$ penalties of LASSO and Ridge:
    \begin{align}
        \min_\beta\left(||\vec{y}-X\vec{\beta}||_2^2 + \lambda_1||\vec{\beta}||_1 + \lambda_2||\vec{\beta}||_2^2\right).
    \end{align}
    It inherits variable selection from LASSO while benefiting from the stability of Ridge in the presence of highly correlated predictors. The tuning parameters $\lambda_1, \lambda_2 \geq 0$ are selected by cross-validation; setting $\lambda_2 = 0$ recovers LASSO and $\lambda_1 = 0$ recovers Ridge.
\end{definition}

\begin{remark}
    In all regularized regression methods, the penalty parameters $\lambda$ (or $\lambda_1, \lambda_2$ for the elastic net) govern the bias-variance tradeoff and are chosen by $k$-fold cross-validation: the value minimizing the average out-of-sample MSE across folds is selected.
\end{remark}

\begin{remark}[Regularization in Portfolio Construction]
    In mean-variance portfolio optimization, the optimal portfolio weights $\vec{w}^* = \Sigma^{-1}\vec{\mu}$ (up to scaling) are extremely sensitive to estimation error in $\Sigma$ and $\vec{\mu}$: the optimizer loads aggressively on the smallest eigenvalues of $\hat{\Sigma}$, which are precisely the worst-estimated directions. Shrinking the sample covariance toward a structured target is the standard remedy. The Ledoit--Wolf estimator
    \begin{align}
        \hat{\Sigma}^{\mathrm{LW}} = (1-\alpha)\hat{\Sigma} + \alpha \bar{\lambda}I, \qquad \bar{\lambda} = \tfrac{1}{n}\mathrm{tr}(\hat{\Sigma}),
    \end{align}
    shrinks $\hat{\Sigma}$ toward a scaled identity, pulling the eigenvalues toward their grand mean while leaving the eigenvectors untouched; the optimal $\alpha$ is available in closed form. This is the direct covariance analogue of Ridge. LASSO-penalized portfolio weights instead enforce sparsity, producing portfolios that hold only a subset of available assets, reducing transaction costs.

    The modern refinement is to shrink each eigenvalue \emph{individually} rather than toward a common target. Random matrix theory (the Marchenko--Pastur law and the Silverstein equation) characterizes how the sample eigenvalues of $\hat{\Sigma}$ are dispersed relative to the population eigenvalues when $n/T \to q > 0$, and the Ledoit--P\'ech\'e rotationally invariant estimator applies the optimal nonlinear shrinkage function to each eigenvalue while preserving the sample eigenvectors. This dominates linear shrinkage when $n$ and $T$ are of comparable order, which is the standard regime in equity covariance estimation.
\end{remark}

\section{Factor Models}

To estimate $Y_t = \beta_0 + \beta_1 X_{1t}+\cdots+ \beta_nX_{nt}+U_t$ for $t= 1,\ldots,T$, or $Y = X\beta+U$ as matrices, if there are more features than data points ($n>T$), then the OLS solution fails, since the covariance matrix is not full-rank.

\begin{definition}[Factor Model]
    To address regression in high dimensions, reduce the dimension of the data set:
    \begin{align}
        \vec{X}_t = \vec{\mu}_t + \Lambda \vec{F}_t + \vec{V}_t,
    \end{align}
    where $\vec{X}_t, \vec{\mu}_t$ are $n$-dimensional vectors, $\vec{F}_t$ is a $k \ll n$-dimensional vector of factors, $\Lambda$ is an unobserved $n \times k$ matrix of factor loadings, and $\vec{V}_t$ is an $n$-dimensional zero-mean vector of idiosyncratic shocks.
\end{definition}

% [CORRECTED] The original declared X = [X_1 ... X_T] (observations as
% COLUMNS, so X is n x T) but then used Sigma = X^T X / T as n x n and
% Z = X Gamma_k with Gamma_k n x k -- neither conforms. Fixed to the
% rows-as-observations convention (X is T x n), which is what is actually used.
% NOTE: Theorem 4.2 below reverts to X being n x T; see the remark there.
\begin{definition}[Principal Component Analysis]
    Let $X \in \mathbb{R}^{T \times n}$ be a data matrix whose $t$th \emph{row} is the observation $\vec{X}_t^T \in \mathbb{R}^n$, so there are $T$ observations of $n$ features. In practice, $X$ should first be centered (subtract the column means) so that the PCA directions reflect genuine covariance structure rather than shifts in the mean. Call the covariance of centered data the $n\times n$ matrix $\hat{\Sigma} = \frac{1}{T} X^TX$. To find the $k$ directions that explain the most variance of the data, find the vector maximizing $\text{Var}(X\vec{\gamma})$:
    \begin{align}
        \max_{\vec{\gamma}\in \mathbb{R}^n} \vec{\gamma} \cdot \hat{\Sigma}\vec{\gamma}, \quad ||\vec{\gamma}||_2=1\\
        \Rightarrow \nabla_{\vec{\gamma}}[\vec{\gamma} \cdot \hat{\Sigma}\vec{\gamma}-\lambda\vec{\gamma} \cdot \vec{\gamma}]=0\\
        \hat{\Sigma}\vec{\gamma} = \lambda \vec{\gamma}.
    \end{align}
    So the eigenvector of the covariance matrix maximizes this, and its variance is the eigenvalue. Hence the first principal component is chosen as the covariance eigenvector with largest eigenvalue $\vec{\gamma}_1$. The other principal components are chosen as variance-maximizing directions orthogonal to the first, so the other highest eigenvectors are selected. Note that this is what SVD has been doing.

    Taking $k$ principal components, call $\Gamma_k = \left[\vec{\gamma}_1\cdots\vec{\gamma}_k\right]$ an $n \times k$ matrix. Then the transformed dataset, as a $T \times k$ matrix, is
    \begin{align}
        Z_k = X\Gamma_k = \left[\vec{z}_1\cdots\vec{z}_k\right].
    \end{align}
    The scores $\vec{z}_i = X\vec{\gamma}_i$ are mutually orthogonal, with
    \begin{align}
        \vec{z}_i \cdot \vec{z}_j = \vec{\gamma}_i^T (X^TX) \vec{\gamma}_j = T\,\vec{\gamma}_i^T \hat{\Sigma}\vec{\gamma}_j = T\lambda_j\,\delta_{ij},
    \end{align}
    so that $\text{Var}(\vec{z}_i) = \lambda_i$. (The original claim $\vec{z}_i\cdot\vec{z}_j = \delta_{ij}$ would force every component to have unit variance, contradicting $\text{Var}(\vec{z}_i)=\lambda_i$: it is the \emph{loadings} $\vec{\gamma}_i$ that are orthonormal, not the scores.)
\end{definition}

\begin{figure}[h]
    \centering
    \includegraphics[width=0.5\linewidth]{pca.png}
    \caption{Principal component axes}
    \label{fig:placeholder}
\end{figure}

\begin{remark}
    To check what an appropriate value for $k$ is, run a full $k=n$ PCA, and plot $\alpha_j = \frac{\lambda_j}{\sum_i \lambda_i}$ to check the relative proportion of variance that each PC explains. Alternatively, stop at the component after which there is the greatest drop. Alternatively, if PCA is being used in a model, an information criterion can be used.
\end{remark}

\begin{definition}[Identifiability]
A parameter $\theta$ is \textit{point identified} by a model $P$ if
\begin{align}
    \theta_1 \neq \theta_2 \Rightarrow P_{\theta_1} \neq P_{\theta_2}.
\end{align}
If a model has \textit{latent factors}, which are not immediately physical or observable, there can be issues with identifiability unless a constraint is chosen. This is the equivalent of having gauge freedom.
\end{definition}

% [CORRECTED] The three stated conditions do NOT identify the parameters --
% they only reduce the gauge group from GL(k) to O(k), since any ORTHOGONAL H
% preserves all three. A further normalization is required, plus a condition on
% the idiosyncratic covariance, or Lambda*Lambda^T and Omega are not separable.
\begin{theorem}[Factor Identification]
    Consider the factor model $\vec{X}_t = \vec{\mu}_t + \Lambda \vec{F}_t + \vec{V}_t$. Impose
    \begin{align}
        \mathbb{E}[\vec{F}_t] &= 0, \quad \mathbb{E}[\vec{F}_t \vec{F}_t^T]= I_k, \quad \mathbb{E}[\vec{F}_t \vec{V}_t^T]=0,
    \end{align}
    which gives the covariance decomposition
    \begin{align}
        \Sigma_{\vec{X}_t} &= \Lambda \Lambda^T + \Omega, \qquad \Omega = \mathbb{E}[\vec{V}_t \vec{V}_t^T].
    \end{align}
    These conditions pin down $\vec{\mu}_t$ and the \emph{column space} of $\Lambda$, but they are \emph{not} sufficient to identify $\Lambda$ and $\vec{F}_t$ individually. Two further restrictions are required:
    \begin{enumerate}[noitemsep]
        \item \textbf{Rotational normalization.} $\Lambda^T\Lambda$ is diagonal with strictly decreasing diagonal entries. (Equivalently: the columns of $\Lambda$ are the ordered eigenvectors of $\Lambda\Lambda^T$.) This fixes the residual orthogonal freedom, up to a sign flip of each column.
        \item \textbf{Separability of the idiosyncratic component.} $\Omega$ is diagonal (the \emph{strict} factor model), or, more realistically, $\Omega$ is sufficiently sparse and its eigenvalues remain bounded as $n \to \infty$ while the $k$ largest eigenvalues of $\Lambda\Lambda^T$ diverge (the \emph{approximate} factor model of Chamberlain and Rothschild, 1983). Without such a condition the decomposition $\Sigma = \Lambda\Lambda^T + \Omega$ is not unique and even $k$ is not identified.
    \end{enumerate}
\end{theorem}
\begin{proof}
    Setting $\mathbb{E}[\vec{F}_t]=0$ gives $\mathbb{E}[\vec{X}_t] = \vec{\mu}_t$, so the mean is identified; this is WLOG since we may always recenter $\Tilde{F} = F-\mathbb{E}[F]$ and absorb the shift into $\vec{\mu}_t$. For the remaining parameters, observe that for any invertible $H \in \mathbb{R}^{k\times k}$,
    \begin{align}
        \vec{X}_t = \vec{\mu}_t + \Lambda \vec{F}_t + \vec{V}_t = \vec{\mu}_t + (\Lambda H)(H^{-1} \vec{F}_t) + \vec{V}_t,
    \end{align}
    so $(\Lambda, \vec{F}_t)$ and $(\Lambda H, H^{-1}\vec{F}_t)$ are observationally equivalent: the model has a $GL(k,\mathbb{R})$ gauge freedom. Computing the covariance,
    \begin{align}
        \Sigma = \Lambda\mathbb{E}[\vec{F}_t \vec{F}_t^T]\Lambda^T + \Lambda \mathbb{E}[\vec{F}_t \vec{V}_t^T] + \mathbb{E}[\vec{V}_t\vec{F}_t ^T]\Lambda ^T + \mathbb{E}[\vec{V}_t\vec{V}_t ^T],
    \end{align}
    and imposing the three stated conditions kills the cross terms and yields $\Sigma = \Lambda\Lambda^T + \Omega$.

    However, the normalization $\mathbb{E}[\vec{F}_t\vec{F}_t^T] = I_k$ is preserved by exactly those $H$ satisfying $H^{-1}(H^{-1})^T = I_k$, i.e.\ by the \emph{orthogonal} matrices $H \in O(k)$; and for such $H$, $(\Lambda H)(\Lambda H)^T = \Lambda H H^T \Lambda^T = \Lambda\Lambda^T$, so the covariance decomposition is unchanged. The gauge group has been reduced from $GL(k)$ to $O(k)$ but not eliminated. Requiring $\Lambda^T\Lambda$ diagonal with distinct, ordered entries fixes the remaining $O(k)$ freedom up to column sign, which is why the loadings recovered by PCA are unique only up to a sign convention.
\end{proof}

\begin{remark}[Gauge-Fixing Analogy]
    The sequence is exactly a gauge-fixing procedure. The raw model has redundancy group $GL(k)$. Normalizing the factor covariance to the identity is a partial gauge fixing that leaves a residual $O(k)$ stabilizer --- the analogue of choosing Lorenz gauge and still retaining the freedom to add harmonic functions. Diagonalizing $\Lambda^T\Lambda$ with an ordering convention is the complete gauge fixing, leaving only the discrete $\{\pm 1\}^k$ sign group. This is why any economic interpretation attached to an individual PCA factor (``this is the value factor'') is a statement about a chosen gauge, not about the model.
\end{remark}

\begin{theorem}[Factor Estimation]
    The optimal choice of factors and loadings minimizes the total unexplained variance in the data. The solution is given by the principal components:
    \begin{align}
        \hat{F}_k = \sqrt{T}\Gamma_k(\Sigma), \quad \hat{\Lambda}_k = \frac{X\hat{F}_k }{T},
    \end{align}
    where $\Sigma = X^TX/n$ and $\Gamma_k$ gives the top $k$ eigenvectors.
\end{theorem}
\begin{proof}
    \emph{Convention note:} this proof takes $X$ to be $n \times T$ (features in rows, observations in columns), which is the transpose of the convention fixed in Definition~4.2 above. This is what makes $\|X(I_T - FF^T/T)\|_F$ conformable. Be careful transcribing formulae between the two.

    The optimal parameters minimize
    \begin{align}
       q(\Lambda, F) = \sum_{i=1}^n \sum_{t=1}^T (X_{it} - \vec{\lambda}_i \cdot \vec{F}_t)^2 = ||X - \Lambda F^T||_F^2.
    \end{align}
    If the factors are fixed, the solution is OLS $\hat{\lambda}_i =(F^TF)^{-1}F^T\vec{X}_i$. Imposing the constraint $\frac{F^TF}{T}=I_k$ gives
    \begin{align}
       q(\Lambda, F) =  \sum_{i=1}^n \sum_{t=1}^T \left(X_{it} - \frac{\vec{X}_i^TF}{T}\vec{F}_t\right)^2 = \left\|X \left(I_T- \frac{FF^T}{T}\right)\right\|_F^2.
    \end{align}
    This is minimizing the orthogonal complement from the projection onto factor space, the part of the data that the model cannot explain. Using the Frobenius norm,
    \begin{align}
        \left\|X \left(I_T- \frac{FF^T}{T}\right)\right\|_F^2 &= \text{tr}\left[\left(I_T- \frac{FF^T}{T}\right)^T X^TX \left(I_T- \frac{FF^T}{T}\right)\right]\\
        &= \text{tr}\left[\left(I_T- \frac{FF^T}{T}\right) X^TX\right],
    \end{align}
    exploiting the cyclicity of the trace and idempotency of the projection. So
    \begin{align}
        \left\|X \left(I_T- \frac{FF^T}{T}\right)\right\|_F^2 = C - \frac{1}{T}\text{tr}[F^T X^TXF].
    \end{align}
    Thus, the optimization problem is
    \begin{align}
        \max_F \text{tr}[F^T X^TXF], \quad\frac{F^TF}{T}=I_k.
    \end{align}
    Rescaling to $U = \frac{F}{\sqrt{T}}, S = X^TX$, this is
    \begin{align}
        \max_U \text{tr}[U^TSU], \quad U^TU=I_k.
    \end{align}
    This is maximizing $\sum_{j=1}^k \vec{u}_j \cdot S\vec{u}_j$ subject to $\vec{u}_i \cdot \vec{u}_j = \delta_{ij}$, which is equivalent to finding the first $k$ principal components. This result gives the theorem, keeping track of coefficients.
\end{proof}

\begin{remark}
    When using this to predict with new data, project the new vector to the factor space $\Tilde{f}_{t+1} = (\Lambda^T\Lambda)^{-1}\Lambda^T\vec{x}_{t+1}$, then run the new vector through the model.
\end{remark}

\section{Ensemble Methods}

\begin{definition}[Greedy Algorithm]
    A \textit{greedy algorithm} is any algorithm that follows the problem-solving heuristic of making the locally optimal choice at each stage. A greedy heuristic can yield locally optimal solutions that approximate a globally optimal solution in a reasonable amount of time. A greedy algorithm tends to work well for the class of problems satisfying the following two properties:
    \begin{enumerate}
        \item (Greedy choice property) Whichever choice seems best at a given moment can be made and then (recursively) solve the remaining sub-problems. The choice made by a greedy algorithm may depend on choices made so far, but not on future choices or all the solutions to the subproblem. It iteratively makes one greedy choice after another, reducing each given problem into a smaller one. In other words, a greedy algorithm never reconsiders its choices. This is the main difference from dynamic programming.
        \item (Optimal substructure) An optimal solution to the problem is constructible from optimal solutions to the sub-problems.
    \end{enumerate}
\end{definition}

\begin{definition}[Additive Model]
    An \textit{additive model} takes the form
    \begin{align}
        Y_t = f_M(X_t)=\sum_{m=1}^M \beta_m f_m(X_t) + U_t,
    \end{align}
    where $U_t$ is random noise and $f_m$ is a basis function or weak learner (it does not do considerably better than random guessing).
\end{definition}

\begin{definition}[Function Estimation Problem]
    For variable $\vec{Y}$ the observations and $X$ the collection of feature observations, define the optimal function for some loss function as
    \begin{align}
        f^* = \arg \min_f \mathbb{E}[L\left(\vec{Y}-f(X)\right)].
    \end{align}
    In practice, for finite numbers of observations, not distributions, this is finding the function that minimizes the empirical risk function:
    \begin{align}
        f^* = \arg \min_f \frac{1}{T} \sum_{t=1}^T L\left(Y_t - f(\vec{X}_t)\right).
    \end{align}
\end{definition}

\begin{definition}[Gradient Boosting]
    Given the above problem, restrict the class of functions to additive models. Construct the additive model sequentially:
    \begin{align}
        f_m(X) = f_{m-1}(X) + \beta_m h_m(X).
    \end{align}
    At each step, greedily choose the new function to minimize the loss of the model on the data --- that is, fit the new function to match the existing residuals:
    \begin{align}
        h_m = \arg\min_h \sum_{t=1}^T L\left(\left[Y_t - f_{m-1}(X_t)\right] - h(X_t)\right).
    \end{align}
    More generally, define the ``pseudo-residuals'' as 
    \begin{align}
        g^{(m)}_t = -\left[\frac{\partial L(Y_t,f(X_t) )}{\partial f(X_t)}\right]_{f=f_{m-1}},
    \end{align}
    since this resembles the residuals for MSE loss, and fit $h$ to match this. To calculate the coefficient, find
    \begin{align}
        \beta_m = \arg \min_\beta \sum_{t=1}^T L\left(Y_t,  f_{m-1}(X_t) +\beta h_m(X_t)\right).
    \end{align}
    Then iterating with some shrinkage/learning rate gives
    \begin{align}
        f_m(X) = f_{m-1}(X) + \nu \beta_m h_m(X).
    \end{align}
    To speed up the process and include randomness, one can use stochastic gradient boosting, whereby a randomly selected subset of the data is chosen to optimize over at each step.
\end{definition}

\begin{figure}[h]
    \centering
    \includegraphics[width=0.6\linewidth]{bagging_boosting.png}
    \caption{Bagging and boosting}
    \label{fig:placeholder}
\end{figure}

\begin{definition}[Bagging/Bootstrap Aggregating]
    For data $\mathcal{D}$ and desired predictor $\hat{f}(x;\mathcal{D})$, bootstrap $B$ samples from $\mathcal{D}$ (sample with replacement to a certain size), then train separate models $\hat{f}^{(b)}(x;\mathcal{D}^{(b)})$ on each bootstrapped dataset, and return an average:
    \begin{align}
        \hat{f}^{(\text{bag})}(x; \mathcal{D}) = \frac{1}{B}\sum_{b=1}^B \hat{f}^{(b)}(x;\mathcal{D}^{(b)}).
    \end{align}
    This is essentially Monte Carlo integrating the expectation of the estimator over the data distribution by bootstrap.
\end{definition}

\begin{proposition}[Variance Reduction by Bagging]\label{prop:bagging}
    Suppose the $B$ bootstrapped estimators $\hat{f}^{(b)}$ each have variance $\sigma^2$ and pairwise correlation $\rho \in [0,1]$. Then the variance of the bagged estimator is
    \begin{align}
        \text{Var}\!\left(\hat{f}^{(\text{bag})}\right) = \rho \sigma^2 + \frac{1-\rho}{B}\sigma^2.
    \end{align}
    As $B \to \infty$, the second term vanishes and the variance floor is $\rho\sigma^2$. Bagging therefore reduces variance most when the individual trees are uncorrelated.
\end{proposition}
\begin{proof}
    Write $\hat{f}^{(\text{bag})} = \frac{1}{B}\sum_{b=1}^B \hat{f}^{(b)}$. Then
    \begin{align}
        \text{Var}\!\left(\hat{f}^{(\text{bag})}\right) &= \frac{1}{B^2}\text{Var}\!\left(\sum_{b=1}^B \hat{f}^{(b)}\right) = \frac{1}{B^2}\left[\sum_{b=1}^B \text{Var}(\hat{f}^{(b)}) + \sum_{b \neq b'}\text{Cov}(\hat{f}^{(b)},\hat{f}^{(b')})\right].
    \end{align}
    Since each term has variance $\sigma^2$ and each distinct pair has covariance $\rho\sigma^2$,
    \begin{align}
        \text{Var}\!\left(\hat{f}^{(\text{bag})}\right) &= \frac{1}{B^2}\left[B\sigma^2 + B(B-1)\rho\sigma^2\right] = \frac{\sigma^2}{B} + \frac{B-1}{B}\rho\sigma^2 = \rho\sigma^2 + \frac{1-\rho}{B}\sigma^2.
    \end{align}
\end{proof}

\section{Tree-Based Methods}

Regression trees are a flexible, non-parametric predictive model that recursively partition the input space via a greedy algorithm. Their main advantages are interpretability, the ability to handle categorical variables naturally, and robustness to monotone transformations of features. Their main disadvantage is instability: small perturbations to the data can lead to drastically different splits. Boosting and bagging resolve instability at the cost of interpretability.

\begin{definition}[Regression Tree]
    Let $\vec{X}_t = (X_{t,1},\dots,X_{t,p})^T$ and $Y_t$ be such that
    \begin{align}
        Y_t = \mathbb{E}[Y_t|\vec{X}_t] + U_t = f(\vec{X}_t) + U_t.
    \end{align}
    A \textit{regression tree} approximates $f$ by
    \begin{align}
        H(\vec{X}_t) = \sum_{k=1}^K \beta_k I_k(\vec{X}_t), \quad I_k(\vec{X}_t) = \begin{cases} 1 & \text{if } \vec{X}_t \in R_k \\ 0 & \text{otherwise,} \end{cases}
    \end{align}
    where $R_1,\dots,R_K$ are disjoint regions (leaves) partitioning the input space. The approximation is equivalent to a linear regression on $K$ indicator (dummy) variables. The prediction in each region is the sample average of $Y_t$ over all observations falling in that region.
\end{definition}

\begin{remark}
    More formally, a regression tree can be written as $Y_t = \sum_{i \in \mathcal{T}} \beta_i B_i(\vec{X}_t) + u_t$, where $\mathcal{T}$ is the set of terminal nodes (leaves) and $B_i(\vec{X}_t) = \mathbf{1}\{\vec{X}_t \in R_i\}$. Each region $R_i$ is defined by a sequence of binary splits $X_{s_j} \leq c_j$ or $X_{s_j} > c_j$ along the path from the root to leaf $i$, so $B_i(\vec{X}_t) = \prod_{j \in J_i} \mathbf{1}\{\vec{X}_t \text{ follows split } j\}$, and $\sum_{i \in \mathcal{T}} B_i(\vec{X}_t) = 1$.
\end{remark}

\begin{definition}[Recursive Partitioning]
    Rather than solving the combinatorially intractable problem of finding the globally optimal partition, a regression tree is grown via a greedy algorithm with sequential binary splitting. At each step, given the current observations in a region $R$, find the split variable $j$ and threshold $c$ minimizing total within-child squared error by testing over all factors and thresholds (or a grid, in practice):
    \begin{align}
        \min_{j,\, c} \left[ \sum_{t:\, X_{t,j} < c} (Y_t - \hat{\beta}_{R_1})^2 + \sum_{t:\, X_{t,j} \geq c} (Y_t - \hat{\beta}_{R_2})^2 \right],
    \end{align}
    where $\hat{\beta}_{R_1}$ and $\hat{\beta}_{R_2}$ are the sample means of $Y$ in the two child regions, which are the optimal constant predictions by the Best Estimator Theorem. Each split creates two new child regions:
    \begin{align}
        R_1 = \{\vec{X} \mid X_j < c\}, \quad R_2 = \{\vec{X} \mid X_j \geq c\}.
    \end{align}
    Partitioning continues until the sum of squares within each region falls below a pre-specified value or the number of observations in each leaf reaches a minimum. Note that only axis-aligned splits are permitted; diagonal boundaries are not allowed.
\end{definition}

\begin{figure}[h]
    \centering
    \includegraphics[width=0.7\linewidth]{trees.jpeg}
    \caption{Tree partitioning}
    \label{fig:placeholder}
\end{figure}
 
\begin{remark}[Correctness and Limitations of Recursive Partitioning]
The algorithm commits to each split permanently before recursing on the child regions. The globally optimal partition into $K$ regions would require searching over all combinations of regions simultaneously, which is intractable. The greedy approximation has no optimality guarantee: a split that appears suboptimal locally might open up much better splits downstream, and vice versa. This is the fundamental source of instability: small perturbations to the data can change the best split at the root, which then propagates into an entirely different tree structure. Boosting and bagging address this by averaging over the instability rather than fixing it.

\end{remark}

\begin{definition}[Pruning]
    To avoid overfitting, a large tree $T_0$ is first grown, then pruned by solving a penalized loss problem. For each subtree $T \subseteq T_0$ with $K$ leaves, define the cost-complexity criterion:
    \begin{align}
        C_\alpha(T) = \sum_{k=1}^K Q_k(T) + \alpha K, \quad Q_k(T) = \sum_{\vec{X}_t \in R_k} (Y_t - \hat{\beta}_k)^2,
    \end{align}
    where $\alpha \geq 0$ is a tuning parameter selected by cross-validation. For a given $\alpha$, there exists a unique smallest subtree $T_\alpha \subseteq T_0$ minimizing $C_\alpha(T)$. Large $\alpha$ penalizes complexity heavily and yields smaller trees; small $\alpha$ permits larger trees. The optimal $T_\alpha$ is found by \emph{weakest-link pruning}: successively collapse the internal node that produces the smallest per-node increase in $\sum_k Q_k(T)$, generating a sequence of nested subtrees that contains $T_\alpha$.
\end{definition}

\begin{definition}[Variable Importance]
    The relative \textit{importance} of the $i$th predictor for a tree $T$ is
    \begin{align}
        I_i(T) = \sum_{j=1}^N \hat{\delta}^2_j \cdot \mathbf{1}(s_j = i),
    \end{align}
    where $N$ is the number of internal (parent) nodes, $s_j$ is the index of the split variable at node $j$, and $\hat{\delta}^2_j$ is the reduction in residual sum of squares from split $j$:
    \begin{align}
        \hat{\delta}^2_j = \underbrace{\sum_{t \in R_j}(Y_t - \bar{Y}_j)^2}_{\text{RSS before split}} - \left[\underbrace{\sum_{t \in R_{j,L}}(Y_t - \bar{Y}_{j,L})^2}_{\text{left child}} + \underbrace{\sum_{t \in R_{j,R}}(Y_t - \bar{Y}_{j,R})^2}_{\text{right child}}\right].
    \end{align}
    A variable is deemed important if it yields large reductions in prediction error across splits. For ensemble methods, the importance of variable $i$ is averaged over all trees: $I_i = \frac{1}{M}\sum_{m=1}^M I_i(T_m)$.
\end{definition}

\begin{definition}[Partial Dependence Plot]
    For a fitted model $\hat{f}(\vec{X})$ and a subset of predictors $\vec{X}_S$ (with complement $\vec{X}_C$), the \textit{partial dependence function} is
    \begin{align}
        \hat{f}_S(\vec{x}_S) = \mathbb{E}_{\vec{X}_C}\!\left[\hat{f}(\vec{x}_S, \vec{X}_C)\right] \approx \frac{1}{T}\sum_{t=1}^T \hat{f}(\vec{x}_S, \vec{X}_{C,t}),
    \end{align}
    where the expectation is taken over the marginal distribution of $\vec{X}_C$, estimated empirically. A partial dependence plot traces $\hat{f}_S(\vec{x}_S)$ as $\vec{x}_S$ varies, displaying the average effect of $\vec{X}_S$ on the prediction after marginalizing out all other features. This provides interpretability for ensemble methods where the individual trees are not directly readable.
\end{definition}

As previously discussed, gradient boosting builds an additive model sequentially, fitting each new weak learner to the pseudo-residuals of the current ensemble. When the weak learners are regression trees, the algorithm is as follows.

\begin{definition}[Boosted Regression Tree Algorithm]
    Given a loss function $L$ and learning rate $0 < \nu < 1$:
    \begin{enumerate}
        \item Initialize $f_0(\vec{X}) = \arg\min_c \sum_{t=1}^T L(Y_t, c)$.
        \item For $m = 1, \dots, M$:
        \begin{enumerate}
            \item Compute pseudo-residuals $r_{tm} = -\left[\frac{\partial L(Y_t, f(\vec{X}_t))}{\partial f(\vec{X}_t)}\right]_{f = f_{m-1}}$.
            \item Fit a regression tree to $\{r_{tm}\}$ with $K_m$ leaves, yielding terminal regions $R_{jm}$.
            \item For each leaf $j$, compute the optimal step:
            \begin{align}
                c_{jm} = \arg\min_c \sum_{\vec{X}_t \in R_{jm}} L\!\left[Y_t,\, f_{m-1}(\vec{X}_t) + c\right].
            \end{align}
            \item Update $f_m(\vec{X}) = f_{m-1}(\vec{X}) + \nu \sum_{j=1}^{K_m} c_{jm}\,\mathbf{1}(\vec{X} \in R_{jm})$.
        \end{enumerate}
        \item Output $\hat{f}(\vec{X}) = f_M(\vec{X})$.
    \end{enumerate}
    For squared error loss, $r_{tm} = Y_t - f_{m-1}(\vec{X}_t)$, so the tree is fit directly to the residuals.
\end{definition}

\begin{remark}
    Two tuning parameters govern the boosted tree: the tree size (controlling the order of variable interaction) and the learning rate $\nu$ paired with the number of iterations $M$. In practice, small $\nu < 0.1$ with $M$ chosen by early stopping is preferred. Writing $J$ for the number of \emph{leaves} (consistent with the notation $K$ for leaves above, but stated separately to avoid the common conflation of leaves with depth), a tree with $J$ leaves has $J-1$ splits and can capture interactions of order at most $J-1$. Hence $J=2$ (a stump) is purely additive and captures \emph{no} interactions at all; $4 \leq J \leq 8$ is the typical range, which suffices for the low-order interactions that dominate in return prediction. To speed up training and introduce regularizing randomness, \emph{stochastic gradient boosting} selects a random subsample of the data at each iteration.
\end{remark}

\begin{remark}
    A popular choice for robustness to outliers is the \textit{Huber loss}:
    \begin{align}
        L[Y_t, f(\vec{X}_t)] = \begin{cases} \frac{1}{2}[Y_t - f(\vec{X}_t)]^2 & \text{if } |Y_t - f(\vec{X}_t)| \leq \delta \\ \delta|Y_t - f(\vec{X}_t)| - \frac{1}{2}\delta^2 & \text{otherwise,} \end{cases}
    \end{align}
    where $\delta$ is taken as the $\alpha$-quantile of the absolute residuals. The Huber loss behaves like squared error for small residuals but grows linearly for large ones, limiting the influence of outliers. The pseudo-residual is then
    \begin{align}
        r_{tm} = \begin{cases} Y_t - f_{m-1}(\vec{X}_t) & \text{if } |Y_t - f_{m-1}(\vec{X}_t)| \leq \delta \\ \delta \cdot \text{sign}[Y_t - f_{m-1}(\vec{X}_t)] & \text{otherwise.} \end{cases}
    \end{align}
\end{remark}

\begin{definition}[Random Forest]
    A \textit{random forest} is a bagged ensemble of regression trees with an additional randomization over split variables, motivated by the variance reduction result of Proposition~\ref{prop:bagging}: since bagging cannot reduce variance below $\rho\sigma^2$, one must reduce $\rho$ by decorrelating the trees. The algorithm is:
    \begin{enumerate}
        \item For $b = 1, \dots, B$:
        \begin{enumerate}
            \item Draw a bootstrap sample of size $T$ from the original data (with replacement).
            \item Fit a regression tree $T_b$ to the bootstrap sample, but at each split, restrict the candidate split variables to a randomly chosen subset of $m \leq p$ variables. Grow the tree until a minimum leaf size is reached; do not prune.
        \end{enumerate}
        \item Output the ensemble prediction:
        \begin{align}
            \hat{f}^{(\text{RF})}(\vec{X}) = \frac{1}{B}\sum_{b=1}^B T_b(\vec{X}).
        \end{align}
    \end{enumerate}
    When $m = p$, this reduces to ordinary bagging. The canonical default is $m = \lfloor p/3 \rfloor$ for \emph{regression} and $m = \lfloor\sqrt{p}\rfloor$ for \emph{classification} (Breiman, 2001); since return prediction is a regression problem, $\lfloor p/3\rfloor$ is the relevant starting point, though $m$ should be tuned by out-of-bag error in practice.
\end{definition}

\begin{figure}[h]
    \centering
    \includegraphics[width=0.5\linewidth]{randomforest.png}
    \caption{Random forest}
    \label{fig:placeholder}
\end{figure}

\begin{remark}
    The number of trees $B$ can be monitored using the \emph{out-of-bag} (OOB) error: for each observation $z_t = (Y_t, \vec{X}_t^T)^T$, the OOB prediction is the average of $T_b(\vec{X}_t)$ over only those bootstrap samples that did not include $z_t$. Since each observation is excluded from roughly one-third of bootstrap samples, this provides a nearly unbiased estimate of the generalization error without a separate validation set.
\end{remark}

\begin{remark}[Tree-Based Methods in Finance]
% [CORRECTED] The original claimed trees BEAT neural networks in GKX (2020).
% This reverses their headline result: NNs win, and performance peaks at
% three hidden layers before declining ("shallow learning" finding).
    Gu, Kelly, and Xiu (2020) conduct a comprehensive comparison of machine learning methods for stock return prediction. Their central finding is that the \emph{nonlinear} methods --- tree ensembles and neural networks --- decisively outperform linear methods, and that among the nonlinear methods, neural networks come out on top. In their monthly out-of-sample results, the ranking at the top is roughly NN3 $>$ NN2 $\approx$ NN4 $>$ GBRT $\approx$ RF, with the best neural network achieving an OOS $R^2$ of around $0.4\%$ per month against approximately $0.33$--$0.34\%$ for the tree ensembles; simple OLS on the full predictor set is negative.

    Two features of this result matter more than the ranking itself. First, performance across NN1--NN5 is \emph{hump-shaped}, peaking at three hidden layers and then \emph{degrading} with further depth --- the opposite of the scaling behaviour seen in vision or language, and a direct consequence of the low signal-to-noise ratio of returns: deeper networks overfit before they can exploit compositional structure. This is the paper's ``shallow learning'' point and is the thing worth remembering. Second, the gains of both trees and networks over linear models come overwhelmingly from capturing \emph{interactions} between characteristics (momentum $\times$ volatility, size $\times$ value) and mild nonlinearities, not from any exotic representational capacity. Variable importance scores confirm that a modest number of characteristics --- price trends, liquidity, and volatility --- carry most of the predictive content, providing a data-driven check on hand-specified factor models.
\end{remark}

The bias-variance decomposition provides a unifying lens for all the estimators covered in the course. The following table summarizes the qualitative tradeoffs:

\begin{table}[H]
\centering
\renewcommand{\arraystretch}{1.3}
\begin{tabular}{@{}llll@{}}
\toprule
\textbf{Method} & \textbf{Bias} & \textbf{Variance} & \textbf{Key mechanism} \\
\midrule
OLS & Low (none if model correct) & High in high dimensions & No regularization \\
Ridge & Introduced by $\lambda$ & Reduced vs.\ OLS & Shrinks all $\hat{\beta}_j$ toward 0 \\
LASSO & Introduced by $\lambda$ & Reduced vs.\ OLS & Zeros out irrelevant $\hat{\beta}_j$ \\
Deep tree & Low (flexible) & High (unstable) & Overfits training data \\
Pruned tree & Higher & Lower & Complexity penalty $\alpha$ \\
Random Forest & Low & Low & Averaging decorrelated trees \\
Boosting & Low & Low (with tuning) & Sequential bias reduction \\
\bottomrule
\end{tabular}
\caption{Qualitative bias-variance tradeoffs across regression methods.}
\end{table}



\section{Sieves}

For $Y_t, \vec{X}_t$ having relationship $Y_t = h(\vec{X}_t) + U_t$, estimating the function $h(\cdot)$ over an infinite-dimensional function space $\mathcal{H}$,
\begin{align}
    \hat{h} = \arg \min_{h \in \mathcal{H}} \sum_{t=1}^T \left(Y_t - h(\vec{X}_t)\right)^2,
\end{align}
is infeasible because we do not have an efficient technique to search over all of $\mathcal{H}$.

\begin{remark}[Function Spaces]
    The function space $\mathcal{H}$ may take several forms:
    \begin{itemize}[noitemsep]
        \item $\mathcal{C}$: the space of continuous functions on $[a,b]$.
        \item $\mathcal{C}^2$: the space of twice continuously differentiable functions.
        \item $\mathcal{B}$: the space of Borel measurable functions.
    \end{itemize}
    Of course, one can always restrict the space. For example, imposing linearity gives $\mathcal{H} = \mathcal{L} = \{\vec{\beta}'\vec{X}_t,\, \vec{\beta} \in \mathbb{R}^p\}$, the space of all linear functions of $\vec{X}_t$. This is the simplest case: $\mathcal{L}$ is finite-dimensional with dimension $p$. Beyond this simple case, the method of sieves offers a principled way to handle the general problem.
\end{remark}

\begin{definition}[Method of Sieves]
    Pick a sequence of finite-dimensional spaces $\{\mathcal{H}_D\}$, $D = 1, 2, 3, \ldots$, that converges to $\mathcal{H}$. Define the approximating functions as
    \begin{align}
        h^D(\vec{X}_t) = \sum_{j=1}^{J_T} \beta_j h_j(\vec{X}_t),
    \end{align}
    where $h_j(\cdot)$ is the $j$th basis function for $\mathcal{H}_D$ and can be either fully known (linear sieves) or indexed by a vector of parameters $\theta_j$, such that $h_j(\vec{X}_t) = h(\vec{X}_t; \theta_j)$ (nonlinear sieves). Note that the number of basis functions $J := J_T$ depends on the sample size $T$. The \textit{sieve estimator} is then
    \begin{align}
        \hat{h}^D = \arg\min_{h \in \mathcal{H}_D} \sum_{t=1}^T \left(Y_t - h(\vec{X}_t)\right)^2.
    \end{align}
    The trick is to replace the intractable $\mathcal{H}$ by a simpler finite-dimensional space $\mathcal{H}_D$.
\end{definition}

\begin{remark}
    The dimension $D$ of the space $\mathcal{H}_D$ also depends on the sample size: $D := D_T$. When the basis functions have no free parameters, $D = J$; when they are indexed by parameters $\theta$, $D > J$, as the parameter count exceeds the number of basis functions. Popular sieves are typically compact, non-decreasing ($\mathcal{H}_T \subseteq \mathcal{H}_{T+1} \subseteq \cdots \subseteq \mathcal{H}$), and are constructed so that for any $h \in \mathcal{H}$, there exists a sequence $\pi_T h \in \mathcal{H}_T$ satisfying $d(h, \pi_T h) \to 0$ as $T \to \infty$, where $\pi_T$ is a projection map from $\mathcal{H}$ into $\mathcal{H}_T$.
\end{remark}


The construction of the approximating sequence $\{\mathcal{H}_D\}$ relies on the concept of dense sets.

\begin{definition}[Dense Set]
    If $A$ and $B$ are subsets of a metric space $\mathcal{X}$, then $A$ is \textit{dense} in $B$ if for any $\varepsilon > 0$ and any $X \in B$, there exists $Y \in A$ such that $\|X - Y\|_{\mathcal{X}} < \varepsilon$.
\end{definition}

\begin{theorem}[Stone--Weierstrass]
    Let $\mathcal{H} = \mathcal{C} = $ space of continuous functions on $[a,b]$. From the theory of approximating functions, the proper subset $\mathcal{P} \subset \mathcal{C}$ of polynomials is dense in $\mathcal{C}$. The key result is the \textit{Stone--Weierstrass theorem}: for any continuous function $m(\cdot)$ with compact support $\mathcal{X}$ on $\mathbb{R}$, and any $\varepsilon > 0$, there exists a polynomial function $P(X)$ such that $\sup_{X \in \mathcal{X}} |m(X) - P(X)| < \varepsilon$. It is therefore natural to define the sequence of approximating spaces $\mathcal{H}_D$, $D = 1, 2, 3, \ldots$ by taking $\mathcal{H}_D$ to be the set of polynomials of degree at most $D - 1$ (including a constant term), so that $\dim(\mathcal{H}_D) = D < \infty$. In the limit, this sequence converges to the infinite-dimensional space of polynomials, which is dense in $\mathcal{C}$.
\end{theorem}


When the basis functions are fully known and the approximating function is linear in the parameters, estimation reduces to familiar linear regression methods.

\begin{definition}[Polynomial Sieve]
    For $X \in \mathbb{R}$, let $\mathrm{Pol}(J_T)$ denote the space of polynomials on $[0,1]$ of degree at most $J_T$:
    \begin{align}
        \mathrm{Pol}(J_T) = \left\{ \sum_{j=0}^{J_T} \beta_j X^j,\; X \in [0,1] : \beta_j \in \mathbb{R} \right\}.
    \end{align}
    For $p = 1$, the \textit{polynomial sieve} approximates $h(X_t)$ by
    \begin{align}
        h(X_t) = \beta_0 + \beta_1 X_t + \beta_2 X_t^2 + \cdots + \beta_J X_t^J.
    \end{align}
    For $X \in \mathbb{R}^p$ with $p > 1$, the polynomial sieve is generalized by including cross terms. For $p = 2$:
    \begin{align}
        h(\vec{X}_t) = \beta_0 + \beta_1 X_1 + \beta_2 X_1^2 + \beta_3 X_2 + \beta_4 X_2^2 + \beta_5 X_1 X_2 + \cdots
    \end{align}
    For general $p > 1$, the polynomial sieve of order $J_T$ includes all monomials of total degree at most $J_T$:
    \begin{align}
        h(\vec{X}_t) = \beta_0 + \sum_{i_1=1}^{p}\beta_{i_1}X_{i_1} + \sum_{i_1=1}^{p}\sum_{i_2=i_1}^{p}\beta_{i_1 i_2}X_{i_1}X_{i_2} + \cdots + \sum_{i_1=1}^{p}\cdots\sum_{i_{J_T}=i_{J_T-1}}^{p}\beta_{i_1\cdots i_{J_T}}X_{i_1}\times\cdots\times X_{i_{J_T}}.
    \end{align}
    An alternative orthogonal basis is given by Legendre polynomials, which are orthogonal on $[-1,1]$. Rodrigues' formula gives the $j$th Legendre polynomial as $p_j(\tau) = \frac{1}{2^j j!}\frac{d^j}{d\tau^j}[(\tau^2-1)^j]$, indexed from $j=0$: $p_0 = 1$, $p_1 = \tau$, $p_2 = \frac{3\tau^2 - 1}{2}$, $p_3 = \frac{5\tau^3-3\tau}{2}$. (The original indexed these from $1$, which is inconsistent with the formula.) To model a function on $[0, T_{\max}]$, one applies the change of variable $\tau \mapsto \frac{2\tau}{T_{\max}} - 1$.
\end{definition}

\begin{proposition}[Estimation of Linear Sieves]
    When the basis functions are fully known (linear sieves), the problem is linear in the parameters. If $J \ll T$, the parameter vector $\vec{\beta} = (\beta_1, \ldots, \beta_J)'$ is estimated by OLS:
    \begin{align}
        \hat{\vec{\beta}} = \left(X_J' X_J\right)^{-1} X_J' \vec{Y},
    \end{align}
    where $X_J$ is the $T \times (J+1)$ design matrix with $(t,j)$-entry $h_j(X_t)$, and $\vec{Y} = (Y_1, \ldots, Y_T)'$. If $J$ is large relative to $T$, penalized estimators such as Ridge or LASSO can be used. The number of basis functions $J$ is selected by cross-validation or information criteria (AIC, BIC).
\end{proposition}

\begin{remark}[Curse of Dimensionality]
    For $p > 1$, the number of monomials of total degree at most $J$ grows as $\binom{p+J}{J} \sim p^J / J!$, which can become very large. In $p = 10$ dimensions with $J = 3$, there are already $\binom{13}{3} = 286$ basis functions, making LASSO regularization or orthogonal bases essential for feasible estimation.
\end{remark}

When the basis functions are themselves indexed by parameters $\theta_j$, the problem is nonlinear in the full parameter vector and requires nonlinear least squares.

\begin{definition}[Sigmoid Neural Network Sieve]
    The \textit{single hidden layer sigmoid artificial neural network (sANN) sieve} of width $J_T$ is defined as
    \begin{align}
        \mathrm{sANN}(J_T) = \left\{ \gamma_0 + \sum_{j=1}^{J_T} \alpha_j S(\vec{\gamma}_j' \vec{X} + \gamma_{0,j}) : \vec{\gamma}_j \in \mathbb{R}^p,\; \alpha_j, \gamma_0, \gamma_{0,j} \in \mathbb{R} \right\},
    \end{align}
    where $S: \mathbb{R} \to \mathbb{R}$ is a sigmoid activation function, i.e., a bounded nondecreasing function such that $\lim_{u \to -\infty} S(u) = 0$ and $\lim_{u \to +\infty} S(u) = 1$. The canonical example is the logistic function $S(x) = 1/(1+e^{-x})$.
\end{definition}

\begin{remark}[Neural Networks as Nonlinear Sieves]
    In the sANN sieve, the parameters $\vec{\gamma}_j$ and $\gamma_{0,j}$ index the basis functions themselves, making the problem nonlinear: $D > J_T$. As $J_T \to \infty$, the sANN sieve fills out the full space of continuous functions by the Universal Approximation Theorem. The dimension $J_T$ plays the same role as the polynomial degree $D$: it controls the expressiveness of the approximating space and must grow with $T$ at an appropriate rate to balance bias and variance. For nonlinear sieves, information criteria are less tractable (degrees of freedom are difficult to determine), making regularization via weight decay or Bayesian priors the preferred approach for controlling overfitting.
\end{remark}

The method of sieves generalizes naturally to semi-parametric models, which combine a finite-dimensional parametric component with an infinite-dimensional nonparametric component.

\begin{definition}[Semi-Parametric Sieve Model]
    Consider the general \textit{semi-parametric model}
% [CORRECTED] The original had the SAME regressor X_t in both the parametric
% and nonparametric terms, which is unidentified: h_0 (ranging over an
% unrestricted infinite-dimensional space) simply absorbs beta_0' X_t, so
% (beta_0, h_0) and (0, h_0 + beta_0'x) are observationally identical.
% The partially linear model (Robinson, 1988) requires DISTINCT regressors.
    \begin{align}
        Y_t = \underbrace{\vec{\beta}_0' \vec{X}_t}_{\text{parametric}} + \underbrace{h_0(\vec{Z}_t)}_{\text{nonparametric}} + U_t = m_0(\vec{X}_t, \vec{Z}_t; \theta_0) + U_t,
    \end{align}
    where $\vec{X}_t \in \mathbb{R}^p$ and $\vec{Z}_t \in \mathbb{R}^q$ are \emph{distinct} sets of regressors, and $\{U_t\}$ is a sequence of disturbances with $\mathbb{E}[U_t | \vec{X}_t, \vec{Z}_t] = 0$ and $\mathbb{E}[U_t^2 | \vec{X}_t, \vec{Z}_t] = \sigma^2(\vec{X}_t,\vec{Z}_t) < \infty$. The separation of regressors is essential, not cosmetic: if the same $\vec{X}_t$ entered both terms, then since $h_0$ ranges over an unrestricted infinite-dimensional space it would simply absorb $\vec{\beta}_0'\vec{X}_t$, making $(\vec{\beta}_0, h_0)$ and $(\vec{0}, h_0 + \vec{\beta}_0'\,\cdot)$ observationally equivalent and leaving $\vec{\beta}_0$ unidentified. This is the \emph{partially linear model} of Robinson (1988), and identification of $\vec{\beta}_0$ additionally requires that $\vec{X}_t$ not be a deterministic function of $\vec{Z}_t$, i.e.\ that $\mathrm{Var}(\vec{X}_t - \mathbb{E}[\vec{X}_t \mid \vec{Z}_t])$ be nonsingular.

    The full parameter is $\theta_0 = (\vec{\beta}_0', h_0)' \in \mathcal{B} \times \mathcal{H} \equiv \Theta$, where $\mathcal{B}$ is a finite-dimensional compact parameter space and $\mathcal{H}$ is infinite-dimensional. The population criterion is uniquely minimized at $\theta_0$:
    \begin{align}
        \theta_0 = \arg\min_{\theta \in \Theta} \mathbb{E}\!\left[\left(Y - m(\vec{X};\theta)\right)^2\right],
    \end{align}
    and the sieve estimator replaces $\Theta$ with an approximating sequence $\Theta_T$:
    \begin{align}
        \hat{\theta} = \arg\min_{\theta \in \Theta_T} \sum_{t=1}^T \left(Y_t - m(\vec{X}_t, \vec{Z}_t;\theta)\right)^2.
    \end{align}
\end{definition}


A particularly important application of sieve methods in finance is the estimation of the term structure of interest rates.

\section{Clustering}

Unlike supervised learning, clustering is an unsupervised problem: there is no response variable $y$, and the goal is to discover latent structure in the feature data $X$ alone. The central question is how to partition observations into groups such that members of the same group are more similar to one another than to members of other groups. Because there is no ground truth to evaluate against, model selection and validation are more challenging than in the supervised setting.

\begin{definition}[Dissimilarity Measure]
    Let $\vec{x}_i, \vec{x}_j \in \mathbb{R}^p$. A \textit{dissimilarity measure} $d(\vec{x}_i, \vec{x}_j) \geq 0$ quantifies how unlike two observations are, satisfying $d(\vec{x}_i, \vec{x}_i) = 0$ and the symmetry $d(\vec{x}_i, \vec{x}_j) = d(\vec{x}_j, \vec{x}_i)$. Note that a dissimilarity need not be a \emph{metric}: the most common choice, the squared Euclidean distance, violates the triangle inequality and so is not one.
    \begin{align}
        d(\vec{x}_i, \vec{x}_j) = \|\vec{x}_i - \vec{x}_j\|_2^2 = \sum_{k=1}^p (x_{ik} - x_{jk})^2.
    \end{align}
    Features should be standardized to zero mean and unit variance before computing dissimilarities, as variables measured on larger scales will otherwise dominate.
\end{definition}

\begin{definition}[$K$-Means Clustering]
    Given observations $\{\vec{x}_i\}_{i=1}^n$ and a number of clusters $K$, \textit{$K$-means clustering} seeks a partition $C: \{1,\dots,n\} \to \{1,\dots,K\}$ minimizing the within-cluster sum of squares (WCSS):
    \begin{align}
        \min_{C,\{\vec{\mu}_k\}_{k=1}^K} \sum_{k=1}^K \sum_{C(i)=k} \|\vec{x}_i - \vec{\mu}_k\|_2^2,
    \end{align}
    where $\vec{\mu}_k = \frac{1}{|C_k|}\sum_{C(i)=k}\vec{x}_i$ is the centroid of cluster $k$. This is solved via \textit{Lloyd's algorithm}:
    \begin{enumerate}
        \item Initialize $K$ centroids $\vec{\mu}_1,\dots,\vec{\mu}_K$ (e.g., by selecting $K$ observations at random).
        \item Assign each observation to its nearest centroid: $C(i) = \arg\min_k \|\vec{x}_i - \vec{\mu}_k\|_2^2$.
        \item Recompute centroids as the mean of their assigned observations.
        \item Repeat steps 2--3 until assignments do not change.
    \end{enumerate}
\end{definition}

\begin{figure}[h]
    \centering
    \includegraphics[width=0.6\linewidth]{kmeans.png}
    \caption{K-means clustering}
    \label{fig:placeholder}
\end{figure}

\begin{proposition}[Convergence of $K$-Means]
    Lloyd's algorithm decreases the WCSS objective at every step and therefore converges in a finite number of iterations.
\end{proposition}
\begin{proof}
    The assignment step minimizes the objective over $C$ with centroids held fixed: assigning each observation to its nearest centroid cannot increase the total squared distance. The update step minimizes the objective over $\{\vec{\mu}_k\}$ with $C$ held fixed: for fixed $C$, $\sum_{C(i)=k}\|\vec{x}_i - \vec{\mu}_k\|_2^2$ is a sum of squared deviations minimized by $\vec{\mu}_k = \frac{1}{|C_k|}\sum_{C(i)=k}\vec{x}_i$. Since the objective is non-negative and strictly decreases at each step unless already at a fixed point, and there are finitely many partitions, convergence is guaranteed.
\end{proof}

\begin{remark}
    $K$-means converges to a local, not necessarily global, minimum. In practice, the algorithm is run many times with different random initializations, and the solution with the lowest WCSS is retained. A smarter initialization, \emph{$K$-means++}, selects the first centroid uniformly at random, then selects each subsequent centroid with probability proportional to its squared distance from the nearest already-chosen centroid. This tends to produce better solutions and faster convergence.
\end{remark}

\begin{remark}
    Choosing $K$ is a model selection problem. A common heuristic is the \emph{elbow method}: plot the WCSS as a function of $K$ and look for the point of diminishing returns. A more principled alternative is the \emph{gap statistic}, which compares the observed WCSS to its expectation under a null reference distribution with no cluster structure, selecting the $K$ that maximizes the gap.
\end{remark}

\begin{definition}[Hierarchical Clustering]
    \textit{Hierarchical clustering} produces a nested sequence of partitions, represented as a \emph{dendrogram}, without requiring $K$ to be specified in advance. The agglomerative (bottom-up) approach begins with each observation in its own cluster and successively merges the two closest clusters until all observations form a single group. At each step, the distance between two clusters $G$ and $H$ is computed via a \emph{linkage function}:
    \begin{align}
        d_{\text{single}}(G,H) &= \min_{i \in G, j \in H} d(\vec{x}_i, \vec{x}_j), \\
        d_{\text{complete}}(G,H) &= \max_{i \in G, j \in H} d(\vec{x}_i, \vec{x}_j), \\
        d_{\text{average}}(G,H) &= \frac{1}{|G||H|}\sum_{i \in G}\sum_{j \in H} d(\vec{x}_i, \vec{x}_j).
    \end{align}
    A partition into $K$ clusters is obtained by cutting the dendrogram at the appropriate height.
\end{definition}

\begin{figure}[h]
    \centering
    \includegraphics[width=0.6\linewidth]{dendrogram.png}
    \caption{Dendrogram}
    \label{fig:placeholder}
\end{figure}

\begin{remark}
    Complete linkage tends to produce compact, roughly equal-sized clusters and is less susceptible to chaining (where a single outlier causes two large clusters to merge). Single linkage can identify elongated clusters but is sensitive to outliers. Average linkage is a compromise. The choice of linkage is substantive and should be guided by the application.
\end{remark}

\begin{definition}[Gaussian Mixture Model]
    A \textit{Gaussian mixture model (GMM)} assumes observations are drawn from a mixture of $K$ Gaussian distributions:
    \begin{align}
        p(\vec{x}) = \sum_{k=1}^K \pi_k \, \mathcal{N}(\vec{x};\vec{\mu}_k, \Sigma_k),
    \end{align}
    where $\pi_k > 0$ are the mixing proportions satisfying $\sum_k \pi_k = 1$, and $\mathcal{N}(\vec{x};\vec{\mu}_k,\Sigma_k)$ is the multivariate Gaussian density with mean $\vec{\mu}_k$ and covariance $\Sigma_k$. The parameters $\theta = \{\pi_k, \vec{\mu}_k, \Sigma_k\}_{k=1}^K$ are estimated by maximum likelihood.
\end{definition}

\begin{definition}[Expectation-Maximization Algorithm]
    Since the component assignments are unobserved, direct maximization of the GMM log-likelihood is intractable. The \textit{EM algorithm} maximizes the expected complete-data log-likelihood by iterating two steps. Let $r_{ik} = \mathbb{P}(Z_i = k \mid \vec{x}_i, \theta)$ be the \emph{responsibility} of component $k$ for observation $i$:
    \begin{enumerate}
        \item E-step: Compute responsibilities given current parameters:
        \begin{align}
            r_{ik} = \frac{\pi_k \, \mathcal{N}(\vec{x}_i; \vec{\mu}_k, \Sigma_k)}{\sum_{j=1}^K \pi_j \, \mathcal{N}(\vec{x}_i; \vec{\mu}_j, \Sigma_j)}.
        \end{align}
        \item M-step: Update parameter estimates using the weighted sufficient statistics. Let $n_k = \sum_i r_{ik}$ and define the update with the following prior class probability, mean, and covariance: 
        \begin{align}
            \hat{\pi}_k &= \frac{n_k}{n}, \quad
            \hat{\vec{\mu}}_k = \frac{1}{n_k}\sum_i r_{ik}\vec{x}_i, \quad
            \hat{\Sigma}_k = \frac{1}{n_k}\sum_i r_{ik}(\vec{x}_i - \hat{\vec{\mu}}_k)(\vec{x}_i - \hat{\vec{\mu}}_k)^T.
        \end{align}
    \end{enumerate}
    The EM algorithm is guaranteed to increase the observed-data log-likelihood at each iteration.
\end{definition}

\begin{figure}[h]
    \centering
    \includegraphics[width=0.7\linewidth]{em_algo.jpeg}
    \caption{EM algorithm}
    \label{fig:placeholder}
\end{figure}

\begin{proposition}[EM Monotonically Increases Log-Likelihood]
    Let $\theta^{(t)}$ and $\theta^{(t+1)}$ be the parameters before and after one EM iteration. Then $\ell(\theta^{(t+1)}) \geq \ell(\theta^{(t)})$, where $\ell(\theta) = \sum_i \log p(\vec{x}_i;\theta)$ is the observed-data log-likelihood.
\end{proposition}
\begin{proof}[Proof sketch]
    Define the complete-data log-likelihood $\ell_c(\theta) = \sum_i \log p(\vec{x}_i, Z_i; \theta)$ and the auxiliary function $Q(\theta|\theta^{(t)}) = \mathbb{E}_{Z|X,\theta^{(t)}}[\ell_c(\theta)]$. The E-step computes $Q(\theta|\theta^{(t)})$, and the M-step chooses $\theta^{(t+1)} = \arg\max_\theta Q(\theta|\theta^{(t)})$. By Jensen's inequality applied to the KL divergence between the posterior $p(Z|X,\theta^{(t)})$ and $p(Z|X,\theta)$, it can be shown that $\ell(\theta) \geq Q(\theta|\theta^{(t)}) + H(\theta^{(t)})$ for an entropy term $H$ that does not depend on $\theta$. Since the M-step ensures $Q(\theta^{(t+1)}|\theta^{(t)}) \geq Q(\theta^{(t)}|\theta^{(t)})$, the observed log-likelihood cannot decrease.
\end{proof}

\begin{remark}
    $K$-means is a special case of EM for GMMs where all covariance matrices are constrained to $\Sigma_k = \sigma^2 I$ and responsibilities are hard (i.e., each observation is assigned to exactly one component). GMMs provide a probabilistic, soft-assignment generalization, and their log-likelihood can be used with BIC for model selection over $K$.
\end{remark}

\begin{remark}[Clustering in Finance]
    GMMs are used in finance for \emph{regime detection}: fitting a 2- or 3-component mixture to market returns identifies ``bull,'' ``bear,'' and ``crisis'' regimes with distinct mean returns and volatilities. The posterior probabilities $r_{ik}$ give a real-time estimate of the current regime, which can be used to adjust portfolio allocations (e.g., reducing equity exposure when the ``crisis'' component has high responsibility). Hierarchical clustering of the correlation matrix of assets is used for portfolio diversification: assets within the same cluster are substitutes, so a well-diversified portfolio should hold assets from different clusters. This ``hierarchical risk parity'' approach has gained traction as an alternative to mean-variance optimization.
\end{remark}

\section{Classification}

Classification is the supervised learning problem where the response $Y$ takes values in a finite set of classes $\{1,\dots,K\}$. Logistic regression (Section~2) provides one approach for $K=2$.

\begin{definition}[Bayes Classifier]
    The \textit{Bayes classifier} assigns each observation $\vec{x}$ to the class with highest posterior probability:
    \begin{align}
        \hat{C}(\vec{x}) = \arg\max_k \mathbb{P}(Y = k \mid \vec{X} = \vec{x}).
    \end{align}
    By Bayes' theorem, $\mathbb{P}(Y=k|\vec{x}) \propto \mathbb{P}(Y=k)\mathbb{P}(\vec{x}|Y=k)=\pi_k f_k(\vec{x})$, where $\pi_k = \mathbb{P}(Y=k)$ is the prior class probability and $f_k(\vec{x})$ is the class-conditional density $\mathbb{P}(\vec{x}|Y=k)$. The Bayes classifier achieves the lowest possible misclassification rate (the Bayes error rate), but is infeasible in practice since $f_k$ is unknown.
\end{definition}

\begin{proposition}[Optimality of the Bayes Classifier]
    The Bayes classifier minimizes the misclassification probability among all classifiers.
\end{proposition}
\begin{proof}
    Let $C^*(\vec{x}) = \arg\max_k \mathbb{P}(Y=k|\vec{x})$ be the Bayes classifier and $\hat{C}$ any other classifier. The misclassification probability of $\hat{C}$ is
    \begin{align}
        \mathbb{P}(\hat{C}(\vec{X}) \neq Y) &= \mathbb{E}\left[\mathbb{P}(\hat{C}(\vec{X}) \neq Y | \vec{X})\right] = \mathbb{E}\left[1 - \mathbb{P}(Y = \hat{C}(\vec{X})|\vec{X})\right].
    \end{align}
    Since $C^*$ maximizes $\mathbb{P}(Y = k|\vec{X})$ pointwise, $\mathbb{P}(Y = C^*(\vec{X})|\vec{X}) \geq \mathbb{P}(Y = \hat{C}(\vec{X})|\vec{X})$ for all $\vec{X}$, so $\mathbb{P}(C^*(\vec{X}) \neq Y) \leq \mathbb{P}(\hat{C}(\vec{X}) \neq Y)$.
\end{proof}

\begin{definition}[Linear Discriminant Analysis]
    \textit{Linear discriminant analysis (LDA)} estimates the Bayes classifier by modeling each class-conditional density as Gaussian with a shared covariance matrix:
    \begin{align}
        \vec{X} \mid Y = k \sim \mathcal{N}(\vec{\mu}_k, \Sigma).
    \end{align}
    Parameters are estimated by MLE: $\hat{\pi}_k = n_k/n$, $\hat{\vec{\mu}}_k = \frac{1}{n_k}\sum_{C(i)=k}\vec{x}_i$, and $\hat{\Sigma} = \frac{1}{n-K}\sum_k\sum_{C(i)=k}(\vec{x}_i - \hat{\vec{\mu}}_k)(\vec{x}_i - \hat{\vec{\mu}}_k)^T$. Then, examine the log-posteriors from the Bayesian $\mathbb{P}(Y=k|\vec{x}) \propto \pi_k f_k(\vec{x})$:
    \begin{align}
        \ln\left(\pi_k \frac{1}{(2\pi)^{p/2}|\Sigma_k|^{1/2}}e^{-\frac 12 (\vec{x}-\vec{\mu}_k)^T\Sigma_k^{-1}(\vec{x}-\vec{\mu}_k)}\right) = \ln{\pi_k}- \ln\left((2\pi)^{p/2}|\Sigma_k|^{1/2}\right)-\frac 12 (\vec{x}-\vec{\mu}_k)^T\Sigma_k^{-1}(\vec{x}-\vec{\mu}_k).
    \end{align}
    The Bayesian classifier assigns the data point to the class with the highest posterior probability, so neglect terms that are identical across classes. Since the covariance matrix is shared, $\Sigma_k = \Sigma$. The LDA decision rule thus assigns $\vec{x}$ to the class $k$ maximizing the \emph{linear discriminant function}:
    \begin{align}
        \delta_k(\vec{x}) = \vec{x}^T\hat{\Sigma}^{-1}\hat{\vec{\mu}}_k - \frac{1}{2}\hat{\vec{\mu}}_k^T\hat{\Sigma}^{-1}\hat{\vec{\mu}}_k + \log\hat{\pi}_k.
    \end{align}
\end{definition}

\begin{figure}[h]
    \centering
    \includegraphics[width=0.5\linewidth]{linear_disc.jpeg}
    \caption{Linear discriminant analysis, showing projected Gaussians with the same variance}
    \label{fig:placeholder}
\end{figure}

\begin{proposition}[LDA Decision Boundary]
    The LDA decision boundary between classes $k$ and $l$ is linear in $\vec{x}$.
\end{proposition}
\begin{proof}
    The boundary is the set of $\vec{x}$ where $\delta_k(\vec{x}) = \delta_l(\vec{x})$, i.e.,
    \begin{align}
        \vec{x}^T\hat{\Sigma}^{-1}(\hat{\vec{\mu}}_k - \hat{\vec{\mu}}_l) = \frac{1}{2}(\hat{\vec{\mu}}_k^T\hat{\Sigma}^{-1}\hat{\vec{\mu}}_k - \hat{\vec{\mu}}_l^T\hat{\Sigma}^{-1}\hat{\vec{\mu}}_l) - \log\frac{\hat{\pi}_k}{\hat{\pi}_l}.
    \end{align}
    The left side is linear in $\vec{x}$ and the right side is a scalar constant, so this defines a hyperplane.
\end{proof}

\begin{example}[LDA for Market Regime Classification]
    An analyst classifies monthly market conditions as ``bull'' ($Y=1$) or ``bear'' ($Y=0$) using two features: the trailing 3-month market return ($x_1$) and the VIX level ($x_2$). From historical data: $\hat{\vec{\mu}}_1 = (3.5, 15)^T$ (bull), $\hat{\vec{\mu}}_0 = (-2.1, 28)^T$ (bear), $\hat{\pi}_1 = 0.7$, $\hat{\pi}_0 = 0.3$. The shared covariance is $\hat{\Sigma} = \left(\begin{smallmatrix} 12 & -8 \\ -8 & 25 \end{smallmatrix}\right)$. For a new month with $\vec{x} = (1.0, 20)^T$, computing $\delta_1(\vec{x})$ and $\delta_0(\vec{x})$ and comparing them yields the classification. The decision boundary $\delta_1(\vec{x}) = \delta_0(\vec{x})$ is a line in the (return, VIX) plane: roughly, months with high returns and low VIX are classified as bull markets.
\end{example}

\begin{definition}[Quadratic Discriminant Analysis]
    \textit{Quadratic discriminant analysis (QDA)} relaxes the shared covariance assumption of LDA, allowing each class its own covariance $\Sigma_k$:
    \begin{align}
        \vec{X} \mid Y = k \sim \mathcal{N}(\vec{\mu}_k, \Sigma_k).
    \end{align}
    The Bayesian classifier method of assigning the data point to the class with the highest posterior probability still holds, but now terms including $\Sigma_k$ cannot be neglected. The decision rule hence assigns $\vec{x}$ to the class maximizing:
    \begin{align}
        \delta_k(\vec{x}) = -\frac{1}{2}\log|\Sigma_k| - \frac{1}{2}(\vec{x}-\vec{\mu}_k)^T\Sigma_k^{-1}(\vec{x}-\vec{\mu}_k) + \log\pi_k.
    \end{align}
    The decision boundary between any two classes is now quadratic in $\vec{x}$. QDA requires estimating $K$ separate covariance matrices, which is more flexible but has higher variance than LDA.
\end{definition}

\begin{remark}
    LDA and QDA represent a bias-variance tradeoff: LDA imposes a strong constraint (shared $\Sigma$) that reduces variance at the cost of potential bias if class covariances differ substantially. LDA is preferred when $n$ is small relative to $p$; QDA is preferred when class covariances differ and $n$ is large enough to estimate them reliably.
\end{remark}

\begin{definition}[Confusion Matrix and Performance Metrics]
    For a binary classifier with $Y \in \{0,1\}$, the \emph{confusion matrix} tabulates the four possible outcomes:
    \begin{align}
        \text{Accuracy} &= \frac{TP + TN}{TP + TN + FP + FN}, \\
        \text{Precision} &= \frac{TP}{TP + FP}, \qquad \text{Recall} = \frac{TP}{TP + FN}, \\
        F_1 &= \frac{2 \cdot \text{Precision} \cdot \text{Recall}}{\text{Precision} + \text{Recall}},
    \end{align}
    where TP, TN, FP, FN denote true positives, true negatives, false positives, and false negatives respectively. The \emph{receiver operating characteristic} (ROC) curve plots the true positive rate (Recall) against the false positive rate $FP/(FP+TN)$ as the classification threshold parametrically varies. The \emph{area under the ROC curve} (AUC) summarizes classifier performance across all thresholds, with AUC $= 1$ indicating a perfect classifier and AUC $= 0.5$ indicating random guessing.
\end{definition}

\begin{remark}
    Accuracy is a misleading metric when class frequencies are imbalanced; a classifier that predicts the majority class for every observation can achieve high accuracy while completely failing to identify the minority class. Precision, recall, $F_1$, and AUC are more informative in such settings, and the choice among them depends on the relative costs of false positives and false negatives in the application.
\end{remark}

The classification methods covered so far model the conditional probability $\mathbb{P}(Y=1|\vec{X})$ explicitly. Support vector methods take a different route: they search directly for the decision boundary that best separates the classes, without committing to a full probabilistic model. The resulting estimators are robust in high-dimensional settings where $d \gg n$.

\begin{definition}[Linear Classifier]
    Let $Y \in \{-1, +1\}$ and $\vec{X} \in \mathbb{R}^d$. A \textit{linear classifier} takes the form
    \begin{align}
        f(\vec{X};\vec{w},b) = \mathrm{sgn}(\vec{w}\cdot\vec{X} + b),
    \end{align}
    where $\mathrm{sgn}(z) = +1$ if $z \geq 0$ and $-1$ otherwise, $\vec{w} \in \mathbb{R}^d$ is the normal vector to the decision hyperplane, and $b \in \mathbb{R}$ is the offset. The hyperplane $\{\vec{x} : \vec{w}\cdot\vec{x} + b = 0\}$ partitions $\mathbb{R}^d$ into a positive half-space ($\vec{w}'\vec{x} + b > 0$) and a negative half-space ($\vec{w}\cdot\vec{x} + b < 0$). The classifier assigns $\hat{Y} = +1$ to observations on the positive side and $\hat{Y} = -1$ to those on the negative side.
\end{definition}

\begin{definition}[Linear Separability and Support Vectors]
    A training set $\{(\vec{X}_i, Y_i)\}_{i=1}^n$ is \emph{linearly separable} if there exists $(\vec{w},b)$ such that $\vec{w} \cdot \vec{X}_i + b > 0$ whenever $Y_i = +1$ and $\vec{w} \cdot\vec{X}_i + b < 0$ whenever $Y_i = -1$. Equivalently, $Y_i(\vec{w} \cdot \vec{X}_i + b) > 0$ for all $i$. In this case, logit classification blows up in $\beta$, so alternative approaches must be employed.

    For a given $\vec{w}$, the distance from a point $\vec{X}$ to the hyperplane is $|\vec{w}\cdot\vec{X}+b|/\|\vec{w}\|$. The \emph{positive support vector} $\vec{X}^+(\vec{w})$ is the positive-class point closest to the hyperplane, and the \emph{negative support vector} $\vec{X}^-(\vec{w})$ is the negative-class point closest to it. By rescaling $(\vec{w},b)$ (which leaves the hyperplane unchanged), we may normalize so that
    \begin{align}
        \vec{w}\cdot\vec{X}^+ + b = +1, \qquad \vec{w}\cdot\vec{X}^- + b = -1.
    \end{align}
    Under this normalization, the \emph{margin}, the distance between the two margin hyperplanes $\{\vec{w}\cdot\vec{x}+b=\pm 1\}$, equals
    \begin{align}
        m(\vec{w}) = \frac{\vec{w}\cdot(\vec{X}^+ - \vec{X}^-)}{\|\vec{w}\|} = \frac{2}{\|\vec{w}\|}.
    \end{align}
\end{definition}

\begin{figure}[h]
    \centering
    \includegraphics[width=0.6\linewidth]{svm.png}
    \caption{Hard- and soft-margin SVM}
    \label{fig:placeholder}
\end{figure}

\begin{theorem}[Hard-Margin SVM]
    The \emph{hard-margin SVM} maximizes the margin over all perfectly-separating hyperplanes, or equivalently solves the quadratic program
    \begin{align}
        \min_{\vec{w},b} \frac{1}{2}\|\vec{w}\|^2 \quad \text{subject to} \quad Y_i(\vec{w}\cdot\vec{X}_i + b) \geq 1, \quad \forall i.
    \end{align}
    If a solution exists it is unique and is called the \emph{optimal separating hyperplane}. The solution satisfies
    \begin{align}
        \vec{w}_0 = \sum_{i=1}^n \alpha_i^0 Y_i \vec{X}_i,
    \end{align}
    where $\alpha_i^0 \geq 0$ are Lagrange multipliers that are nonzero only for the support vectors (observations on the margin boundary, where $Y_i(\vec{w}_0\cdot\vec{X}_i + b_0) = 1$).
\end{theorem}
\begin{proof}
    Introduce Lagrange multipliers $\alpha_i \geq 0$ for each constraint and form the Lagrangian
    \begin{align}
        \mathcal{L}(\vec{w},b,\vec{\alpha}) = \frac{1}{2}\vec{w}\cdot\vec{w} - \sum_{i=1}^n \alpha_i[Y_i(\vec{w}\cdot\vec{X}_i+b)-1].
    \end{align}
    The KKT stationarity conditions in $(\vec{w},b)$ give
    \begin{align}
        \frac{\partial \mathcal{L}}{\partial \vec{w}} = \vec{w} - \sum_i \alpha_i Y_i \vec{X}_i = 0, \qquad \frac{\partial \mathcal{L}}{\partial b} = -\sum_i \alpha_i Y_i = 0,
    \end{align}
    so $\vec{w}_0 = \sum_i \alpha_i Y_i \vec{X}_i$ and $\sum_i \alpha_i Y_i = 0$. Substituting these back into $\mathcal{L}$ yields the dual objective
    \begin{align}
        W(\vec{\alpha}) = \sum_{i=1}^n \alpha_i - \frac{1}{2}\sum_{i=1}^n\sum_{j=1}^n \alpha_i\alpha_j Y_i Y_j \vec{X}_i^T\vec{X}_j,
    \end{align}
    which is maximized over $\alpha_i \geq 0$, $\sum_i \alpha_i Y_i = 0$. By complementary slackness, $\alpha_i[Y_i(\vec{w}_0\cdot\vec{X}_i+b_0)-1]=0$ for all $i$, so $\alpha_i > 0$ only for support vectors. The optimal classifier is
    \begin{align}
        f(\vec{X};\vec{w}_0,b_0) = \mathrm{sgn}\!\left(\sum_{i=1}^n \alpha_i^0 Y_i \vec{X}_i^T\vec{X} + b_0\right).
    \end{align}
\end{proof}

\begin{remark}[Why Large Margins Generalize]
    The margin $2/\|\vec{w}_0\|$ measures how robustly the training data are separated: a large margin means that small perturbations of the data points do not alter the classification. This geometric stability translates directly into better out-of-sample prediction; the hard-margin SVM implicitly regularizes by constraining the norm $\|\vec{w}\|$, in the same spirit as Ridge regression.
\end{remark}

When the data are not linearly separable, no feasible solution to the hard-margin problem exists. The \emph{soft-margin SVM} relaxes the constraints by introducing slack variables.

\begin{definition}[Soft-Margin SVM]
    Introduce nonnegative slack variables $\xi_i \geq 0$ that measure the degree of constraint violation, and solve
    \begin{align}
        \min_{\vec{w},b,\xi_i \geq 0} \frac{1}{2}\|\vec{w}\|^2 + C\sum_{i=1}^n \xi_i \quad \text{subject to} \quad Y_i(\vec{w}\cdot\vec{X}_i + b) \geq 1 - \xi_i, \quad \forall i.
    \end{align}
    The regularization parameter $C \in [0,\infty]$ controls the bias-variance tradeoff:
    \begin{itemize}[noitemsep]
        \item $C \to 0$: all constraints are effectively dropped (any hyperplane is admissible);
        \item small $C$: wide soft margin, tolerating many misclassifications;
        \item large $C$: narrow soft margin, penalizing misclassification heavily;
        \item $C \to \infty$: recovers the hard-margin problem.
    \end{itemize}
    If a solution exists it is unique and is called the \emph{optimal generalized hyperplane}.
\end{definition}

\begin{proposition}[Equivalence to Hinge Loss Minimization]
    The soft-margin SVM is equivalent to the unconstrained minimization of the \emph{hinge loss} plus an $\ell^2$ regularization penalty. Setting $\lambda = 1/(2nC)$,
    \begin{align}
        \min_{(\vec{w},b) \in \mathbb{R}^{d+1}} \frac{1}{n}\sum_{i=1}^n \ell(\vec{X}_i, Y_i;\vec{w},b) + \lambda\|\vec{w}\|^2, \qquad \ell(\vec{X}_i,Y_i;\vec{w},b) := \max\!\left(0,\, 1 - Y_i(\vec{w}\cdot\vec{X}_i + b)\right).
    \end{align}
    The hinge loss is zero for correctly classified points lying beyond the margin and grows linearly with the distance of misclassified (or within-margin) points from the correct margin hyperplane.
\end{proposition}
\begin{proof}
    Since the objective is increasing in $\xi_i$ for $C \neq 0$, at optimality $\xi_i = \max(0, 1 - Y_i(\vec{w}\cdot\vec{X}_i+b))$. Substituting eliminates $\xi_i$ from the problem, giving the unconstrained form. Dividing through by $2nC$ and setting $\lambda = 1/(2nC)$ yields the stated expression.
\end{proof}


\begin{definition}[Kernel Trick]
    The SVM extends to nonlinear classifiers via a two-step procedure:
    \begin{enumerate}[noitemsep]
        \item Map $\vec{X} \in \mathbb{R}^d$ to an expanded feature space $\vec{Z} = \varphi(\vec{X}) \in \mathbb{R}^K$, $K \gg d$, using a fixed nonlinear transformation $\varphi$.
        \item Find the optimal separating hyperplane in the feature space $\{(\varphi(\vec{X}_i), Y_i)\}_{i=1}^n$.
    \end{enumerate}
    The key observation is that both the dual objective $W(\vec{\alpha})$ and the optimal classifier depend on inputs only through inner products $\langle\varphi(\vec{X}_i),\varphi(\vec{X}_j)\rangle$. A \emph{kernel function} $K(\vec{u},\vec{v}) = \langle\varphi(\vec{u}),\varphi(\vec{v})\rangle$ computes this inner product in the feature space without explicitly constructing $\varphi$. The dual objective becomes
    \begin{align}
        W(\vec{\alpha}) = \sum_i \alpha_i - \frac{1}{2}\sum_{i,j} \alpha_i\alpha_j Y_i Y_j K(\vec{X}_i,\vec{X}_j),
    \end{align}
    and the classifier is $f(\vec{X}) = \mathrm{sgn}\!\left(\sum_i \alpha_i^0 Y_i K(\vec{X}_i, \vec{X}) + b_0\right)$.
\end{definition}

\begin{theorem}[Mercer's Theorem]
    A continuous symmetric function $K(\vec{u},\vec{v})$ corresponds to an inner product $\langle\varphi(\vec{u}),\varphi(\vec{v})\rangle$ in some feature space if and only if $K$ is positive semi-definite: for all $n$, all $\vec{X}_1,\ldots,\vec{X}_n$, and all $c_1,\ldots,c_n \in \mathbb{R}$,
    \begin{align}
        \sum_{i=1}^n\sum_{j=1}^n c_i c_j K(\vec{X}_i,\vec{X}_j) \geq 0.
    \end{align}
    This theorem justifies choosing a kernel directly: any positive semi-definite $K$ implicitly defines a valid (possibly infinite-dimensional) feature map $\varphi$, so one need never construct $\varphi$ explicitly.
\end{theorem}

\begin{figure}[h]
    \centering
    \includegraphics[width=0.5\linewidth]{kernel.png}
    \caption{Data becomes linearly separable under the kernel mapping}
    \label{fig:placeholder}
\end{figure}

\begin{definition}[Common Kernel Functions]
    The most widely used kernels are:
    \begin{itemize}[noitemsep]
        \item \textbf{Linear:} $K(\vec{u},\vec{v}) = \vec{u}\cdot\vec{v}$. No transformation of features; equivalent to the original SVM. Works well in high-dimensional settings with low overfitting risk.
        \item \textbf{Polynomial of degree $k$:} $K(\vec{u},\vec{v}) = (1 + \vec{u}\cdot\vec{v})^k$. Captures all monomial interactions up to degree $k$. For $p = 2, k = 2$, the induced feature map is $\varphi(\vec{x}) = [1,\, \sqrt{2}x_1,\, \sqrt{2}x_2,\, \sqrt{2}x_1x_2,\, x_1^2,\, x_2^2]^T$.
        \item \textbf{Gaussian RBF:} $K(\vec{u},\vec{v}) = \exp(-\gamma\|\vec{u}-\vec{v}\|_2^2)$. Highly flexible; each support vector generates a Gaussian bump centered at $\vec{X}_i$, and the decision boundary is a sum of localized effects. Overfits when $\gamma$ is large (very narrow bumps). The induced feature space is infinite-dimensional.
        \item \textbf{Neural network:} $K(\vec{u},\vec{v}) = \tanh(\gamma_1 \vec{u}\cdot\vec{v} + \gamma_2)$.
    \end{itemize}
    Kernel choice is fundamentally an empirical problem: it is equivalent to choosing a notion of similarity between observations. The recommended strategy is to start with a linear kernel and increase flexibility only when cross-validation confirms improvement.
\end{definition}

The SVM framework extends naturally to regression of a real-valued response $Y$ via a modified loss function that is insensitive to small errors.

\begin{definition}[$\varepsilon$-Insensitive Loss and SVR]
    The \emph{support vector regression} (SVR) estimator minimizes
    \begin{align}
        H(\vec{w}, b_0) = \sum_{i=1}^n V_\varepsilon(y_i - f(\vec{x}_i)) + \frac{\lambda}{2}\|\vec{w}\|^2, \qquad f(\vec{x}) = \vec{w}\cdot\vec{x} + b_0,
    \end{align}
    where the \emph{$\varepsilon$-insensitive loss} is
    \begin{align}
        V_\varepsilon(r) = \begin{cases} 0 & \text{if } |r| < \varepsilon, \\ |r| - \varepsilon & \text{otherwise.} \end{cases}
    \end{align}
    The loss is zero for residuals smaller than $\varepsilon$ (analogous to correctly classified points in the SVM) and grows linearly beyond that threshold (analogous to the hinge loss for soft-margin misclassifications). The $\ell^2$ term $\frac{\lambda}{2}\|\vec{w}\|^2$ is a Ridge-type regularization.
\end{definition}

\begin{remark}
    The two parameters $\varepsilon$ and $\lambda$ play distinct roles. The parameter $\varepsilon$ determines which observations are negligible in estimating $(\vec{w},b_0)$: larger $\varepsilon$ means fewer support vectors, since more residuals fall inside the insensitive band. Because $V_\varepsilon$ is sensitive to the scale of $y$, practitioners typically normalize $y$ to unit variance before fitting. The regularization parameter $\lambda$ plays the usual bias-variance role and is selected by cross-validation.
\end{remark}

\begin{theorem}[SVR Solution and Sparsity]
    The minimizers $\hat{\vec{w}}$ and $\hat{b}_0$ of $H$ admit the representation
    \begin{align}
        \hat{\vec{w}} = \sum_{i=1}^n (\hat{\alpha}_i^* - \hat{\alpha}_i)\vec{x}_i, \qquad \hat{f}(\vec{x}) = \sum_{i=1}^n (\hat{\alpha}_i^* - \hat{\alpha}_i)\langle\vec{x},\vec{x}_i\rangle + b_0,
    \end{align}
    where $\hat{\alpha}_i^*, \hat{\alpha}_i \geq 0$ solve the quadratic dual program
    \begin{align}
        \min_{\alpha^*,\alpha} \; \varepsilon\sum_{i=1}^n(\alpha_i^* + \alpha_i) - \sum_{i=1}^n y_i(\alpha_i^* - \alpha_i) + \frac{1}{2}\sum_{i,i'=1}^n (\alpha_i^* - \alpha_i)(\alpha_{i'}^* - \alpha_{i'})\langle\vec{x}_i,\vec{x}_{i'}\rangle,
    \end{align}
    subject to $0 \leq \alpha_i, \alpha_i^* \leq 1/\lambda$, $\sum_i(\alpha_i^*-\alpha_i) = 0$, $\alpha_i\alpha_i^* = 0$. The complementary slackness condition $\alpha_i\alpha_i^* = 0$ guarantees that only a sparse subset of observations have $\hat{\alpha}_i^* - \hat{\alpha}_i \neq 0$; these are the \emph{support vectors of the regression}.
\end{theorem}

\begin{remark}[SVR vs.\ Ridge in High Dimensions]
    When $d \gg n$, Ridge regression estimates a $d$-dimensional coefficient vector $\vec{\beta}$ by shrinking all components toward zero. SVR exchanges this for estimating at most $n$ dual weights $(\hat{\alpha}_i^* - \hat{\alpha}_i)$ attached to the observed feature vectors $\{\vec{x}_i\}$: the effective dimension of the problem drops from $d$ to (at most) $n$. SVR thus has a flavor of elastic net, but instead of eliminating individual covariates from $\vec{x}$ it eliminates entire observation vectors from the sum. This makes SVR particularly attractive in text and NLP applications where $d$ (vocabulary size) can be orders of magnitude larger than $n$ (number of time periods).
\end{remark}

\begin{definition}[Nonlinear SVR via Kernels]
    Since the SVR prediction $\hat{f}(\vec{x}) = \sum_i (\hat{\alpha}_i^* - \hat{\alpha}_i)\langle\vec{x},\vec{x}_i\rangle + b_0$ and the dual objective both depend on inputs only through inner products $\langle\vec{x}_i,\vec{x}_j\rangle$, the kernel trick applies directly. Replacing the inner product by a kernel $K(\vec{x},\vec{x}') = \langle\varphi(\vec{x}),\varphi(\vec{x}')\rangle$ yields the nonlinear SVR predictor
    \begin{align}
        \hat{f}(\vec{x}) = \sum_{i=1}^n (\hat{\alpha}_i^* - \hat{\alpha}_i) K(\vec{x},\vec{x}_i) + b_0.
    \end{align}
    Common choices are the polynomial kernel $(1 + \langle\vec{x},\vec{x}'\rangle)^k$ and the Gaussian RBF $\exp(-\gamma\|\vec{x}-\vec{x}'\|^2)$.
\end{definition}


\begin{remark}[SVM and SVR in Finance]
    The SVM and SVR inherit the regularization philosophy of Ridge and LASSO but express it geometrically: the margin constraint (SVM) or the $\varepsilon$-insensitive band (SVR) introduces sparsity in the dual rather than the primal, automatically selecting the most informative observations as support vectors. This dual sparsity is particularly valuable in high-dimensional text and alternative-data applications where $d \gg n$ is the rule rather than the exception. The kernel trick then enables nonlinear extensions at essentially no additional computational cost, since the dual problem always has dimension $n$ regardless of the feature space dimensionality.

    One limitation relative to LASSO or elastic net is that SVR does not perform variable (covariate) selection: the solution $\hat{\vec{w}} = \sum_t \hat{\alpha}_t \vec{x}_t$ uses all $d$ covariates for each support vector $\vec{x}_t$, so inference on individual word coefficients is not directly available. Standard errors are also difficult to compute, as the dual representation does not map to a standard regression framework.
\end{remark}

\section{Neural Networks}

Neural networks formalize and extend the sieve estimation framework of Section~7. A single hidden layer sigmoid neural network (sANN) was already introduced as a canonical nonlinear sieve; here we develop the architecture, estimation, and regularization theory in detail, and connect the approach to the empirical asset-pricing and option-pricing literatures.

\begin{definition}[Feed-Forward Neural Network]
    A \textit{single hidden layer feed-forward neural network} with $J_T$ hidden units approximates an unknown function $f : \mathbb{R}^p \to \mathbb{R}$ by
    \begin{align}
        H(\vec{X};\theta) = \beta_0 + \sum_{j=1}^{J_T} \beta_j S(\vec{\gamma}_j' \vec{X} + \gamma_{0,j}),
    \end{align}
    where:
    \begin{itemize}[noitemsep]
        \item $\vec{X} \mapsto \vec{\gamma}_j'\vec{X} + \gamma_{0,j}$ is an affine transformation (linear combination plus shift) of the input vector;
        \item $S : \mathbb{R} \to \mathbb{R}$ is an \emph{activation function};
        \item $\theta := (\beta_0, \beta_1, \ldots, \beta_{J_T}, \vec{\gamma}_1', \ldots, \vec{\gamma}_{J_T}', \gamma_{0,1}, \ldots, \gamma_{0,J_T})'$ is the full vector of parameters (called \emph{weights}).
    \end{itemize}
    The scalars $\beta_0$ and $\gamma_{0,j}$ are called \emph{biases}; $\vec{\gamma}_j$ controls the slope of $S$ and $\gamma_{0,j}$ shifts the curve horizontally.
\end{definition}

\begin{remark}[Connection to Sieves]
    The sANN is precisely the nonlinear sieve: the basis functions $h_j(\vec{X}) = S(\vec{\gamma}_j\cdot\vec{X} + \gamma_{0,j})$ are parameterized by $(\vec{\gamma}_j, \gamma_{0,j})$, making the problem nonlinear. The \textit{Universal Approximation Theorem} guarantees that, as $J_T \to \infty$, the class $\mathrm{sANN}(J_T)$ is dense in the space of continuous functions (Cybenko, 1989; Hornik, Stinchcombe, and White, 1989).
\end{remark}

The network can be represented in matrix form. Define augmented vectors $\widetilde{\vec{X}}_t = (1, \vec{X}_t^T)^T$ and $\widetilde{\vec{\gamma}}_j = (\gamma_{0,j}, \vec{\gamma}_j^T)^T$, and assemble them into matrices $X \in \mathbb{R}^{T \times (p+1)}$ and $\Gamma \in \mathbb{R}^{(p+1) \times J_T}$. Apply the activation function elementwise to form the \emph{output matrix}
\begin{align}
    O(X\Gamma)_{[T \times (J_T+1)]} =
    \begin{pmatrix}
        1 & S(\widetilde{\vec{\gamma}}_1\cdot\widetilde{\vec{X}}_1) & \cdots & S(\widetilde{\vec{\gamma}}_{J_T}\cdot\widetilde{\vec{X}}_1) \\
        \vdots & \vdots & \ddots & \vdots \\
        1 & S(\widetilde{\vec{\gamma}}_1\cdot\widetilde{\vec{X}}_T) & \cdots & S(\widetilde{\vec{\gamma}}_{J_T}\cdot\widetilde{\vec{X}}_T)
    \end{pmatrix}.
\end{align}
Setting $\vec{\beta} = (\beta_0, \beta_1, \ldots, \beta_{J_T})'$, the network output vector is compactly written as $H(X;\theta) = O(X\Gamma)\vec{\beta}$.

\begin{definition}[Activation Functions]
    An \textit{activation function} $S : \mathbb{R} \to \mathbb{R}$ introduces the nonlinearity that makes neural networks universal approximators. A \emph{squashing} (sigmoid) function satisfies $\lim_{u \to +\infty} S(u) = b$ and $\lim_{u \to -\infty} S(u) = a$ for finite $a < b$. Common choices include:
    \begin{itemize}[noitemsep]
        \item \textbf{Logistic:} $S(x) = \frac{1}{1 + e^{-x}}$, mapping $\mathbb{R} \to (0,1)$.
        \item \textbf{Hyperbolic tangent:} $S(x) = \tanh(x) = \frac{e^x - e^{-x}}{e^x + e^{-x}}$, mapping $\mathbb{R} \to (-1,1)$.
        \item \textbf{Rectified Linear Unit (ReLU):} $S(x) = \max(0,x)$, a non-squashing function with no saturation for $x > 0$.
        \item \textbf{Radial Basis Function (RBF):} $S(x) = \exp(-x^2)$, a non-squashing function with localized response.
    \end{itemize}
    The logistic and hyperbolic tangent functions are bounded and smooth; ReLU is piecewise linear and is the dominant choice in modern deep networks because it avoids vanishing gradients and produces sparse activations.
\end{definition}

\begin{remark}[Geometric Interpretation]
    Each hidden unit $j$ computes a scalar $\vec{\gamma}_j\cdot\vec{X}_t + \gamma_{0,j}$, whose level set $\{\vec{x} : \vec{\gamma}_j\cdot\vec{x} = \gamma_{0,j}\}$ is a hyperplane in $\mathbb{R}^p$ with normal $\vec{\gamma}_j$ and signed distance $\gamma_{0,j}/\|\vec{\gamma}_j\|$ from the origin. The activation function $S$ maps this signed distance to a nonlinear response. With $J$ hidden units, the covariate space is partitioned into at most $M(J,p)$ polyhedral regions, where $M(J,p) = M(J-1,p)+M(J-1,p-1)$ with $M(1,p)=2$ and $M(J,1)=J+1$. Each region receives a distinct affine prediction, so the network is equivalent to a piecewise-affine function.
\end{remark}


\begin{definition}[Empirical Risk Minimization for Neural Networks]
    Given a loss function $L$ and sample $\{(Y_t, \vec{X}_t)\}_{t=1}^T$, the parameters are estimated by minimizing the empirical risk:
    \begin{align}
        R_T(\theta) = \frac{1}{T}\sum_{t=1}^T L\!\left(Y_t,\, H(\vec{X}_t;\theta)\right).
    \end{align}
    For the quadratic loss, $R_T(\theta) = \frac{1}{T}\|\vec{Y} - H(X;\theta)\|_2^2$. Because $H$ is nonlinear in $\theta$, no closed-form minimizer exists and the problem is non-convex, admitting multiple local minima.
\end{definition}

\begin{definition}[Gradient Descent]
    Starting from an initial guess $\theta^{(0)}$, \textit{gradient descent} iterates
    \begin{align}
        \theta^{(k+1)} = \theta^{(k)} - \eta_k \nabla_\theta R_T\!\left(\theta^{(k)}\right),
    \end{align}
    where $\eta_k > 0$ is the \emph{learning rate}. Choosing $\eta_k$ too large causes oscillation; too small causes slow convergence. \emph{Stochastic gradient descent} (SGD) replaces the full gradient with a gradient computed on a random ``mini-batch'' of observations at each iteration, which reduces computational cost and introduces noise that can help escape shallow local minima.
\end{definition}

The gradient $\nabla_\theta R_T$ is computed analytically via the \emph{backpropagation} algorithm, which applies the chain rule along the computational graph $\vec{\gamma}_j \to z_{j,t} \to a_{j,t} \to H(\vec{X}_t;\theta) \to R_T(\theta)$, where $z_{j,t} = \vec{\gamma}_j\cdot\vec{X}_t + \gamma_{0,j}$ and $a_{j,t} = S(z_{j,t})$.

\begin{proposition}[Backpropagation Gradients]
    For the quadratic loss with a single hidden layer, define the output error $\delta_t = -\frac{2}{T}(Y_t - H(\vec{X}_t;\theta))$. Then
    \begin{align}
        \frac{\partial R_T}{\partial \beta_j} &= \sum_{t=1}^T \delta_t \, S(z_{j,t}), \\
        \frac{\partial R_T}{\partial \vec{\gamma}_j} &= \sum_{t=1}^T \delta_t \beta_j S'(z_{j,t})\,\vec{X}_t.
    \end{align}
    The error $\delta_t$ at the output layer is propagated backward through the hidden layer via the chain rule, weighted by $\beta_j S'(z_{j,t})$. This computation scales linearly in $T$ and $J_T$, making training feasible for large networks.
\end{proposition}


\begin{definition}[Regularization and Complexity Control]
    The number of hidden units $J_T$ controls the bias-variance tradeoff of the network:
    \begin{itemize}[noitemsep]
        \item If $J_T$ is too small, the network underfits: high bias, low variance.
        \item If $J_T$ is too large, the network may overfit: low training error, high out-of-sample error.
    \end{itemize}
    Overfitting is controlled via \emph{weight decay} (an $\ell^2$ penalty $\lambda\|\theta\|_2^2$ added to $R_T$) and \emph{early stopping} (halting training when validation error begins to rise). Practical guidance is to start with a moderately large $J_T$ and rely on regularization rather than deliberate underfitting; the network architecture determines \emph{potential} complexity while regularization determines \emph{effective} complexity.
\end{definition}

\begin{remark}
    The tuning parameter $J_T$ may be selected by cross-validation over a grid (e.g., $J_T \in \{10, 20, 50, 100, 200\}$), choosing the value minimizing out-of-sample MSE. Information criteria are harder to apply here because the effective degrees of freedom of a nonlinear sieve are not easily characterized analytically.
\end{remark}


A deep neural network (DNN) extends the single hidden layer architecture by stacking $L > 1$ hidden layers. While the universal approximation theorem holds for shallow networks, depth allows compositional functions to be approximated with exponentially fewer parameters than a shallow network of equivalent expressiveness (Mhaskar, Liao, and Poggio, 2017).

\begin{definition}[Deep Neural Network]
    Let $J_\ell$ denote the number of hidden units in layer $\ell \in \{1,\ldots,L\}$, and let $\Gamma_\ell \in \mathbb{R}^{(J_{\ell-1}+1) \times J_\ell}$ collect the weights for layer $\ell$. Define the output matrix of layer $\ell$ recursively by
    \begin{align}
        O_\ell\!\left(O_{\ell-1}(\cdot)\Gamma_\ell\right)_{[T \times (J_\ell+1)]} :=
        \begin{pmatrix}
            1 & S(\widetilde{\vec{\gamma}}_{1\ell}\cdot O_{1,\ell-1}) & \cdots & S(\widetilde{\vec{\gamma}}_{J_\ell \ell}\cdot O_{1,\ell-1}) \\
            \vdots & \vdots & \ddots & \vdots \\
            1 & S(\widetilde{\vec{\gamma}}_{1\ell}\cdot O_{T,\ell-1}) & \cdots & S(\widetilde{\vec{\gamma}}_{J_\ell \ell}'\cdot O_{T,\ell-1})
        \end{pmatrix},
    \end{align}
    with $O_0 = X$ (the raw input). The full network output is the composition
    \begin{align}
        H(X;\theta) = O_L\!\left(O_{L-1}\!\left(\cdots O_1(X\Gamma_1)\cdots\right)\Gamma_L\right)\vec{\beta}.
    \end{align}
    The total parameter count is $\sum_{\ell=1}^L (J_{\ell-1}+1)J_\ell + J_L + 1$, which can easily exceed the sample size $T$, requiring strong regularization.
\end{definition}

\begin{figure}[h]
    \centering
    \includegraphics[width=0.5\linewidth]{neural.png}
    \caption{Neural network}
    \label{fig:placeholder}
\end{figure}

\begin{definition}[Dropout]
    \textit{Dropout} is a regularization technique for deep networks. At each training step, each neuron in each hidden layer is independently set to zero with probability $1-q$ (typically $q = 0.5$). Formally, the layer-$\ell$ output under dropout is
    \begin{align}
        y_i^{(\ell+1)} &= \vec{\gamma}_i^{(\ell+1)}\!\left(\vec{r}^{(\ell)} \odot \vec{z}^{(\ell)}\right) + b_i^{(\ell+1)}, \quad \vec{r}^{(\ell)} \sim \mathrm{Bernoulli}(q)^{J_\ell},
    \end{align}
    where $\odot$ denotes elementwise multiplication. At test time, the full network is used with weights scaled by $q$, recovering the same expected output as at training time. Dropout is equivalent to training an ensemble of $2^{\sum_\ell J_\ell}$ ``thinned'' sub-networks, which significantly reduces overfitting.
\end{definition}

\begin{remark}[Challenges of Deep Network Training]
    Deep networks present two main challenges beyond those of shallow networks. First, the large parameter count demands correspondingly large training datasets and significant computational resources. Second, standard gradient descent suffers from \emph{vanishing gradients}: when squashing activations (logistic, tanh) are saturated, $S'(z_{j,t}) \approx 0$, causing the backpropagated error signal to diminish exponentially with depth. ReLU activations largely resolve this by preserving gradient magnitude for positive pre-activations. Stochastic gradient descent with dropout and appropriate learning-rate schedules is the standard training recipe.
\end{remark}

\begin{table}[H]
\centering
\renewcommand{\arraystretch}{1.3}
\begin{tabular}{@{}lll@{}}
\toprule
\textbf{Architecture} & \textbf{Key advantage} & \textbf{Key challenge} \\
\midrule
Shallow (1 hidden layer) & Universal approximation; interpretable & May require very large $J_T$ \\
Deep ($L > 1$ hidden layers) & Compositional efficiency; fewer units & Vanishing gradients; overfit risk \\
ReLU activation & No gradient saturation; fast training & Non-smooth; unbounded outputs \\
Dropout regularization & Ensemble effect; reduces overfit & Longer training; tuning $q$ \\
\bottomrule
\end{tabular}
\caption{Comparison of neural network design choices.}
\end{table}

\begin{remark}[Neural Networks in Finance: Broader Perspective]
    The two empirical examples above illustrate a recurring theme: neural networks deliver the largest gains over parametric benchmarks when the true function has \emph{unknown interaction structure} (return prediction, where characteristic interactions are complex and data-driven) or \emph{known smooth structure with parametric misspecification} (option pricing, where the residual surface is smooth and well-approximated by a low-dimensional network). In both settings, the key design choice is whether to use the network as a direct universal approximator or as a residual corrector layered on top of a theoretically motivated parametric model. The latter approach is often preferable when interpretability or no-arbitrage constraints matter, as in derivatives pricing and risk management.
\end{remark}

\section{Large Language Models}

Large language models (LLMs) are neural networks trained to model the probability distribution over sequences of text. They are the foundation of modern natural language processing and have demonstrated remarkable generalization across tasks ranging from translation and summarization to code generation and mathematical reasoning. LLMs are relevant both as a research topic and as a measurement tool in empirical finance: many economically relevant state variables (cash-flow news, managerial sentiment, regulatory uncertainty, discount-rate expectations) are observed only indirectly through language.


\begin{definition}[Language Model]
    A language model defines a probability distribution over sequences of tokens $(t_1, t_2, \dots, t_T)$ drawn from a vocabulary $\mathcal{V}$. By the chain rule of probability:
    \begin{align}
        \mathbb{P}(t_1, t_2, \dots, t_T) = \prod_{\tau=1}^T \mathbb{P}(t_\tau \mid t_1, \dots, t_{\tau-1}).
    \end{align}
    The fundamental task is \emph{next-token prediction}: the model learns a sequence of conditional distributions, from which generation proceeds one token at a time by sampling.
\end{definition}

\begin{remark}[LLMs as Conditional Probability Models]
    An LLM is a very large conditional probability model. The object being estimated is a sequence of conditional distributions $\{P(\cdot \mid x_1,\ldots,x_{t-1})\}_{t=1}^T$, making the framework naturally statistical. The model is generative because a joint distribution is built out of those conditionals via the chain rule. What is unusual relative to familiar econometric settings is the dimension and flexibility of the function class, not the logic of estimation.
\end{remark}


\begin{definition}[Tokenization]
    Raw text must first be converted into a finite alphabet of \emph{tokens}. Tokens need not be whole words; typical tokens include words, subwords, punctuation, and special symbols. For example, ``unbelievable'' $\to$ (``un'', ``believ'', ``able''), which allows the alphabet to be smaller, since many words are constructed from these roots. The token vocabulary is finite, so tokenization maps arbitrary text into a manageable discrete alphabet, so that the input text can be given as a set of indices $(t_1, t_2,...)$ pointing to alphabet entries.
\end{definition}

\begin{figure}[h]
    \centering
    \includegraphics[width=0.6\linewidth]{embedding_matrix.png}
    \caption{The embedding matrix}
    \label{fig:placeholder}
\end{figure}

\begin{definition}[Token Embedding]
    Each token index $i \in \mathcal{V}$ is mapped to a learned continuous vector $\vec{e}(i) \in \mathbb{R}^d$ via an \emph{embedding matrix} $E \in \mathbb{R}^{|\mathcal{V}| \times d}$, where the ''one-hot'' index vector picks a row. Similar contexts tend to push tokens toward similar numerical representations; this is the first stage of representation learning. Since the model must also track the position of each token in the sequence, a \emph{positional encoding} $\vec{p}_\tau \in \mathbb{R}^d$ is added:
    \begin{align}
        \vec{h}_\tau^{(0)} = E_{t_\tau} + \vec{p}_\tau.
    \end{align}
    The standard sinusoidal positional encoding uses $p_{\tau, 2k} = \sin(\tau / 10000^{2k/d})$ and $p_{\tau, 2k+1} = \cos(\tau / 10000^{2k/d})$, ensuring that the encoding varies smoothly with position. In modern models, learned position schemes are also common.
\end{definition}

\begin{remark}[Statistical Interpretation of Embeddings]
    Embeddings are a learned basis expansion: the model replaces a categorical variable by a dense, low-dimensional vector representation. This is analogous to latent-factor constructions such as PCA, but unlike PCA, the representation is learned jointly with the prediction objective rather than estimated from a covariance matrix. In finance language, embeddings can be viewed as learned latent characteristics.
\end{remark}

\begin{remark}[Necessity of Positional Information]
    Transformers process tokens largely in parallel, which is computationally attractive but removes any built-in notion of sequence order. Without position information, the model cannot distinguish ``rates affect inflation'' from ``inflation affect rates.''
\end{remark}


\begin{definition}[Scaled Dot-Product Attention]
    Let $X \in \mathbb{R}^{T \times d}$ collect the current token representations of the input sequence in the embedding space. Self-attention forms three linear projections:
    \begin{align}
        Q = XW^Q, \quad K = XW^K, \quad V = XW^V,
    \end{align}
    where $W^Q, W^K \in \mathbb{R}^{d \times d_k}$ and $W^V \in \mathbb{R}^{d \times d_v}$ are learned projection matrices. The roles are:
    \begin{itemize}[noitemsep]
        \item \textbf{Query}: what a token is looking for.
        \item \textbf{Key}: how a token advertises itself as relevant.
        \item \textbf{Value}: the information the token contributes.
    \end{itemize}
    The attention output is:
    \begin{align}
        \text{Attention}(Q, K, V) = \text{softmax}\!\left(\frac{QK^T}{\sqrt{d_k}}\right)V.
    \end{align}
    The $(i,j)$ entry of $QK^T / \sqrt{d_k}$ is the attention score of position $i$ attending to position $j$, measuring their compatibility. Softmax normalizes these scores into a probability distribution over positions with $\text{softmax}(\vec{z})_i = \frac{e^{z_i}}{\sum_je^{z_j}}$. Dividing by $\sqrt{d_k}$ stabilizes gradients by preventing the dot products from growing large in magnitude.
\end{definition}

\begin{remark}[One-Token Formula]
    For a single token $i$, the updated representation is thus
    \begin{align}
        \tilde{v}_i = \sum_{j=1}^T w_{ij} v_j, \qquad w_{ij} = \frac{\exp(q_i\cdot k_j / \sqrt{d_k})}{\sum_{\ell=1}^T \exp(q_i\cdot k_\ell / \sqrt{d_k})}.
    \end{align}
    The weights $w_{ij}$ are endogenous and observation-specific: different tokens in the same sentence can use very different weighting schemes. Each token's representation becomes a context-dependent weighted average of all other tokens' value vectors.
\end{remark}

\begin{remark}[Attention as a Learned Kernel Smoother]
    Defining $W(X) = \text{softmax}(QK^T/\sqrt{d_k})$, attention can be written as $\text{Attention}(Q,K,V) = W(X)V$, where $W(X)$ is a data-dependent weighting matrix with row $i$ containing the weights used by token $i$ to aggregate information from all other tokens. This is a learned, context-dependent smoothing operator that resembles a nonlinear kernel smoother, but the similarity metric is itself learned from data. Hence, attention is adaptive weighting under a learned geometry.
\end{remark}

\begin{remark}[Causal Masking]
    In autoregressive language modeling, token $t$ must not look ahead to future tokens. This is enforced by \emph{causal masking}: setting the attention scores for $j > t$ to $-\infty$ before the softmax, so those positions receive zero attention weight.
\end{remark}

\begin{remark}[Self-Attention vs.\ Cross-Attention]
    In self-attention, queries, keys, and values all come from the same sequence $X$: the sequence updates itself using information from itself. In \emph{cross-attention}, queries come from one sequence while keys and values come from another. Cross-attention is important in encoder--decoder architectures (e.g., machine translation), where the decoder attends to the encoder's output.
\end{remark}

\begin{figure}[h]
    \centering
    \includegraphics[width=0.6\linewidth]{attn_is_all.png}
    \caption{The attention mechanism}
    \label{fig:placeholder}
\end{figure}

\begin{definition}[Multi-Head Attention]
    Rather than computing a single attention function, \emph{multi-head attention} runs $H$ attention functions in parallel with different learned projections, then concatenates and projects the results:
    \begin{align}
        \text{head}_h &= \text{Attention}(XW_h^Q,\, XW_h^K,\, XW_h^V), \quad h = 1, \ldots, H, \\
        \text{MultiHead}(X) &= \text{Concat}(\text{head}_1, \dots, \text{head}_H)\,W^O,
    \end{align}
    where $W_h^Q, W_h^K \in \mathbb{R}^{d \times d_k}$, $W_h^V \in \mathbb{R}^{d \times d_v}$, and $W^O \in \mathbb{R}^{Hd_v \times d}$ are learned projection matrices. Different heads can specialize: one head may focus on local syntax, another on long-range dependencies, another on entity identification or causal relations. A single head is like one learned signal extractor; multi-head attention lets the model examine the same information set through several learned lenses at once.
\end{definition}

\begin{definition}[Transformer Block]
    A \emph{transformer} is a stack of $L$ identical blocks. Each block applies multi-head attention followed by a position-wise feedforward network, with residual connections and layer normalization at each step:
    \begin{align}
        \vec{H}' &= \text{LayerNorm}\!\left(\vec{H}^{(\ell-1)} + \text{MultiHead}(\vec{H}^{(\ell-1)})\right), \\
        \vec{H}^{(\ell)} &= \text{LayerNorm}\!\left(\vec{H}' + \text{FFN}(\vec{H}')\right),
    \end{align}
    where $\text{FFN}(\vec{h}) = W_2 \, \sigma(W_1 \vec{h} + \vec{b}_1) + \vec{b}_2$ is applied identically to each position, with $W_1 \in \mathbb{R}^{d \times d_{ff}}$ and $W_2 \in \mathbb{R}^{d_{ff} \times d}$.
\end{definition}

\begin{remark}[Residual Connections]
    The \textit{residual connection} $\text{Output} = x + F(x)$ means the layer learns a correction to the input. If no change is needed, $F(x) \approx 0$. Residual connections preserve information from earlier layers, improve gradient flow in deep networks, and stabilize training. The analogy is: new estimate $=$ old estimate $+$ adjustment.
\end{remark}

\begin{remark}[Attention vs.\ Feedforward: Communication vs.\ Computation]
    Attention mixes information \emph{across} tokens (inter-position communication), while the feedforward block transforms information \emph{within} each token (intra-position computation). The feedforward network maps $\mathbb{R}^{d_{\text{model}}} \to \mathbb{R}^{d_{ff}} \to \mathbb{R}^{d_{\text{model}}}$, introducing nonlinearity through the activation function $\sigma$. Without the FFN, the model would be largely linear. Globally: attention $=$ communication, FFN $=$ computation.
\end{remark}

\begin{remark}[Layer Normalization]
    Layer normalization normalizes the activations across the feature dimension for each token independently: $\text{LayerNorm}(\vec{h}) = \gamma \odot \frac{\vec{h} - \mu}{\sigma} + \beta$, where $\mu$ and $\sigma$ are the mean and standard deviation of $\vec{h}$, and $\gamma, \beta$ are learned scale and shift parameters. This stabilizes training by preventing internal covariate shift, analogous to how batch normalization stabilizes activations across the batch dimension.
\end{remark}


\begin{definition}[Output Layer]
    After $L$ transformer blocks, the final hidden state $H \in \mathbb{R}^{T \times d}$ is projected to the vocabulary via a linear map $W_{\text{vocab}} \in \mathbb{R}^{d \times |\mathcal{V}|}$:
    \begin{align}
        Z = HW_{\text{vocab}} \in \mathbb{R}^{T \times |\mathcal{V}|}, \qquad P(x_t \mid x_{<t}) = \text{softmax}(Z_t) \in \mathbb{R}^{|\mathcal{V}|}.
    \end{align}
    Each row of the output gives a probability distribution over the entire vocabulary for the next token at that position.
\end{definition}


\begin{remark}[Architectural Variants]
    Transformer architectures come in three forms:
    \begin{itemize}[noitemsep]
        \item \textbf{Encoder}: processes the full input sequence bidirectionally (e.g., BERT). Useful for classification and embedding tasks.
        \item \textbf{Decoder}: generates text autoregressively with causal masking (e.g., GPT). This is the architecture underlying modern LLMs.
        \item \textbf{Encoder--Decoder}: maps input sequences to output sequences via cross-attention (e.g., the original Transformer for translation).
    \end{itemize}
    Modern LLMs are typically decoder-only models, trained autoregressively.
\end{remark}

\begin{figure}[h]
    \centering
    \includegraphics[width=0.9\linewidth]{transformers.png}
    \caption{Various transformer architectures}
    \label{fig:placeholder}
\end{figure}


\begin{remark}[Transformers vs.\ Recurrent Neural Networks]
    In a recurrent neural network (RNN), information passes sequentially through hidden states, so learning long-range dependencies requires many steps and suffers from vanishing gradients. In a transformer, tokens interact directly via attention: any token can access any other in one step. This yields easier learning of long-range dependencies, more stable optimization, and highly parallelizable computation. The transformer has essentially replaced the RNN as the dominant sequence model.
\end{remark}

\begin{remark}[LLMs as Nonlinear Sieves]
    The transformer architecture is a direct generalization of the nonlinear sieve framework from Section~7. A single hidden layer sigmoid neural network (sANN) was introduced as a nonlinear sieve with parameterized basis functions $h_j(\vec{X}) = S(\vec{\gamma}_j\cdot\vec{X} + \gamma_{0,j})$. The transformer replaces the fixed nonlinearity $S$ with the attention mechanism, a data-dependent, adaptive weighting operator, while retaining the same fundamental structure: a composition of learned nonlinear maps whose expressiveness grows with width (number of heads, hidden dimension) and depth (number of layers). The Universal Approximation Theorem guarantees that sufficiently large transformers can approximate any continuous function, just as for shallow networks, but depth allows compositional functions to be represented with exponentially fewer parameters.
\end{remark}

\begin{definition}[Cross-Entropy Loss]
    The model is trained to maximize the log-likelihood of observed text, equivalently minimizing the \emph{cross-entropy loss}:
    \begin{align}
        \hat{\theta} = \arg\min_\theta \left[-\sum_{\tau=1}^T \log P_\theta(t_\tau \mid t_1, \dots, t_{\tau-1})\right].
    \end{align}
    This \emph{next-token prediction} objective is self-supervised: the labels are derived from the data itself, requiring no human annotation. Training an LLM is maximum likelihood estimation.
\end{definition}

\begin{proposition}[Cross-Entropy and KL Divergence]
    If $p^*(\cdot \mid x_{<t})$ denotes the true conditional distribution and $p_\theta(\cdot \mid x_{<t})$ the model, then the expected cross-entropy decomposes as
    \begin{align}
        \mathcal{H}(p^*, p_\theta) = \mathcal{H}(p^*) + D_{\mathrm{KL}}(p^* \,\|\, p_\theta).
    \end{align}
    The entropy term $\mathcal{H}(p^*)$ does not depend on $\theta$. Minimizing expected cross-entropy is therefore equivalent to minimizing KL divergence from the true distribution.
\end{proposition}
\begin{proof}
    By definition,
    \begin{align}
        \mathcal{H}(p^*, p_\theta) &= -\sum_x p^*(x) \log p_\theta(x)
        = -\sum_x p^*(x) \log p^*(x) + \sum_x p^*(x) \log \frac{p^*(x)}{p_\theta(x)}\\
        &= \mathcal{H}(p^*) + D_{\mathrm{KL}}(p^* \,\|\, p_\theta).
    \end{align}
\end{proof}

\begin{remark}[Empirical Scaling Laws]
    The scaling behavior of LLMs is empirically well-described by \emph{power laws}. For a model with $N$ parameters trained on $D$ tokens using compute $C$, the test cross-entropy loss scales approximately as
    \begin{align}
        \mathcal{L}(N,D,C) \approx aN^{-\alpha} + bD^{-\beta} + mC^{-\gamma} + c,
    \end{align}
    for constants $a, b, m, c, \alpha, \beta, \gamma > 0$. Performance improves smoothly and predictably with scale, but with diminishing returns. These scaling laws allow researchers to optimally allocate compute between model size, data, and training to minimize prediction error.
\end{remark}

\begin{remark}[Implications of Scaling]
    Several empirical findings are notable. First, larger models can be more sample-efficient, learning more per token of training data. Second, compute-efficient training often stops before full convergence. Third, the naive strategy of ``train a smaller model much longer'' is frequently suboptimal relative to training a larger model for fewer steps. Fourth, aspects of final model performance can be predicted from smaller pilot runs, giving firms a way to forecast whether a major training run is likely to justify its cost. The exact formula is model-dependent, but the broader lesson is robust: parameters, data, and compute all matter jointly, and none should be thought of in isolation.
\end{remark}

\begin{remark}[The Alignment Problem]
    Next-token prediction is a powerful objective, but it is not the same thing as optimizing for helpfulness, truthfulness, or safety. Internet text contains noise, conflicts, and undesirable behavior; likelihood maximization does not directly encode preference for truthful or helpful outputs. A model can be statistically good at predicting internet text while still being poor at following a user's intent. Post-training alignment addresses this gap through a three-stage pipeline: (1) supervised fine-tuning on curated responses, (2) reward model estimation from human preference rankings, and (3) reinforcement learning from human feedback.
\end{remark}

\begin{definition}[Pre-training and Fine-tuning]
    Modern LLMs are trained in two stages. In \emph{pre-training}, the model is trained on a large corpus of unlabeled text using the next-token prediction objective, learning general linguistic structure, world knowledge, and reasoning patterns. In \emph{fine-tuning} (also called supervised fine-tuning, SFT), the pre-trained model is further trained on a smaller dataset of curated prompt--response pairs, adapting its behavior to follow instructions. This \emph{transfer learning} paradigm is highly parameter-efficient: fine-tuning updates far fewer parameters than training from scratch, and often achieves strong performance with limited task-specific data.
\end{definition}

\begin{definition}[Reward Modeling]
    For the same prompt, humans rank several candidate outputs. A separate \emph{reward model} $R_\phi(x, y)$ is trained to predict those preferences, where $x$ is a prompt and $y$ is a response. The reward model is a learned proxy for the human preference relation: it maps (prompt, response) pairs to scalar scores such that $R_\phi(x, y_1) > R_\phi(x, y_2)$ whenever $y_1$ is preferred to $y_2$.
\end{definition}

\begin{definition}[Reinforcement Learning from Human Feedback]
    Starting from the supervised model, \emph{reinforcement learning from human feedback} (RLHF) optimizes the language model policy $\pi_\theta$ to maximize expected reward subject to a KL divergence penalty:
    \begin{align}
        \max_\theta \mathbb{E}_{\vec{y} \sim \pi_\theta(\cdot|\vec{x})}\!\left[R_\phi(\vec{x}, \vec{y})\right] - \beta \, D_{\text{KL}}\!\left(\pi_\theta(\cdot|\vec{x}) \,\|\, \pi_{\text{ref}}(\cdot|\vec{x})\right),
    \end{align}
    where $\pi_{\text{ref}}$ is the supervised fine-tuned model. The reward term encourages preferred outputs; the KL penalty prevents the policy from drifting too far from the reference, which could lead to reward hacking (exploiting weaknesses in the reward model to produce pathological but high-scoring outputs).
\end{definition}

\begin{remark}[Economic Interpretation of RLHF]
    The components of RLHF have natural economic analogues:
    \begin{itemize}[noitemsep]
        \item Reward model $\approx$ learned utility function.
        \item Human rankings $\approx$ revealed preferences over outputs.
        \item KL regularization $\approx$ adjustment cost or penalty for departing from reference behavior, acting as a trust-region or model-risk control.
        \item RLHF $\approx$ constrained policy optimization under noisy preference data.
    \end{itemize}
    RLHF changes the behavior of the model but does not replace the knowledge acquired during pretraining. It adds a preference-alignment layer on top of the statistical model learned from data.
\end{remark}

\begin{remark}[Limitations]
    Even highly capable LLMs still display important limitations: hallucinations (plausible-sounding but factually incorrect outputs), brittleness outside the training distribution, sensitivity to prompt phrasing, and incomplete reliability in high-stakes settings. These limitations require careful human oversight in any production financial application.
\end{remark}

\begin{definition}[Econometric Mapping from Text to Returns]
    Let $\mathcal{T}_{i,t}$ denote the text associated with firm $i$ at time $t$. An LLM produces an embedding
    \begin{align}
        \vec{\phi}_{i,t} = \Phi(\mathcal{T}_{i,t}) \in \mathbb{R}^k,
    \end{align}
    where $\Phi$ is the LLM's representation function. A predictive specification can then be written as
    \begin{align}
        R_{i,t+1} = \alpha + \vec{\beta}'\vec{\phi}_{i,t} + \vec{\gamma}'\vec{Z}_{i,t} + \varepsilon_{i,t+1},
    \end{align}
    where $\vec{Z}_{i,t}$ collects standard firm characteristics. The LLM is not the asset-pricing model itself; it is a \emph{measurement device} that maps raw text into latent predictors. The downstream predictive regression remains an econometric object. More generally, a nonlinear factor model formulation is
    \begin{align}
        R_{i,t+1} = f\!\left(\Phi(\mathcal{T}_{i,t}),\, \vec{Z}_{i,t}\right) + \varepsilon_{i,t+1},
    \end{align}
    where the embedding vector plays the role of latent text factors and $f(\cdot)$ determines how those factors load into returns.
\end{definition}

\begin{remark}[What Latent Text Factors May Capture]
    The embedding $\vec{\phi}_{i,t}$ potentially summarizes information about cash-flow news, discount-rate news, managerial sentiment, uncertainty, disagreement, exposure to macro narratives, and risk related to regulation or technology. Unlike PCA-based factors, these representations are context-sensitive: the same word (``bank,'' ``volatility,'' ``exposure'') receives different embeddings depending on the surrounding text.
\end{remark}

\begin{remark}[Empirical Pipeline]
    A standard empirical pipeline for LLM-based asset pricing proceeds as follows: (1) collect text at date $t$; (2) preserve timing to avoid look-ahead bias; (3) obtain embeddings or prompt-based labels; (4) combine with standard firm characteristics; (5) estimate predictive models for $R_{i,t+1}$; (6) evaluate out of sample and compute economic value. Good empirical design remains essential: one must respect the information set available at time $t$, benchmark against strong simpler models, use genuine out-of-sample validation, assess whether gains survive transaction costs, and distinguish statistical significance from economic significance.
\end{remark}

\end{document}